% Sensibilisation sur les rôles de la méiose et de la fécondation dans le maintien de la diversité génétique des individus au sein d'une espèce — SVT Tle D — Cours signé OnBuch+
% Programme officiel MINESEC. Compiler avec : tectonic N04-brassage-genetique-diversite.tex
\newcommand{\DOCMATIERE}{SVT}
\newcommand{\DOCNIVEAU}{Tle D}
\newcommand{\DOCMODULE}{Module I : Le Monde Vivant — Brassage génétique et unicité génétique des individus}
\newcommand{\DOCLECON}{N04}
\newcommand{\DOCTITRE}{Sensibilisation sur les rôles de la méiose et de la fécondation dans le maintien de la diversité génétique des individus au sein d'une espèce}
% OnBuch+ — préambule « cours signés » ENRICHI (compilé avec tectonic / XeLaTeX)
% Système visuel « Pop » d'OnBuch+ : fond crème (jamais blanc), bordures encre
% épaisses, ombres « dures » décalées, titres Archivo Black en capitales.
% Version MATHÉMATIQUES : graphes (pgfplots/tikz), boîtes pédagogiques
% (définition, propriété, méthode, attention, le savais-tu, exercice résolu,
% exercice, TP, bilan), page de garde et signature OnBuch+ ancrée en bas.
%
% Macros attendues dans main.tex AVANT \input{preamble} :
%   \DOCMATIERE  \DOCNIVEAU  \DOCMODULE  \DOCLECON (numéro)  \DOCTITRE
\documentclass[11pt]{article}

\usepackage{fontspec}
\usepackage[french]{babel}
\usepackage[a4paper,margin=2cm,top=3.4cm,bottom=2.8cm,headheight=1.4cm,footskip=1.3cm]{geometry}
\usepackage{amsmath,amssymb,amsfonts}
\usepackage{mathtools} % \xmapsto, \coloneqq... souvent utilisés par les modèles
\usepackage{cancel} % simplifications barrées (bilans de forces)
% \abs{z} (module, valeur absolue) et \norm{u} (norme), utilisés par les modèles
\providecommand{\abs}{}\let\abs\relax\DeclarePairedDelimiter{\abs}{\lvert}{\rvert}
\providecommand{\norm}{}\let\norm\relax\DeclarePairedDelimiter{\norm}{\lVert}{\rVert}
\usepackage[table]{xcolor} % [table] : fournit aussi \rowcolors (alternance de lignes)
\usepackage{pagecolor}
\usepackage{tikz}
\usetikzlibrary{arrows.meta,positioning,calc,shapes.geometric,shapes.misc,shapes.arrows,decorations.pathmorphing,decorations.pathreplacing,decorations.markings,patterns,shadows,mindmap,trees,fit,backgrounds,matrix,intersections,angles,quotes}
\usepackage{pgfplots}
\usepackage[version=4]{mhchem} % équations nucléaires \ce{^{235}_{92}U}
\usepackage[european,siunitx]{circuitikz} % schémas électriques
\pgfplotsset{compat=1.18}
\usepgfplotslibrary{fillbetween,groupplots} % aires entre courbes (calcul intégral)
\usepackage{enumitem}
\usepackage{fancyhdr}
% [most] charge toutes les bibliothèques tcolorbox (listings, minted, raster,
% documentation...) et gonfle inutilement la mémoire TeX sur un document long
% (~300 boîtes) : on ne charge que ce qui est réellement utilisé.
\usepackage[skins,breakable]{tcolorbox}
\usepackage{graphicx}
\usepackage{colortbl}
\usepackage{booktabs}
\usepackage{tabularx}
\usepackage{diagbox} % case d'en-tête diagonale des tableaux à double entrée
% Colonnes extensibles centrée / gauche / droite (C, L, R), souvent utilisées par les modèles
\newcolumntype{C}{>{\centering\arraybackslash}X}
\newcolumntype{L}{>{\raggedright\arraybackslash}X}
\newcolumntype{R}{>{\raggedleft\arraybackslash}X}
\usepackage{array}
\usepackage{multicol}
\usepackage{siunitx}
% parse-numbers=false : accepte aussi \SI{2,5 \times 10^{-3}}{...}
\sisetup{locale=FR,output-decimal-marker={,},group-minimum-digits=4,parse-numbers=false}
\DeclareSIUnit\ohm{\ensuremath{\symup{\Omega}}} % U+2126 absent de la police
\DeclareSIUnit\division{div} % réglages d'oscilloscope (ms/div, V/div)
\DeclareSIUnit\tr{tr} % tours (vitesse de rotation, ex. \SI{1500}{\tr\per\minute})
\DeclareSIUnit\molar{M} % concentration molaire (ex. \SI{3}{\molar}, \SI{1}{\micro\molar})
\DeclareSIUnit\an{an} % année (le modèle confond parfois avec \year, un compteur TeX) — ex. \SI{5}{\kilo\meter\per\an}
\DeclareSIUnit\FCFA{FCFA} % franc CFA (ex. \SI{500}{\FCFA\per\kilo\gram})
\DeclareSIUnit\degreeBrix{\textdegree Brix} % teneur en sucre d'un sirop/jus (réfractométrie)
\DeclareSIUnit\month{mois} % le modèle utilise parfois \month, un compteur TeX, comme unité de durée
\usepackage{qrcode}
% \degree : commande fréquemment utilisée par les modèles pour un angle ou une
% température en dehors de \SI{}{} (non fournie par les paquets chargés).
\providecommand{\degree}{^{\circ}}
% \qed (fin de démonstration), utilisé par les modèles sans amsthm
\providecommand{\qed}{\hfill$\square$}
% \barre : double barre des valeurs interdites dans un tableau de variations (tkz-tab)
\providecommand{\barre}{\ensuremath{\|}}
% environnement proof (amsthm non chargé) utilisé par les modèles
\ifdefined\proof\else\newenvironment{proof}[1][Démonstration]{\par\noindent\textit{#1.}\ }{\qed\par}\fi
\usepackage{lastpage}
\usepackage{hyperref}
\hypersetup{hidelinks}

% ============================================================
% Palette « Pop » OnBuch+ — mêmes hex que la classe P (plus_ui.dart)
% ============================================================
\definecolor{poporange}{HTML}{F59321}
\definecolor{popdark}{HTML}{E07A0C}
\definecolor{popcream}{HTML}{FFF6E8}
\definecolor{popink}{HTML}{1C1714}
\definecolor{poppurple}{HTML}{7A5AE0}
\definecolor{popgreen}{HTML}{2E9E6A}
\definecolor{poppink}{HTML}{E0457B}
\definecolor{popblue}{HTML}{2F7DE1}
\definecolor{popgold}{HTML}{C98A12}
\definecolor{popmuted}{HTML}{8A7F76}
\definecolor{popline}{HTML}{EADFD2}
\definecolor{popgray}{HTML}{8A7F76}
\colorlet{popgrey}{popgray}
% Alias tolérants (le modèle nomme parfois les couleurs différemment)
\colorlet{popwhite}{white}
\colorlet{popblack}{popink}
\colorlet{popred}{poppink}
\colorlet{popyellow}{popgold}
% Teintes claires (fonds de tableaux/encadrés) — le modèle nomme parfois
% les couleurs avec un suffixe L (ex. popblueL) sans les avoir définies.
\colorlet{poporangeL}{poporange!20}
\colorlet{popdarkL}{popdark!20}
\colorlet{poppurpleL}{poppurple!20}
\colorlet{popgreenL}{popgreen!20}
\colorlet{poppinkL}{poppink!20}
\colorlet{popblueL}{popblue!20}
\colorlet{popgoldL}{popgold!20}
\colorlet{popredL}{popred!20}
\colorlet{popgrayL}{popgray!20}
% Teintes claires pour fonds de tableaux / zones
\colorlet{popblueL}{popblue!10!white}
\colorlet{poporangeL}{poporange!14!white}
\colorlet{popgreenL}{popgreen!12!white}
\colorlet{poppurpleL}{poppurple!12!white}
\colorlet{poppinkL}{poppink!12!white}
\colorlet{popgoldL}{popgold!14!white}

\pagecolor{popcream}

% ============================================================
% Typographie
% ============================================================
\defaultfontfeatures{Ligatures=TeX}
\setmainfont{PlusJakartaSans-Medium.ttf}[Path=../../../fonts/,BoldFont=PlusJakartaSans-Bold.ttf]
% Maths en sans-serif (Fira Math) avec chiffres et lettres droites de Plus
% Jakarta, pour que les formules se fondent dans le texte.
\usepackage[math-style=ISO,bold-style=ISO]{unicode-math}
\setmathfont{FiraMath-Regular.otf}[Path=../../../fonts/]
\setmathfont{PlusJakartaSans-Medium.ttf}[Path=../../../fonts/,range={up/num,up/{Latin,latin},\mathup}]
\usepackage{icomma} % virgule décimale sans espace en mode math
% Caractères mathématiques Unicode absents des polices de texte (Plus Jakarta,
% Archivo Black) : rendus par leur équivalent mathématique, en texte comme en
% formule — sinon ils s'affichent en carré vide (ex. « Arithmétique dans ℤ »).
\usepackage{newunicodechar}
\newunicodechar{ℝ}{\ensuremath{\mathbb{R}}}
\newunicodechar{ℤ}{\ensuremath{\mathbb{Z}}}
\newunicodechar{ℂ}{\ensuremath{\mathbb{C}}}
\newunicodechar{ℕ}{\ensuremath{\mathbb{N}}}
\newunicodechar{ℚ}{\ensuremath{\mathbb{Q}}}
\newunicodechar{𝔻}{\ensuremath{\mathbb{D}}}
\newunicodechar{α}{\ensuremath{\alpha}}
\newunicodechar{β}{\ensuremath{\beta}}
\newunicodechar{γ}{\ensuremath{\gamma}}
\newunicodechar{δ}{\ensuremath{\delta}}
\newunicodechar{ε}{\ensuremath{\varepsilon}}
\newunicodechar{θ}{\ensuremath{\theta}}
\newunicodechar{λ}{\ensuremath{\lambda}}
\newunicodechar{μ}{\ensuremath{\mu}}
\newunicodechar{σ}{\ensuremath{\sigma}}
\newunicodechar{τ}{\ensuremath{\tau}}
\newunicodechar{φ}{\ensuremath{\varphi}}
\newunicodechar{ψ}{\ensuremath{\psi}}
\newunicodechar{ω}{\ensuremath{\omega}}
\newunicodechar{Σ}{\ensuremath{\Sigma}}
% majuscules produites par \MakeUppercase dans les titres (α-aminés -> Α)
\newunicodechar{Α}{\ensuremath{\alpha}}
\newunicodechar{Β}{\ensuremath{\beta}}
\newunicodechar{Γ}{\ensuremath{\Gamma}}
\newunicodechar{Δ}{\ensuremath{\Delta}}
\newunicodechar{Ε}{\ensuremath{\varepsilon}}
\newunicodechar{Θ}{\ensuremath{\Theta}}
\newunicodechar{Λ}{\ensuremath{\Lambda}}
\newunicodechar{Μ}{\ensuremath{\mu}}
\newunicodechar{Τ}{\ensuremath{\tau}}
\newunicodechar{Φ}{\ensuremath{\Phi}}
\newunicodechar{Ψ}{\ensuremath{\Psi}}
\newunicodechar{Ω}{\ensuremath{\Omega}}
\newunicodechar{↦}{\ensuremath{\mapsto}}
\newunicodechar{∈}{\ensuremath{\in}}
\newunicodechar{∉}{\ensuremath{\notin}}
\newunicodechar{∧}{\ensuremath{\wedge}}
\newunicodechar{⇔}{\ensuremath{\Leftrightarrow}}
\newunicodechar{⟺}{\ensuremath{\Longleftrightarrow}}
\newunicodechar{⇒}{\ensuremath{\Rightarrow}}
\newunicodechar{⟹}{\ensuremath{\Longrightarrow}}
\newunicodechar{∘}{\ensuremath{\circ}}
\newunicodechar{∪}{\ensuremath{\cup}}
\newunicodechar{∩}{\ensuremath{\cap}}
\newunicodechar{≡}{\ensuremath{\equiv}}
\newunicodechar{⊂}{\ensuremath{\subset}}
\newunicodechar{⊥}{\ensuremath{\perp}}
\newunicodechar{∃}{\ensuremath{\exists}}
\newunicodechar{∀}{\ensuremath{\forall}}
\newunicodechar{∛}{\ensuremath{\sqrt[3]{\,}}}
\newunicodechar{∜}{\ensuremath{\sqrt[4]{\,}}}
\newunicodechar{⃗}{\ensuremath{{}^{\to}}}
\newunicodechar{ⁿ}{\ifmmode^{n}\else\textsuperscript{\ensuremath{n}}\fi}
\newunicodechar{ˣ}{\ifmmode^{x}\else\textsuperscript{\ensuremath{x}}\fi}
\newunicodechar{ᵇ}{\ifmmode^{b}\else\textsuperscript{\ensuremath{b}}\fi}
\newunicodechar{ʳ}{\ifmmode^{r}\else\textsuperscript{\ensuremath{r}}\fi}
\newunicodechar{ᵘ}{\ifmmode^{u}\else\textsuperscript{\ensuremath{u}}\fi}
\newunicodechar{ʸ}{\ifmmode^{y}\else\textsuperscript{\ensuremath{y}}\fi}
\newunicodechar{ᵃ}{\ifmmode^{a}\else\textsuperscript{\ensuremath{a}}\fi}
\newunicodechar{ᵉ}{\ifmmode^{e}\else\textsuperscript{\ensuremath{e}}\fi}
\newunicodechar{ᵏ}{\ifmmode^{k}\else\textsuperscript{\ensuremath{k}}\fi}
\newunicodechar{ᵖ}{\ifmmode^{p}\else\textsuperscript{\ensuremath{p}}\fi}
\newunicodechar{ᴺ}{\ifmmode^{N}\else\textsuperscript{\ensuremath{N}}\fi}
\newunicodechar{ᶜ}{\ifmmode^{c}\else\textsuperscript{\ensuremath{c}}\fi}
\newunicodechar{ʰ}{\ifmmode^{h}\else\textsuperscript{\ensuremath{h}}\fi}
\newunicodechar{ᵠ}{\ifmmode^{q}\else\textsuperscript{\ensuremath{q}}\fi}
\newunicodechar{ₙ}{\ifmmode_{n}\else\textsubscript{\ensuremath{n}}\fi}
\newunicodechar{ₐ}{\ifmmode_{a}\else\textsubscript{\ensuremath{a}}\fi}
\newunicodechar{ᵢ}{\ifmmode_{i}\else\textsubscript{\ensuremath{i}}\fi}
\newunicodechar{ᵣ}{\ifmmode_{r}\else\textsubscript{\ensuremath{r}}\fi}
\newunicodechar{ₖ}{\ifmmode_{k}\else\textsubscript{\ensuremath{k}}\fi}
\newunicodechar{ᵦ}{\ifmmode_{\beta}\else\textsubscript{\ensuremath{\beta}}\fi}
\newunicodechar{ₓ}{\ifmmode_{x}\else\textsubscript{\ensuremath{x}}\fi}
\newfontfamily\popdisplay{ArchivoBlack-Regular.ttf}[Path=../../../fonts/]
\newfontfamily\poplogo{SpaceGrotesk-Bold.ttf}[Path=../../../fonts/]
\newfontfamily\popsemi{PlusJakartaSans-SemiBold.ttf}[Path=../../../fonts/]
\newfontfamily\popxbold{PlusJakartaSans-ExtraBold.ttf}[Path=../../../fonts/]
\newfontfamily\popxboldls{PlusJakartaSans-ExtraBold.ttf}[Path=../../../fonts/,LetterSpace=3.0]

\setlist[enumerate]{leftmargin=1.7em,itemsep=5pt,topsep=4pt}
\setlist[itemize]{leftmargin=1.7em,itemsep=4pt,topsep=4pt,label={\color{poporange}\textbullet}}
\setlength{\parindent}{0pt}
\setlength{\parskip}{5pt}
\renewcommand{\arraystretch}{1.25}

% pgfplots : style Pop par défaut
\pgfplotsset{
  every axis/.append style={axis line style={popink,line width=1pt},tick style={popink},
    grid style={popline,line width=0.5pt},label style={font=\small},tick label style={font=\footnotesize}},
  popcurve/.style={poporange,line width=1.8pt,no markers,smooth},
}
% Même style disponible dans un \draw TikZ simple (hors pgfplots)
\tikzset{popcurve/.style={poporange,line width=1.8pt}}
\tikzset{popgrid/.style={popline,very thin}}
\colorlet{popcurve}{poporange} % le nom du style est parfois utilisé comme couleur

% ============================================================
% Pill / squiggle
% ============================================================
\newcommand{\poppill}[3]{%
  \tikz[baseline=-0.55ex]\node[draw=popink,line width=1.2pt,rounded corners=2.6pt,%
    fill=#2,inner xsep=6.5pt,inner ysep=3pt]{%
    \popxboldls\fontsize{7}{7}\selectfont\color{#3}\MakeUppercase{#1}};%
}
\newcommand{\popsquiggle}[1]{%
  \tikz[baseline=-0.4ex]{
    \draw[#1,line width=2.6pt,line cap=round]
      plot [smooth] coordinates {(0,0.09) (0.34,0.01) (0.68,0.21) (1.02,0.01) (1.36,0.10)};
  }%
}
% Mot-clé surligné (marqueur orange sous le texte)
\newcommand{\cle}[1]{{\popxbold\color{popdark}#1}}

% ============================================================
% En-tête / pied
% ============================================================
\pagestyle{fancy}
\fancyhf{}
\renewcommand{\headrulewidth}{0pt}
\renewcommand{\footrulewidth}{0pt}
\lhead{\raisebox{-0.15ex}{\poplogo\fontsize{13}{13}\selectfont\color{popink}OnBuch\color{poporange}+}}
\rhead{\poppill{\DOCMATIERE\ \textbullet\ \DOCNIVEAU\ \textbullet\ Leçon \DOCLECON}{white}{popink}}
\lfoot{\popxbold\fontsize{7.6}{7.6}\selectfont\color{popmuted}\MakeUppercase{Cours signé OnBuch+}\ \textbullet\ {\popsemi\DOCTITRE}}
\rfoot{\tikz[baseline=-0.6ex]\node[rounded corners=8.5pt,fill=popink,minimum height=17pt,minimum width=17pt,inner xsep=5pt,inner sep=0]{\popxbold\fontsize{8}{8}\selectfont\color{popcream}\thepage\,/\,\pageref*{LastPage}};}

% ============================================================
% Titres
% ============================================================
\newcommand{\chaptitre}[2]{%
  \begin{tcolorbox}[enhanced,colback=poporange,colframe=popink,boxrule=2.6pt,arc=11pt,
      drop shadow=popink,left=15pt,right=15pt,top=12pt,bottom=11pt,before skip=3pt,after skip=18pt]
    \poppill{#2}{popink}{popcream}\par\vspace{9pt}
    {\raggedright\hyphenpenalty=10000\exhyphenpenalty=10000\popdisplay\fontsize{22}{24}\selectfont\color{popcream}\MakeUppercase{#1}\par}
  \end{tcolorbox}
}
% Section numérotée : pastille orange avec le numéro + titre Archivo
\newcounter{popsec}
\newcommand{\coursec}[1]{%
  \stepcounter{popsec}\setcounter{popsub}{0}%
  \par\vspace{16pt}\noindent
  \begin{minipage}{\linewidth}
  \tikz[baseline=-0.7ex]\node[draw=popink,line width=1.6pt,fill=poporange,rounded corners=4pt,minimum size=22pt,inner sep=0]{\popdisplay\fontsize{12}{12}\selectfont\color{popcream}\thepopsec};%
  \hspace{8pt}{\popdisplay\fontsize{15}{17}\selectfont\color{popink}\MakeUppercase{#1}}\par
  \vspace{4pt}\noindent{\color{poporange}\rule{36pt}{3.6pt}}
  \end{minipage}\par\nopagebreak\vspace{8pt}
}
\newcounter{popsub}
\newcommand{\courssub}[1]{%
  \stepcounter{popsub}%
  \par\vspace{9pt}\noindent{\popxbold\fontsize{11.5}{13}\selectfont\color{popink}{\color{poporange}\thepopsec.\thepopsub}\hspace{6pt}#1}\par\nopagebreak\vspace{4pt}
}

% ============================================================
% Boîtes pédagogiques — même recette : fond blanc, bordure épaisse colorée,
% ombre dure encre, badge plein. Titre optionnel [#1].
% ============================================================
\tcbset{popbase/.style={breakable,enhanced,colback=white,boxrule=2.3pt,arc=9pt,
  shadow={3pt}{-3pt}{0pt}{popink},left=14pt,right=14pt,top=11pt,bottom=11pt,before skip=11pt,after skip=15pt,
  fontupper=\popsemi\fontsize{10.6}{14.5}\selectfont\color{popink},
  fontlower=\popsemi\fontsize{10.6}{14.5}\selectfont\color{popink}}}
% Coche dessinée (le glyphe ✓ n'existe pas dans les polices du document)
\newcommand{\popcheck}[1]{\tikz[baseline=-0.5ex]\draw[#1,line width=1.6pt,line cap=round,line join=round](-0.13,0.0)--(-0.04,-0.09)--(0.13,0.1);}
\providecommand{\coche}{\popcheck{popgreen}}
\providecommand{\resultat}[1]{\par\smallskip\noindent\fcolorbox{popgreen}{popgreenL}{\parbox{\dimexpr\linewidth-2\fboxsep-2\fboxrule\relax}{\textbf{Résultat.} #1}}\par\smallskip}
\newcommand{\popboxhead}[3]{\poppill{#1}{#2}{white}\ifx\relax#3\relax\else\hspace{6pt}{\popxbold\fontsize{10.6}{12}\selectfont\color{popink}#3}\fi\par\vspace{7pt}}

\newtcolorbox{aretenir}[1][]{popbase,colframe=popgreen,before upper={\popboxhead{À retenir}{popgreen}{#1}}}
\newtcolorbox{exemplebox}[1][]{popbase,colframe=poporange,before upper={\popboxhead{Exemple}{poporange}{#1}}}
\newtcolorbox{definition}[1][]{popbase,colframe=popblue,before upper={\popboxhead{Définition}{popblue}{#1}}}
\newtcolorbox{propriete}[1][]{popbase,colframe=poppurple,before upper={\popboxhead{Propriété}{poppurple}{#1}}}
\newtcolorbox{methode}[1][]{popbase,colframe=popgold,before upper={\popboxhead{Méthode}{popgold}{#1}}}
\newtcolorbox{attention}[1][]{popbase,colframe=poppink,before upper={\popboxhead{Attention}{poppink}{#1}}}
\newtcolorbox{experience}[1][]{popbase,colframe=popblue,colback=popblueL,before upper={\popboxhead{Expérience}{popblue}{#1}}}
% « Le savais-tu ? » : carte violette pleine (contexte, culture, Cameroun)
\newtcolorbox{savaistu}[1][]{popbase,colback=poppurple,colframe=popink,
  fontupper=\popsemi\fontsize{10.4}{14.2}\selectfont\color{white},
  before upper={\poppill{Le savais-tu\,?}{white}{poppurple}\ifx\relax#1\relax\else\hspace{6pt}{\popxbold\color{white}#1}\fi\par\vspace{7pt}}}
% Exercice résolu : énoncé en haut, \tcblower puis la solution sur fond crème
\newcounter{popexo}\newcounter{popexr}
\newtcolorbox{exoresolu}[1][]{popbase,colframe=poporange,colbacklower=popcream,
  segmentation style={popink,line width=1.2pt,dashed},
  before upper={\stepcounter{popexr}\popboxhead{Exercice résolu \thepopexr}{poporange}{#1}},
  before lower={\poppill{Solution}{popink}{popcream}\par\vspace{6pt}}}
% Exercice (non résolu en place) — la correction est dans la partie Corrigés
\newtcolorbox{exercice}[1][]{popbase,colframe=popink,
  before upper={\stepcounter{popexo}\popboxhead{Exercice \thepopexo}{popink}{#1}}}
% Corrigé
\newtcolorbox{corrige}[1][]{popbase,colframe=popgreen,colback=popgreenL,
  before upper={\popboxhead{Corrigé}{popgreen}{#1}}}
% Objectifs / prérequis / bilan
\newtcolorbox{objectifs}{popbase,colframe=popink,colback=poporangeL,
  before upper={\popboxhead{Objectifs de la leçon}{popink}{}}}
\newtcolorbox{prerequis}{popbase,colframe=popink,colback=popblueL,
  before upper={\popboxhead{Prérequis}{popblue}{}}}
\newtcolorbox{bilan}{popbase,colframe=popink,colback=white,boxrule=2.8pt,
  before upper={\popboxhead{Fiche bilan}{poporange}{L'essentiel à connaître pour l'examen}}}
% Figure encadrée avec légende
\newtcolorbox{popfigure}[1][]{enhanced,breakable=false,colback=white,colframe=popline,boxrule=1.4pt,arc=8pt,
  left=10pt,right=10pt,top=10pt,bottom=8pt,before skip=10pt,after skip=12pt,halign=center,
  lower separated=false,halign lower=center,
  fontlower=\popsemi\fontsize{9}{11.5}\selectfont\color{popmuted},
  #1}
\newcommand{\legende}[1]{\par\vspace{6pt}{\popsemi\fontsize{9}{11.5}\selectfont\color{popmuted}#1}}

% ============================================================
% Signature OnBuch+
% ============================================================
\newcommand{\poplogoinline}[1]{{\poplogo\fontsize{#1}{#1}\selectfont\color{popink}OnBuch\color{poporange}+}}
\newcommand{\signatureonbuch}{%
  \begin{tcolorbox}[enhanced,colback=popink,colframe=popink,boxrule=2.2pt,arc=10pt,
      drop shadow=poporange,left=14pt,right=14pt,top=10pt,bottom=10pt,before skip=0pt,after skip=0pt]
    \tikz[baseline=-0.7ex]\node[circle,fill=popgreen,draw=popcream,line width=1.3pt,minimum size=22pt,inner sep=0]{\popcheck{white}};%
    \hspace{9pt}\begin{minipage}[c]{0.62\linewidth}
      {\popxbold\fontsize{12}{13}\selectfont\color{popcream}Signé\ \textbullet\ {\poplogo OnBuch\color{poporange}+}}\par\vspace{2pt}
      {\raggedright\popsemi\fontsize{8}{10.5}\selectfont\color{popline}\MakeUppercase{Équipe pédagogique OnBuch+}\\\MakeUppercase{Conforme au programme officiel MINESEC}\par}
    \end{minipage}\hfill
    {\poplogo\fontsize{22}{22}\selectfont\color{popcream}OnBuch\color{poporange}+}
  \end{tcolorbox}%
}

% ============================================================
% Page de garde
% #1 titre · #2 matière · #3 niveau · #4 module · #5 n° leçon · #6 durée ·
% #7 description / objectifs (texte ou itemize)
% ============================================================
\newcommand{\pagegarde}[7]{%
  \thispagestyle{empty}
  \newgeometry{a4paper,margin=1.7cm,top=1.6cm,bottom=1.4cm}%
  \begin{tikzpicture}[remember picture,overlay]
    \fill[poporange!18] ([xshift=-3.2cm,yshift=-3.6cm]current page.north east) circle (4.6cm);
    \fill[poppurple!14] ([xshift=2.4cm,yshift=7.6cm]current page.south west) circle (3.8cm);
  \end{tikzpicture}%
  {\poplogo\fontsize{30}{30}\selectfont\color{popink}OnBuch\color{poporange}+}\hfill
  \raisebox{0.6ex}{\poppill{Cours signé \textbullet\ Premium}{popink}{popcream}}\par
  \vspace{1pt}\popsquiggle{poppurple}\par
  \vspace{8pt}
  \begin{tcolorbox}[enhanced,colback=poporange,colframe=popink,boxrule=2.8pt,arc=13pt,
      drop shadow=popink,left=18pt,right=18pt,top=12pt,bottom=13pt,after skip=10pt]
    \poppill{#2 \textbullet\ #3}{popink}{popcream}\hspace{5pt}\poppill{#4}{white}{popink}\par\vspace{8pt}
    {\popdisplay\fontsize{14}{14}\selectfont\color{popink}LEÇON #5}\par\vspace{4pt}
    {\raggedright\hyphenpenalty=10000\exhyphenpenalty=10000\popdisplay\fontsize{25}{27}\selectfont\color{popcream}\MakeUppercase{#1}\par}
    \vspace{8pt}
    {\popxbold\fontsize{9.5}{9.5}\selectfont\color{popink}\MakeUppercase{Durée conseillée : #6}}
  \end{tcolorbox}
  \begin{tcolorbox}[enhanced,colback=white,colframe=popink,boxrule=2.2pt,arc=10pt,
      drop shadow=popink,left=16pt,right=16pt,top=10pt,bottom=10pt,after skip=8pt]
    \poppill{Dans cette leçon}{poppurple}{white}\par\vspace{5pt}
    \popsemi\fontsize{9.3}{12.6}\selectfont\color{popink}#7
  \end{tcolorbox}
  \noindent\begin{minipage}[t]{0.487\linewidth}
    \begin{tcolorbox}[enhanced,colback=popgreen,colframe=popink,boxrule=2.2pt,arc=10pt,
        drop shadow=popink,left=11pt,right=11pt,top=10pt,bottom=10pt,height=3.1cm,valign=center]
      \tcbox[colback=white,colframe=popink,boxrule=1.3pt,arc=5pt,boxsep=0pt,nobeforeafter,
        left=3pt,right=3pt,top=3pt,bottom=3pt,valign=center]{\qrcode[height=1.9cm]{https://play.google.com/store/apps/details?id=com.luvvix.onbuch&hl=fr}}%
      \hspace{8pt}\begin{minipage}[c]{0.5\linewidth}
        {\popdisplay\fontsize{11}{12.5}\selectfont\color{white}\MakeUppercase{Télécharge OnBuch}}\par\vspace{4pt}
        {\raggedright\popsemi\fontsize{8.4}{11}\selectfont\color{white}Gratuit \textbullet\ Google Play\\Tuteur IA, annales, concours, résultats\par}
      \end{minipage}
    \end{tcolorbox}
  \end{minipage}\hfill
  \begin{minipage}[t]{0.487\linewidth}
    \begin{tcolorbox}[enhanced,colback=poppurple,colframe=popink,boxrule=2.2pt,arc=10pt,
        drop shadow=popink,left=11pt,right=11pt,top=10pt,bottom=10pt,height=3.1cm,valign=center]
      \tikz[baseline=-0.7ex]\node[circle,fill=white,draw=popink,line width=1.3pt,minimum size=1.6cm,inner sep=0]{\popdisplay\fontsize{24}{24}\selectfont\color{poppurple}+};%
      \hspace{8pt}\begin{minipage}[c]{0.6\linewidth}
        {\popdisplay\fontsize{11}{12.5}\selectfont\color{white}\MakeUppercase{Passe à OnBuch+}}\par\vspace{4pt}
        {\raggedright\popsemi\fontsize{8.4}{11}\selectfont\color{white}Tous les cours signés, Tuteur IA illimité, fiches PDF — onglet OnBuch+ de l'app\par}
      \end{minipage}
    \end{tcolorbox}
  \end{minipage}
  \vspace{8pt}
  \popcontenu
  \vfill
  \signatureonbuch
  \restoregeometry
  \newpage
}

% Rangée « Ce cours contient » (page de garde)
\newcommand{\poptile}[3]{% #1 couleur #2 gros texte #3 légende
  \begin{tcolorbox}[enhanced,colback=#1,colframe=popink,boxrule=2pt,arc=9pt,shadow={2.5pt}{-2.5pt}{0pt}{popink},
      left=6pt,right=6pt,top=9pt,bottom=9pt,halign=center,valign=center,height=1.75cm,width=\linewidth,nobeforeafter]
    {\popdisplay\fontsize{12.5}{13}\selectfont\color{white}\MakeUppercase{#2}}\par\vspace{3pt}
    {\centering\popsemi\fontsize{7.8}{9.5}\selectfont\color{white}#3\par}
  \end{tcolorbox}}
\newcommand{\popcontenu}{%
  \par\noindent{\popxboldls\fontsize{7.5}{7.5}\selectfont\color{popmuted}CE COURS SIGNÉ CONTIENT}\par\vspace{7pt}
  \noindent\begin{minipage}[t]{0.235\linewidth}\poptile{popblue}{Cours}{complet, illustré, pas à pas}\end{minipage}\hfill
  \begin{minipage}[t]{0.235\linewidth}\poptile{poporange}{Méthodes}{et exercices résolus}\end{minipage}\hfill
  \begin{minipage}[t]{0.235\linewidth}\poptile{poppink}{Exercices}{type examen + corrigés}\end{minipage}\hfill
  \begin{minipage}[t]{0.235\linewidth}\poptile{popgreen}{Fiche bilan}{l'essentiel à retenir}\end{minipage}\par}

% Sommaire
\newtcolorbox{sommaire}{enhanced,colback=white,colframe=popink,boxrule=2.2pt,arc=10pt,
  drop shadow=popink,left=18pt,right=18pt,top=13pt,bottom=13pt,before skip=6pt,after skip=20pt,
  before upper={\poppill{Sommaire}{poppurple}{white}\par\vspace{9pt}\popsemi\fontsize{10.6}{14.5}\selectfont\color{popink}}}

% Signature de fin de cours — poussée tout en bas de la dernière page
\newcommand{\findecours}{%
  \par\vfill\nopagebreak
  \begin{center}\popsquiggle{poporange}\end{center}\vspace{4pt}
  \signatureonbuch
}
\AtBeginDocument{\def\llbracket{\mathopen{[\mkern-3.5mu[}}\def\rrbracket{\mathclose{]\mkern-3.5mu]}}} % intervalles d'entiers (glyphe ⟦ absent de la police)
% Glyphes absents de Fira Math : redéfinis à partir de glyphes disponibles
\AtBeginDocument{%
  \def\setminus{\mathbin{\backslash}}\let\smallsetminus\setminus
  \def\vdots{\mathord{\vbox{\baselineskip3.2pt\lineskiplimit0pt\kern4pt\hbox{.}\hbox{.}\hbox{.}}}}%
  \def\ddots{\mathinner{\mkern1mu\raise6.5pt\hbox{.}\mkern2mu\raise3.6pt\hbox{.}\mkern2mu\raise0.7pt\hbox{.}\mkern1mu}}%
  \def\triangle{\mathord{\scalebox{0.9}{$\symup{\Delta}$}}}%
  \def\nabla{\mathord{\raisebox{\depth}{\scalebox{1}[-1]{$\symup{\Delta}$}}}}%
  \def\checkmark{\popcheck{popdark}}%
  \def\star{\mathbin{\ast}}\def\bigstar{\mathord{\ast}}%
  \def\diamond{\mathbin{\scalebox{0.6}{$\Diamond$}}}%
  \def\bigcup{\mathop{\mathchoice{\raisebox{-0.3em}{\scalebox{1.7}{$\cup$}}}{\scalebox{1.25}{$\cup$}}{\cup}{\cup}}\displaylimits}%
  \def\bigcap{\mathop{\mathchoice{\raisebox{-0.3em}{\scalebox{1.7}{$\cap$}}}{\scalebox{1.25}{$\cap$}}{\cap}{\cap}}\displaylimits}%
  \def\lesseqqgtr{\mathrel{\lessgtr}}\def\bigoplus{\mathop{\oplus}\displaylimits}%
}
\providecommand{\ssi}{si et seulement si} % abréviation fréquente
\AtBeginDocument{\providecommand{\wideparen}{\overparen}} % arc de cercle

%% FIGFIX-BEGIN v3 — correctifs globaux des figures (docs/onbuch-generation/tools/figfix)
\usepackage{adjustbox}\usepackage{etoolbox}
% toute figure TikZ/pgfplots plus large que la ligne est réduite à la largeur disponible
\BeforeBeginEnvironment{tikzpicture}{\begin{adjustbox}{max width=\linewidth}}
\AfterEndEnvironment{tikzpicture}{\end{adjustbox}}
% légendes pgfplots lisibles par-dessus les courbes
\pgfplotsset{every axis/.append style={legend style={fill=white,fill opacity=0.92,text opacity=1,draw=black!15,inner sep=3pt}}}
% étiquettes d'arêtes (syntaxe edge["..."]) avec fond blanc
\tikzset{every edge quotes/.append style={fill=white,inner sep=1.2pt}}
% bulles de cartes mentales : texte à la ligne dans la bulle, bulle un peu plus grande
\tikzset{every concept/.append style={text width=2.0cm,align=center,minimum size=2.3cm,inner sep=2pt}}
%% FIGFIX-END
\begin{document}

\pagegarde{Sensibilisation sur les rôles de la méiose et de la fécondation dans le maintien de la diversité génétique des individus au sein d'une espèce}{SVT}{Tle D}{Module I : Le Monde Vivant — Brassage génétique et unicité génétique des individus}{N04}{8 h}{Cette leçon étudie les mécanismes du brassage génétique (interchromosomique et intrachromosomique) à partir des expériences historiques de Mendel et de Morgan, et montre comment ils expliquent la diversité des individus au sein d'une espèce. Elle fournit les outils d'interprétation des croisements (mono-, di- et trihybridisme), de calcul des taux de recombinaison et de construction des cartes génétiques, indispensables pour l'épreuve de SVT du Baccalauréat.
\vspace{4pt}
\begin{multicols}{2}\small
\begin{itemize}
\item À la fin de cette leçon, je sais interpréter les résultats d'un croisement de monohybridisme et en déduire les relations de dominance entre allèles
\item interpréter les résultats d'expériences de dihybridisme et de trihybridisme chez les souris, volailles, plantes à fleurs et drosophiles
\item distinguer la ségrégation indépendante des gènes (brassage interchromosomique) du linkage total ou partiel (brassage intrachromosomique)
\item expliquer le rôle de la méiose (brassages) et de la fécondation dans la diversité génétique des individus
\item déterminer le pourcentage de recombinaison entre deux gènes à partir des résultats d'un croisement-test
\item établir la carte génétique d'un chromosome à partir d'un exercice de dihybridisme ou de trihybridisme avec gènes liés
\end{itemize}\end{multicols}}

\begin{sommaire}
\begin{multicols}{2}
\begin{enumerate}
\item Rappels et monohybridisme : les lois de Mendel
\item Dihybridisme et trihybridisme : le brassage interchromosomique
\item Linkage et crossing-over : le brassage intrachromosomique
\item Pourcentage de recombinaison et carte génétique
\item Exceptions à la monogénie : polygénie et pléiotropie
\item Synthèse : méiose, fécondation et diversité génétique
\item Activité d'intégration
\item Méthodes et exercices résolus
\item Exercices
\item Corrigés des exercices
\item Fiche bilan
\end{enumerate}
\end{multicols}
\end{sommaire}

\begin{objectifs}
\begin{itemize}
\item À la fin de cette leçon, je sais interpréter les résultats d'un croisement de monohybridisme et en déduire les relations de dominance entre allèles
\item interpréter les résultats d'expériences de dihybridisme et de trihybridisme chez les souris, volailles, plantes à fleurs et drosophiles
\item distinguer la ségrégation indépendante des gènes (brassage interchromosomique) du linkage total ou partiel (brassage intrachromosomique)
\item expliquer le rôle de la méiose (brassages) et de la fécondation dans la diversité génétique des individus
\item déterminer le pourcentage de recombinaison entre deux gènes à partir des résultats d'un croisement-test
\item établir la carte génétique d'un chromosome à partir d'un exercice de dihybridisme ou de trihybridisme avec gènes liés
\item relever les exceptions à la monogénie (polygénie, pléiotropie) et interpréter les résultats expérimentaux correspondants
\item rédiger un échiquier de croisement et un raisonnement génétique rigoureux conforme aux exigences du Bac
\end{itemize}
\end{objectifs}

\begin{prerequis}
\begin{itemize}
\item Méiose : étapes, chiasmas et rôle dans la production de gamètes haploïdes (notions de Première)
\item Fécondation : rencontre aléatoire des gamètes et restauration de la diploïdie
\item Notions de gène, allèle, locus, génotype, phénotype, homozygote/hétérozygote
\item Chromosomes homologues et caryotype
\item Proportions et calculs de pourcentages (outils mathématiques simples)
\end{itemize}
\end{prerequis}

\newpage

\coursec{Rappels et monohybridisme : les lois de Mendel}

Dans un village de la région du Centre, une famille d'éleveurs de volailles remarque que les poulets issus d'un même couple de reproducteurs présentent des plumages très variés : blancs, noirs, barrés, et même des combinaisons jamais observées chez les parents. Au marché de Ngaoundéré, un producteur de maïs constate que les grains de ses épis hybrides ne ressemblent ni au père ni à la mère. Comment expliquer que des parents identiques produisent des descendants si différents ? Pourquoi certains caractères semblent-ils toujours transmis ensemble alors que d'autres se combinent librement ?

Pour répondre à ces questions, il faut remonter à la source de la génétique moderne : les croisements de pois réalisés par le moine Gregor Mendel au XIX\textsuperscript{e} siècle. Cette première section pose le vocabulaire fondamental, présente l'expérience historique du monohybridisme et dégage les deux premières lois de Mendel, que la méiose — étudiée en classe de Première — permet aujourd'hui d'expliquer au niveau cellulaire.

\courssub{Le vocabulaire indispensable de la génétique}

Avant d'interpréter le moindre croisement, il faut parler le langage des généticiens. Chaque terme ci-dessous sera utilisé constamment dans toute la leçon : prends le temps de les maîtriser.

\begin{definition}[Gène, allèle, locus]
Un \cle{gène} est un segment d'ADN porté par un chromosome, qui conditionne l'expression d'un caractère héréditaire (forme de la graine, couleur des fleurs, groupe sanguin...). Le \cle{locus} est l'emplacement précis occupé par ce gène sur le chromosome. Un gène peut exister sous plusieurs versions appelées \cle{allèles} : par exemple, le gène « forme de la graine de pois » possède l'allèle \emph{lisse} et l'allèle \emph{ridé}.
\end{definition}

\begin{definition}[Génotype et phénotype]
Le \cle{génotype} est l'ensemble des allèles portés par un individu pour le ou les caractères étudiés ; on le note entre parenthèses, par exemple $(L//L)$ ou $(L//r)$. Le \cle{phénotype} est l'expression observable du génotype, notée entre crochets : $[L]$ pour « graines lisses », $[r]$ pour « graines ridées ». Le phénotype résulte du génotype et, pour certains caractères, de l'influence du milieu.
\end{definition}

Un individu est dit \cle{homozygote} pour un gène lorsqu'il porte deux allèles identiques $(L//L)$ ou $(r//r)$ ; il est dit \cle{hétérozygote} lorsqu'il porte deux allèles différents $(L//r)$. Un individu homozygote ne produit qu'\textbf{un seul type} de gamètes, alors qu'un hétérozygote en produit \textbf{deux types} en proportions égales : c'est une conséquence directe de la méiose.

\begin{definition}[Dominance, récessivité, codominance]
Chez un hétérozygote, l'\cle{allèle dominant} est celui qui s'exprime seul dans le phénotype ; l'\cle{allèle récessif} est celui qui ne s'exprime qu'à l'état homozygote. On parle de \cle{codominance} lorsque les deux allèles s'expriment simultanément chez l'hétérozygote, donnant un phénotype nouveau, composite (exemple classique : les groupes sanguins $A$ et $B$ chez l'Homme, où l'individu de génotype $(A//B)$ est de groupe $[AB]$).
\end{definition}

\begin{attention}[Génotype $\neq$ phénotype]
L'erreur la plus fréquente au Baccalauréat consiste à confondre génotype et phénotype. Deux individus de phénotype identique $[L]$ peuvent avoir des génotypes différents : $(L//L)$ ou $(L//r)$. Autre piège : croire que l'allèle récessif a « disparu » chez l'hybride F1. Il est toujours présent dans le génotype, simplement \emph{masqué} dans le phénotype, et il réapparaîtra en F2. Enfin, n'écris jamais $(L/L)$ avec une barre simple : la convention est $(L//L)$, les deux barres symbolisant les deux chromosomes homologues.
\end{attention}

\courssub{Le monohybridisme : définition}

\begin{definition}[Monohybridisme]
Le \cle{monohybridisme} est l'étude de la transmission d'\textbf{un seul caractère} (une seule paire d'allèles) au cours de croisements entre individus de races pures différentes. Les parents, dits de « race pure », sont homozygotes pour le caractère considéré.
\end{definition}

\courssub{L'expérience historique de Mendel sur le pois}

\begin{experience}[Les croisements de Mendel sur le pois (1856--1863)]
\textbf{Matériel et choix de l'organisme.} Gregor Mendel choisit le pois (\emph{Pisum sativum}), plante à fleurs dont les organes reproducteurs sont protégés (autofécondation naturelle), ce qui garantit la pureté des lignées. Il peut aussi réaliser des pollinisations croisées contrôlées en castrant les fleurs (ablation des étamines) puis en déposant le pollen choisi sur le pistil.

\textbf{Protocole.} Mendel croise deux races pures de pois différant par un caractère : des plants à \textbf{graines lisses} $\times$ des plants à \textbf{graines ridées} (et, dans une autre série, fleurs violettes $\times$ fleurs blanches). Il laisse ensuite les hybrides F1 s'autoféconder pour obtenir la génération F2.

\textbf{Résultats chiffrés.}
\begin{itemize}
\item En F1 : \textbf{toutes} les graines sont lisses (uniformité de la F1).
\item En F2, sur \num{7324} graines : \num{5474} lisses et \num{1850} ridées.
\end{itemize}

\textbf{Interprétation.} Le caractère « ridé », absent en F1, réapparaît en F2 : il n'avait donc pas disparu, il était masqué. Le rapport $\dfrac{5474}{1850} \approx 2,96 \approx \dfrac{3}{1}$ suggère que chaque hybride F1 porte les deux déterminants du caractère et les transmet séparément à ses gamètes.
\end{experience}

\courssub{Uniformité de la F1 : la dominance}

L'observation est frappante : quel que soit le sens du croisement (graine lisse $\times$ ridée ou l'inverse), 100 \% des hybrides F1 présentent le phénotype $[L]$ « lisse ». Mendel en déduit que l'allèle \emph{lisse} \textbf{domine} l'allèle \emph{ridé}, dit récessif.

Écrivons le croisement de façon rigoureuse. Les parents de race pure sont $(L//L)$, phénotype $[L]$, et $(r//r)$, phénotype $[r]$. Chacun ne produit qu'un type de gamète :

\begin{center}
\begin{tabularx}{\linewidth}{|l|X|X|}
\hline
\rowcolor{popblueL} & Parent 1 & Parent 2 \\
\hline
Phénotype & $[L]$ lisse & $[r]$ ridée \\
\hline
Génotype & $(L//L)$ & $(r//r)$ \\
\hline
Gamètes & 100 \% $(L)$ & 100 \% $(r)$ \\
\hline
\end{tabularx}
\end{center}

La fécondation donne des zygotes $(L//r)$, tous de phénotype $[L]$ : c'est l'\textbf{uniformité de la F1}.

\begin{propriete}[Première loi de Mendel : uniformité des hybrides F1]
Le croisement de deux individus de race pure différant par un caractère donne une première génération F1 \textbf{homogène} : tous les hybrides présentent le phénotype du parent portant l'allèle dominant. Cette loi, généralisation des résultats expérimentaux de Mendel, est admise ; elle s'explique par la dominance d'un allèle sur l'autre chez l'hétérozygote.
\end{propriete}

\courssub{La F2 : la réapparition du caractère récessif et le rapport 3/4 -- 1/4}

L'autofécondation des hybrides F1 $(L//r)$ est l'étape décisive. Chaque F1 produit, par méiose, deux types de gamètes en proportions égales : 50 \% $(L)$ et 50 \% $(r)$. La rencontre au hasard de ces gamètes lors de la fécondation se prévoit avec un \textbf{échiquier de croisement}.

\begin{popfigure}
\begin{center}
\begin{tikzpicture}[scale=1.05]
% Grille échiquier
\draw[thick, popdark] (0,0) rectangle (4,4);
\draw[thick, popdark] (2,0) -- (2,4);
\draw[thick, popdark] (0,2) -- (4,2);
% Gamètes en-têtes
\node at (-1.1,3) {\small $(L)$};
\node at (-1.1,1) {\small $(r)$};
\node at (-1.1,4.55) {\small Gamètes mâles};
\node at (1,4.55) {\small $(L)$};
\node at (3,4.55) {\small $(r)$};
\node at (2,5.25) {\small Gamètes femelles};
% Génotypes des cases
\node[popblue] at (1,3.2) {$(L//L)$};
\node[popblue] at (3,3.2) {$(L//r)$};
\node[popblue] at (1,1.2) {$(L//r)$};
\node[poporange] at (3,1.2) {$(r//r)$};
% Phénotypes des cases
\node[popblue] at (1,2.6) {\footnotesize $[L]$};
\node[popblue] at (3,2.6) {\footnotesize $[L]$};
\node[popblue] at (1,0.6) {\footnotesize $[L]$};
\node[poporange] at (3,0.6) {\footnotesize $[r]$};
\end{tikzpicture}
\end{center}
\legende{Échiquier de croisement de l'autofécondation des hybrides F1 $(L//r) \times (L//r)$. Chaque case a la même probabilité $\frac{1}{4}$. On obtient $\frac{3}{4}$ de graines lisses $[L]$ (en bleu) et $\frac{1}{4}$ de graines ridées $[r]$ (en orange).}
\end{popfigure}

Les résultats théoriques sont donc : $\frac{1}{4}\,(L//L)$, $\frac{2}{4}\,(L//r)$, $\frac{1}{4}\,(r//r)$, soit en phénotypes $\frac{3}{4}\,[L]$ et $\frac{1}{4}\,[r]$.

\begin{exemplebox}[Vérification sur les données réelles de Mendel]
Sur les \num{7324} graines de la F2, Mendel compte \num{5474} graines lisses et \num{1850} graines ridées. Le rapport observé vaut :
\[ \frac{5474}{1850} \approx 2,96 \approx 3 \]
Les proportions théoriques auraient donné $\frac{3}{4} \times 7324 = \num{5493}$ graines lisses et $\frac{1}{4} \times 7324 = \num{1831}$ graines ridées. L'écart avec les valeurs observées (19 graines de chaque catégorie) est négligeable : les résultats expérimentaux confirment le rapport 3/4 -- 1/4. Plus l'échantillon est grand, plus les proportions observées se rapprochent des proportions théoriques.
\end{exemplebox}

\begin{propriete}[Deuxième loi de Mendel : ségrégation des allèles]
Les deux allèles d'un même gène, réunis chez l'hybride F1, \textbf{se séparent (disjoignent) au cours de la formation des gamètes} : chaque gamète ne reçoit qu'un seul allèle. L'autofécondation des F1 donne en F2 les proportions phénotypiques $\frac{3}{4}$ dominant -- $\frac{1}{4}$ récessif (et génotypiques $\frac{1}{4}$ -- $\frac{2}{4}$ -- $\frac{1}{4}$). Cette loi, appelée aussi « loi de pureté des gamètes », est admise comme généralisation des faits expérimentaux.
\end{propriete}

\begin{methode}[Rédiger un échiquier de croisement sans erreur]
\begin{enumerate}
\item \textbf{Identifier} les phénotypes des parents et vérifier qu'il s'agit de races pures (homozygotes) ou d'hybrides (hétérozygotes).
\item \textbf{Écrire} les génotypes avec la convention $(L//r)$ et préciser la dominance si elle est connue.
\item \textbf{Déterminer les gamètes} produits par chaque parent (un hétérozygote donne 50 \% -- 50 \% ; un homozygote, 100 \%).
\item \textbf{Construire l'échiquier} : gamètes d'un parent en ligne, de l'autre en colonne ; remplir chaque case par le génotype du zygote.
\item \textbf{Regrouper} les génotypes identiques, puis déduire les proportions génotypiques et enfin phénotypiques.
\item \textbf{Conclure} en comparant, si l'énoncé le demande, avec les résultats expérimentaux chiffrés.
\end{enumerate}
\end{methode}

\courssub{Le croisement-test : révéler le génotype caché}

Un problème pratique se pose à l'éleveur comme au généticien : un individu de phénotype dominant $[L]$ est-il homozygote $(L//L)$ ou hétérozygote $(L//r)$ ? Son apparence ne le dit pas. La solution est le croisement-test.

\begin{definition}[Croisement-test (test-cross)]
Le \cle{croisement-test} consiste à croiser l'individu de génotype inconnu (phénotype dominant) avec un \textbf{individu récessif homozygote} $(r//r)$, appelé \emph{testeur}. L'analyse des proportions de la descendance révèle le génotype de l'individu testé.
\end{definition}

Raisonnons. Le testeur $(r//r)$ ne produit que des gamètes $(r)$ : il ne peut « masquer » aucun allèle apporté par l'autre parent. Chaque descendant révèle donc directement le gamète reçu de l'individu testé.

\begin{itemize}
\item Si l'individu testé est $(L//L)$ : 100 \% des descendants sont $(L//r)$, donc $[L]$ — descendance homogène.
\item Si l'individu testé est $(L//r)$ : 50 \% $(L//r)$ $[L]$ et 50 \% $(r//r)$ $[r]$ — descendance en proportions 1/2 -- 1/2.
\end{itemize}

\begin{popfigure}
\begin{center}
\begin{tikzpicture}[
  box/.style={draw=popdark, thick, rounded corners=2pt, align=center, minimum width=2.6cm, minimum height=1cm},
  res/.style={draw, thick, rounded corners=2pt, align=center, minimum width=3.1cm, minimum height=0.9cm}
]
% Hypothèse 1
\node[box, fill=popblueL, align=center] (a) at (0,3){Individu testé\\ $(L//L)$ ?};
\node[box, fill=poporangeL, align=center] (b) at (4.2,3){Testeur\\ $(r//r)$};
\draw[->, thick, popdark] (a) -- (b) node[midway, above]{\small croisement};
\node[res, fill=popgreenL, align=center] (r1) at (2.1,1.2){100 \% $(L//r)$ $[L]$\\ descendance homogène};
\draw[->, thick, popdark] (2.1,2.5) -- (r1);
% Hypothèse 2
\node[box, fill=popblueL, align=center] (c) at (0,-1.2){Individu testé\\ $(L//r)$ ?};
\node[box, fill=poporangeL, align=center] (d) at (4.2,-1.2){Testeur\\ $(r//r)$};
\draw[->, thick, popdark] (c) -- (d) node[midway, above]{\small croisement};
\node[res, fill=poppinkL, align=center] (r2) at (2.1,-3){50 \% $(L//r)$ $[L]$ \quad 50 \% $(r//r)$ $[r]$\\ proportions 1/2 -- 1/2};
\draw[->, thick, popdark] (2.1,-1.7) -- (r2);
\end{tikzpicture}
\end{center}
\legende{Principe du croisement-test. Le testeur récessif homozygote $(r//r)$ révèle les gamètes produits par l'individu testé : une descendance homogène indique un homozygote dominant ; une descendance 1/2 -- 1/2 indique un hétérozygote.}
\end{popfigure}

\begin{exoresolu}[Génotype d'une poule à plumage noir]
Chez la poule, l'allèle $N$ (plumage noir) domine l'allèle $b$ (plumage blanc). Un éleveur possède une poule noire et veut connaître son génotype. Il la croise avec un coq blanc. Sur 40 poussins, on compte 19 noirs et 21 blancs. Quel est le génotype de la poule ?
\tcblower
\textbf{Solution.} Le coq blanc a le phénotype récessif $[b]$, donc le génotype $(b//b)$ : c'est un croisement-test. La poule noire est soit $(N//N)$, soit $(N//b)$.

Si elle était $(N//N)$, tous les poussins seraient $(N//b)$, donc noirs : descendance homogène.

Or la descendance est mixte : $\frac{19}{40} \approx 47,5\,\%$ de noirs et $\frac{21}{40} \approx 52,5\,\%$ de blancs, proportions voisines de 1/2 -- 1/2. La poule a donc fourni deux types de gamètes $(N)$ et $(b)$ : elle est \textbf{hétérozygote $(N//b)$}. L'échiquier le confirme : gamètes $(N)$, $(b)$ $\times$ gamète $(b)$ donnent $\frac{1}{2}\,(N//b)\,[N]$ et $\frac{1}{2}\,(b//b)\,[b]$.
\end{exoresolu}

\courssub{Le fondement chromosomique : la méiose explique les lois de Mendel}

Pourquoi les allèles se séparent-ils dans les gamètes ? La réponse est cellulaire : elle tient au comportement des chromosomes au cours de la méiose.

Les deux allèles d'un gène occupent le même locus sur deux \textbf{chromosomes homologues} : l'un d'origine paternelle, l'autre d'origine maternelle. Or, en \textbf{anaphase I} de méiose, les chromosomes homologues se séparent et migrent vers des pôles opposés de la cellule. Chaque gamète ne reçoit donc qu'\textbf{un seul chromosome de chaque paire}, et par conséquent \textbf{un seul allèle de chaque gène} : c'est exactement la deuxième loi de Mendel, la « pureté des gamètes ».

La fécondation, en réunissant au hasard deux gamètes haploïdes, restaure la diploïdie et crée des combinaisons alléliques nouvelles : le zygote $(L//r)$ n'existait chez aucun des deux parents de race pure. Méiose et fécondation sont ainsi les deux moteurs du brassage génétique, première source de la diversité des individus au sein d'une espèce — le cœur de cette leçon.

\begin{aretenir}[L'essentiel de la section]
\begin{itemize}
\item Le monohybridisme étudie la transmission d'un seul caractère ; les parents de race pure sont homozygotes.
\item F1 homogène (1\textsuperscript{re} loi de Mendel) : l'allèle dominant masque l'allèle récessif chez l'hétérozygote.
\item F2 en proportions 3/4 -- 1/4 (2\textsuperscript{e} loi de Mendel) : les allèles disjoignent lors de la formation des gamètes, car les chromosomes homologues se séparent en anaphase I de méiose.
\item Le croisement-test (individu dominant $\times$ récessif homozygote) révèle le génotype : descendance homogène $\Rightarrow$ homozygote ; descendance 1/2 -- 1/2 $\Rightarrow$ hétérozygote.
\end{itemize}
\end{aretenir}

\begin{savaistu}[Une découverte ignorée pendant 35 ans]
Gregor Mendel présente ses résultats en 1865 à la Société des sciences de Brünn (actuelle Brno) : personne n'y prête attention, et il meurt en 1884 sans avoir été reconnu. Ce n'est qu'en 1900 que trois botanistes — le Néerlandais Hugo de Vries, l'Allemand Carl Correns et l'Autrichien Erich von Tschermak — redécouvrent \textbf{indépendamment} ses lois en menant chacun leurs propres croisements. Mendel est aujourd'hui considéré comme le père de la génétique, et ses pois du jardin du monastère ont changé à jamais notre compréhension de l'hérédité — y compris celle des maïs hybrides du marché de Ngaoundéré.
\end{savaistu}

\coursec{Dihybridisme et trihybridisme : le brassage interchromosomique}

Dans la section précédente, nous avons étudié le \cle{monohybridisme} : le croisement d'individus ne différant que par \emph{un seul} caractère. Les deux premières lois de Mendel (uniformité des hybrides F1, disjonction des allèles en F2) en découlent. Mais un organisme ne se résume pas à un caractère ! Un cultivateur de maïs à Ebolowa s'intéresse à la fois à la couleur des grains \emph{et} à la taille des plants ; un éleveur de volailles à Ngaoundéré sélectionne à la fois le plumage \emph{et} la forme de la crête. Que deviennent les lois de Mendel lorsque \emph{deux} caractères — voire trois — sont transmis simultanément ? Se transmettent-ils « en bloc » ou indépendamment ? C'est toute la question du \cle{dihybridisme}, dont l'étude va révéler un mécanisme fondamental de la diversité génétique : le \cle{brassage interchromosomique}.

\courssub{Le dihybridisme : définition et expérience historique de Mendel}

\begin{definition}[Dihybridisme]
Le \cle{dihybridisme} est l'étude de la transmission de \textbf{deux couples d'allèles} (deux caractères) au cours d'un croisement entre individus de races pures différant par ces deux caractères. Les hybrides F1 obtenus sont des \textbf{dihybrides} (doublement hétérozygotes), de génotype $A/a \;;\; B/b$.
\end{definition}

\begin{experience}[L'expérience de Mendel sur le pois (1865)]
\textbf{Matériel et protocole.} Mendel croise deux lignées pures de pois (\emph{Pisum sativum}) :
\begin{itemize}
\item des plants à graines \textbf{jaunes et lisses} ;
\item des plants à graines \textbf{vertes et ridées}.
\end{itemize}
Les hybrides F1 sont ensuite croisés entre eux (autofécondation) pour donner la génération F2.

\textbf{Observations.}
\begin{itemize}
\item En F1 : \textbf{toutes} les graines sont jaunes et lisses. Les allèles « jaune » ($J$) et « lisse » ($L$) sont dominants ; « vert » ($j$) et « ridé » ($l$) sont récessifs. C'est la première loi de Mendel (uniformité des hybrides F1), vérifiée pour chaque caractère.
\item En F2 : Mendel dénombre 556 graines réparties en \textbf{quatre classes phénotypiques} : 315 jaunes-lisses, 108 vertes-lisses, 101 jaunes-ridées, 32 vertes-ridées, soit des proportions très proches de :
\[ \frac{9}{16} \; [J\;L] \;;\quad \frac{3}{16} \; [j\;L] \;;\quad \frac{3}{16} \; [J\;l] \;;\quad \frac{1}{16} \; [j\;l] \]
\end{itemize}

\textbf{Interprétation.} Deux classes phénotypiques sont \emph{nouvelles} : vertes-lisses et jaunes-ridées. Ces combinaisons inédites prouvent que les deux couples d'allèles se sont \textbf{recombinés} : ils ne se transmettent pas en bloc. De plus, pour chaque caractère pris isolément, on retrouve le rapport $3/4 - 1/4$ du monohybridisme : par exemple, les graines jaunes représentent $\frac{315+101}{556} = \frac{416}{556} \approx 74{,}8\,\% \approx \frac{3}{4}$ de l'effectif. Or $\frac{9}{16} = \frac{3}{4} \times \frac{3}{4}$, $\frac{3}{16} = \frac{3}{4} \times \frac{1}{4}$, etc. : les proportions F2 sont le \emph{produit} des proportions de deux monohybridismes indépendants.

\textbf{Conclusion.} Chaque dihybride F1 produit \textbf{quatre types de gamètes équiprobables} : $JL$, $Jl$, $jL$, $jl$, chacun avec la probabilité $1/4$.
\end{experience}

\courssub{L'échiquier de croisement : dénombrement des phénotypes et génotypes}

La fécondation réunit au hasard un gamète mâle et un gamète femelle. Avec 4 types de gamètes de chaque côté, il y a $4 \times 4 = 16$ combinaisons équiprobables, résumées dans l'\cle{échiquier de croisement}.

\begin{popfigure}
\begin{center}
\begin{tikzpicture}[scale=0.92, every node/.style={font=\small}]
% Grille
\foreach \i in {0,...,4} {
  \draw[popline] (\i,0) -- (\i,-4);
  \draw[popline] (0,-\i) -- (4,-\i);
}
% En-têtes gamètes
\node[font=\small\bfseries, popblue] at (0.5,0.35) {$JL$};
\node[font=\small\bfseries, popblue] at (1.5,0.35) {$Jl$};
\node[font=\small\bfseries, popblue] at (2.5,0.35) {$jL$};
\node[font=\small\bfseries, popblue] at (3.5,0.35) {$jl$};
\node[font=\small\bfseries, popblue, rotate=90] at (-0.35,-0.5) {$JL$};
\node[font=\small\bfseries, popblue, rotate=90] at (-0.35,-1.5) {$Jl$};
\node[font=\small\bfseries, popblue, rotate=90] at (-0.35,-2.5) {$jL$};
\node[font=\small\bfseries, popblue, rotate=90] at (-0.35,-3.5) {$jl$};
% Coloration des cases : 9 [JL] popgreenL ; 3 [jL] poporangeL ; 3 [Jl] poppurpleL ; 1 [jl] poppinkL
% Ligne 1 (JL x ...) : tout dominant -> vert
\fill[popgreenL] (0,-1) rectangle (1,0); \fill[popgreenL] (1,-1) rectangle (2,0);
\fill[popgreenL] (2,-1) rectangle (3,0); \fill[popgreenL] (3,-1) rectangle (4,0);
% Ligne 2 (Jl x ...)
\fill[popgreenL] (0,-2) rectangle (1,-1); \fill[poppurpleL] (1,-2) rectangle (2,-1);
\fill[popgreenL] (2,-2) rectangle (3,-1); \fill[poppurpleL] (3,-2) rectangle (4,-1);
% Ligne 3 (jL x ...)
\fill[popgreenL] (0,-3) rectangle (1,-2); \fill[popgreenL] (1,-3) rectangle (2,-2);
\fill[poporangeL] (2,-3) rectangle (3,-2); \fill[poporangeL] (3,-3) rectangle (4,-2);
% Ligne 4 (jl x ...)
\fill[popgreenL] (0,-4) rectangle (1,-3); \fill[poppurpleL] (1,-4) rectangle (2,-3);
\fill[poporangeL] (2,-4) rectangle (3,-3); \fill[poppinkL] (3,-4) rectangle (4,-3);
% Génotypes
\node at (0.5,-0.5) {$J/J\,L/L$}; \node at (1.5,-0.5) {$J/J\,L/l$}; \node at (2.5,-0.5) {$J/j\,L/L$}; \node at (3.5,-0.5) {$J/j\,L/l$};
\node at (0.5,-1.5) {$J/J\,L/l$}; \node at (1.5,-1.5) {$J/J\,l/l$}; \node at (2.5,-1.5) {$J/j\,L/l$}; \node at (3.5,-1.5) {$J/j\,l/l$};
\node at (0.5,-2.5) {$J/j\,L/L$}; \node at (1.5,-2.5) {$J/j\,L/l$}; \node at (2.5,-2.5) {$j/j\,L/L$}; \node at (3.5,-2.5) {$j/j\,L/l$};
\node at (0.5,-3.5) {$J/j\,L/l$}; \node at (1.5,-3.5) {$J/j\,l/l$}; \node at (2.5,-3.5) {$j/j\,L/l$}; \node at (3.5,-3.5) {$j/j\,l/l$};
% Légende couleurs
\fill[popgreenL] (5.0,-0.2) rectangle (5.4,-0.6); \node[right, font=\small] at (5.5,-0.4) {$[J\;L]$ : 9/16};
\fill[poporangeL] (5.0,-1.0) rectangle (5.4,-1.4); \node[right, font=\small] at (5.5,-1.2) {$[j\;L]$ : 3/16};
\fill[poppurpleL] (5.0,-1.8) rectangle (5.4,-2.2); \node[right, font=\small] at (5.5,-2.0) {$[J\;l]$ : 3/16};
\fill[poppinkL] (5.0,-2.6) rectangle (5.4,-3.0); \node[right, font=\small] at (5.5,-2.8) {$[j\;l]$ : 1/16};
\end{tikzpicture}
\end{center}
\legende{Échiquier de croisement du dihybridisme F1 $\times$ F1 ($J/j\,;\,L/l \times J/j\,;\,L/l$). En ligne et en colonne : les 4 types de gamètes équiprobables produits par chaque dihybride. Les 16 cases donnent 9 génotypes différents regroupés en 4 classes phénotypiques dans le rapport 9 -- 3 -- 3 -- 1.}
\end{popfigure}

\begin{aretenir}[Résultats caractéristiques du dihybridisme avec gènes indépendants]
Pour un croisement F1 $\times$ F1 avec dominance totale aux deux locus :
\begin{itemize}
\item \textbf{4 phénotypes} en F2 dans les proportions $\dfrac{9}{16} - \dfrac{3}{16} - \dfrac{3}{16} - \dfrac{1}{16}$ ;
\item \textbf{9 génotypes} différents : $J/J\,L/L$ (1), $J/J\,L/l$ (2), $J/j\,L/L$ (2), $J/j\,L/l$ (4), $J/J\,l/l$ (1), $J/j\,l/l$ (2), $j/j\,L/L$ (1), $j/j\,L/l$ (2), $j/j\,l/l$ (1) ;
\item \textbf{2 phénotypes nouveaux} (recombinés) apparaissent, absents des parents de race pure.
\end{itemize}
\end{aretenir}

\courssub{Le croisement-test en dihybridisme}

Comment vérifier expérimentalement qu'un dihybride produit bien 4 gamètes équiprobables ? On réalise un \cle{croisement-test} : on croise le dihybride F1 avec un \textbf{double récessif} $j/j\,;\,l/l$, qui ne produit qu'un seul type de gamète ($jl$). Le phénotype de chaque descendant révèle alors \emph{directement} le gamète fourni par le dihybride.

\begin{center}
\begin{tabularx}{\linewidth}{|l|X|X|X|X|}
\hline
\rowcolor{popblueL} Gamète du dihybride & $JL$ & $Jl$ & $jL$ & $jl$ \\
\hline
Génotype du descendant & $J/j\,;\,L/l$ & $J/j\,;\,l/l$ & $j/j\,;\,L/l$ & $j/j\,;\,l/l$ \\
\hline
Phénotype & $[J\,L]$ & $[J\,l]$ & $[j\,L]$ & $[j\,l]$ \\
\hline
Proportion & $1/4$ & $1/4$ & $1/4$ & $1/4$ \\
\hline
\end{tabularx}
\end{center}

\begin{methode}[Reconnaître l'indépendance de deux gènes]
\begin{enumerate}
\item Repérer dans l'énoncé un \textbf{croisement-test} (individu double hétérozygote $\times$ double récessif) ou un croisement F1 $\times$ F1.
\item Calculer les \textbf{proportions} de chaque classe phénotypique (effectif de la classe / effectif total).
\item Conclure :
\begin{itemize}
\item croisement-test donnant \textbf{4 classes équiprobables ($\approx 25\,\%$ chacune)} $\Rightarrow$ gènes \textbf{indépendants} (portés par des chromosomes différents) ;
\item croisement F1 $\times$ F1 donnant \textbf{9 -- 3 -- 3 -- 1} $\Rightarrow$ gènes indépendants ;
\item toute autre proportion (2 classes majoritaires, classes inégales...) $\Rightarrow$ gènes \textbf{liés} (voir section suivante).
\end{itemize}
\end{enumerate}
\end{methode}

\begin{exoresolu}[Croisement-test chez la souris]
Chez la souris, l'allèle $N$ (pelage noir) domine $n$ (brun) et l'allèle $R$ (poils raides) domine $r$ (poils frisés). On croise une souris double hétérozygote $N/n\,;\,R/r$ avec une souris brune à poils frisés. On obtient : 48 noires-raides, 52 noires-frisées, 51 brunes-raides, 49 brunes-frisées. Interpréter.
\tcblower
\textbf{Solution.} Le parent brun frisé est double récessif $n/n\,;\,r/r$ : c'est un croisement-test. Effectif total : $48+52+51+49 = 200$. Proportions : $48/200 = 24\,\%$ ; $52/200 = 26\,\%$ ; $51/200 = 25{,}5\,\%$ ; $49/200 = 24{,}5\,\%$. Les 4 classes sont \textbf{équiprobables} ($\approx 25\,\%$ chacune). La souris dihybride a donc produit 4 types de gamètes en quantités égales ($NR$, $Nr$, $nR$, $nr$) : les deux gènes \textbf{ségrègent indépendamment}, ils sont portés par des paires de chromosomes différentes. Il y a eu brassage interchromosomique.
\end{exoresolu}

\courssub{La troisième loi de Mendel et son explication cytologique}

\begin{propriete}[Troisième loi de Mendel : ségrégation indépendante des caractères]
Lorsque deux gènes sont portés par des \textbf{paires de chromosomes différentes}, leurs allèles se répartissent \textbf{indépendamment} les uns des autres dans les gamètes au cours de la méiose. Un dihybride $A/a\,;\,B/b$ produit alors 4 types de gamètes équiprobables : $AB$, $Ab$, $aB$, $ab$ ($1/4$ chacun). Cette loi, généralisation des expériences de Mendel, n'est valable que pour des gènes \emph{non liés}.
\end{propriete}

D'où vient cette indépendance ? De la \textbf{méiose}. En \textbf{métaphase I}, chaque paire de chromosomes homologues se place sur la plaque équatoriale avec une orientation \textbf{aléatoire}, indépendante de l'orientation des autres paires ; en \textbf{anaphase I}, les homologues de chaque paire migrent ensuite vers des pôles opposés. Pour 2 paires de chromosomes, il existe donc \textbf{2 configurations également probables}, donnant au total 4 combinaisons alléliques : c'est le \cle{brassage interchromosomique}.

\begin{definition}[Brassage interchromosomique]
Le \cle{brassage interchromosomique} est le mécanisme de diversification génétique dû à la \textbf{séparation aléatoire et indépendante des chromosomes homologues} en anaphase I de méiose. Il produit des gamètes porteurs de \textbf{nouvelles combinaisons d'allèles} situés sur des chromosomes différents. Pour $n$ paires de chromosomes, le nombre de types de gamètes possibles est :
\[ 2^{n} \]
\end{definition}

\begin{popfigure}
\begin{center}
\begin{tikzpicture}[scale=0.85, >=Stealth]
% Configuration 1
\node[font=\small\bfseries, popdark] at (2,3.6) {Configuration 1};
\draw[popline, thick] (0,0) ellipse (2 and 1.1);
% Paire 1 : A (bleu) à gauche, a (bleu clair) à droite
\draw[popblue, very thick] (0.6,1.6) -- (0.6,0.2); \draw[popblue, very thick] (0.75,1.6) -- (0.75,0.2);
\node[popblue, font=\small\bfseries] at (0.35,1.9) {$A$};
\draw[popblue!45, very thick] (1.3,1.6) -- (1.3,0.2); \draw[popblue!45, very thick] (1.45,1.6) -- (1.45,0.2);
\node[popblue!60, font=\small\bfseries] at (1.7,1.9) {$a$};
% Paire 2 : B (orange) à gauche, b à droite
\draw[poporange, very thick] (0.6,-0.2) -- (0.6,-1.6); \draw[poporange, very thick] (0.75,-0.2) -- (0.75,-1.6);
\node[poporange, font=\small\bfseries] at (0.35,-1.9) {$B$};
\draw[poporange!45, very thick] (1.3,-0.2) -- (1.3,-1.6); \draw[poporange!45, very thick] (1.45,-0.2) -- (1.45,-1.6);
\node[poporange!70, font=\small\bfseries] at (1.7,-1.9) {$b$};
% Flèches de migration
\draw[->, popdark, thick] (0.67,1.9) -- (0.67,2.7);
\draw[->, popdark, thick] (1.37,1.9) -- (1.37,2.7);
\draw[->, popdark, thick] (0.67,-1.9) -- (0.67,-2.7);
\draw[->, popdark, thick] (1.37,-1.9) -- (1.37,-2.7);
\node[font=\small, popdark] at (1,3.1) {pôle 1};
\node[font=\small, popdark] at (1,-3.1) {pôle 2};
% Résultat config 1
\node[font=\small\bfseries, popgreen] at (4.6,1.6) {$\Rightarrow$ gamètes $AB$};
\node[font=\small\bfseries, popgreen] at (4.6,-1.6) {$\Rightarrow$ gamètes $ab$};

% Configuration 2
\begin{scope}[xshift=8.5cm]
\node[font=\small\bfseries, popdark] at (2,3.6) {Configuration 2};
\draw[popline, thick] (0,0) ellipse (2 and 1.1);
\draw[popblue, very thick] (0.6,1.6) -- (0.6,0.2); \draw[popblue, very thick] (0.75,1.6) -- (0.75,0.2);
\node[popblue, font=\small\bfseries] at (0.35,1.9) {$A$};
\draw[popblue!45, very thick] (1.3,1.6) -- (1.3,0.2); \draw[popblue!45, very thick] (1.45,1.6) -- (1.45,0.2);
\node[popblue!60, font=\small\bfseries] at (1.7,1.9) {$a$};
% Paire 2 inversée : b à gauche, B à droite
\draw[poporange!45, very thick] (0.6,-0.2) -- (0.6,-1.6); \draw[poporange!45, very thick] (0.75,-0.2) -- (0.75,-1.6);
\node[poporange!70, font=\small\bfseries] at (0.35,-1.9) {$b$};
\draw[poporange, very thick] (1.3,-0.2) -- (1.3,-1.6); \draw[poporange, very thick] (1.45,-0.2) -- (1.45,-1.6);
\node[poporange, font=\small\bfseries] at (1.7,-1.9) {$B$};
\draw[->, popdark, thick] (0.67,1.9) -- (0.67,2.7);
\draw[->, popdark, thick] (1.37,1.9) -- (1.37,2.7);
\draw[->, popdark, thick] (0.67,-1.9) -- (0.67,-2.7);
\draw[->, popdark, thick] (1.37,-1.9) -- (1.37,-2.7);
\node[font=\small, popdark] at (1,3.1) {pôle 1};
\node[font=\small, popdark] at (1,-3.1) {pôle 2};
\node[font=\small\bfseries, popgreen] at (4.6,1.6) {$\Rightarrow$ gamètes $Ab$};
\node[font=\small\bfseries, popgreen] at (4.6,-1.6) {$\Rightarrow$ gamètes $aB$};
\end{scope}
\end{tikzpicture}
\end{center}
\legende{Brassage interchromosomique en anaphase I de méiose chez un dihybride $A/a\,;\,B/b$. Les deux paires de chromosomes homologues (paire 1 portant $A/a$, paire 2 portant $B/b$) s'orientent au hasard en métaphase I : deux configurations équiprobables. Au total, 4 types de gamètes équiprobables : $AB$, $ab$, $Ab$, $aB$.}
\end{popfigure}

\courssub{Extension au trihybridisme}

\begin{definition}[Trihybridisme]
Le \cle{trihybridisme} est l'étude de la transmission simultanée de \textbf{trois couples d'allèles} portés par trois paires de chromosomes différentes. Le trihybride F1, de génotype $A/a\,;\,B/b\,;\,C/c$, produit $2^{3} = 8$ types de gamètes équiprobables.
\end{definition}

Le raisonnement se généralise sans difficulté : si chaque paire d'allèles ségrège indépendamment, les proportions s'obtiennent par \emph{produit} des probabilités.

\begin{itemize}
\item \textbf{Gamètes du trihybride} : $ABC$, $ABc$, $AbC$, $Abc$, $aBC$, $aBc$, $abC$, $abc$, chacun avec la probabilité $1/8$.
\item \textbf{Croisement F1 $\times$ F1} : l'échiquier comporterait $8 \times 8 = 64$ cases. Les proportions F2 sont :
\[ \frac{27}{64}\,[A\,B\,C] \;;\quad 3 \times \frac{9}{64}\,[A\,B\,c],\ [A\,b\,C],\ [a\,B\,C] \;;\quad 3 \times \frac{3}{64}\,[A\,b\,c],\ [a\,B\,c],\ [a\,b\,C] \;;\quad \frac{1}{64}\,[a\,b\,c] \]
On vérifie que $\frac{27}{64} = \left(\frac{3}{4}\right)^{3}$ et $\frac{1}{64} = \left(\frac{1}{4}\right)^{3}$ : chaque proportion est le produit des probabilités de trois monohybridismes indépendants.
\item \textbf{Croisement-test} (trihybride $\times$ triple récessif) : \textbf{8 classes phénotypiques équiprobables} à $1/8 = 12{,}5\,\%$ chacune.
\end{itemize}

\begin{exemplebox}[Exemples classiques du programme]
\begin{itemize}
\item \textbf{Chez la volaille} : la forme de la crête (crête « rose » / crête « simple ») et la couleur du plumage peuvent ségréger indépendamment ; un croisement F1 $\times$ F1 donne le rapport 9 -- 3 -- 3 -- 1.
\item \textbf{Chez les plantes à fleurs} : chez le pois comme chez le maïs cultivé au Cameroun, couleur des grains, forme des grains et taille des plants se recombinent librement lorsque les gènes sont sur des chromosomes différents — c'est le principe utilisé par les sélectionneurs pour créer des variétés améliorées.
\item \textbf{Chez la drosophile} : les gènes situés sur des paires de chromosomes différentes ségrègent indépendamment ; ceux portés par le même chromosome sont liés (section suivante).
\end{itemize}
\end{exemplebox}

\begin{popfigure}
\begin{center}
\begin{tikzpicture}
\begin{axis}[
  ybar, bar width=14pt,
  width=0.95\linewidth, height=6.2cm,
  ymin=0, ymax=62,
  ylabel={Fréquence (\%)},
  symbolic x coords={[AB],[Ab],[aB],[ab]},
  xtick=data,
  legend style={at={(0.5,1.02)}, anchor=south, legend columns=2, font=\small},
  nodes near coords, every node near coord/.append style={font=\scriptsize},
  ymajorgrids, grid style={popline!40}
]
\addplot[fill=popblue!70, draw=popblue] coordinates {([AB],56.25) ([Ab],18.75) ([aB],18.75) ([ab],6.25)};
\addplot[fill=poporange!70, draw=poporange] coordinates {([AB],25) ([Ab],25) ([aB],25) ([ab],25)};
\legend{F2 (F1 $\times$ F1) : 9--3--3--1, Croisement-test : 1--1--1--1}
\end{axis}
\end{tikzpicture}
\end{center}
\legende{Comparaison des fréquences des 4 classes phénotypiques en dihybridisme avec gènes indépendants. En F2 (9/16 = 56,25 \% ; 3/16 = 18,75 \% ; 1/16 = 6,25 \%), la classe double dominante est majoritaire ; en croisement-test, les 4 classes sont équiprobables (25 \%), signature de l'indépendance des gènes.}
\end{popfigure}

\begin{attention}[Erreur fréquente : appliquer le 9 -- 3 -- 3 -- 1 à tort]
Les proportions 9 -- 3 -- 3 -- 1 (F2) et 1 -- 1 -- 1 -- 1 (croisement-test) ne sont valables \textbf{que si les deux gènes sont indépendants}, c'est-à-dire portés par des chromosomes différents. Si les gènes sont \textbf{liés} (même chromosome), le croisement-test donne des classes \emph{inégales} : deux classes parentales majoritaires et, éventuellement, deux classes recombinées minoritaires. \textbf{Avant tout calcul, vérifie l'indépendance des gènes} en examinant les proportions observées — c'est souvent la première question de l'exercice au Baccalauréat !
\end{attention}

\courssub{Conséquence : le brassage interchromosomique, moteur de la diversité génétique}

Chaque paire de chromosomes \emph{double} le nombre de combinaisons possibles : $2^{n}$ types de gamètes pour $n$ paires. Et comme la fécondation réunit deux gamètes pris au hasard, le nombre de combinaisons possibles pour un zygote est considérable.

\begin{savaistu}[8,4 millions de gamètes différents... avant même la fécondation]
L'espèce humaine possède $n = 23$ paires de chromosomes. Par brassage interchromosomique seul, chaque individu peut produire $2^{23} \approx \num{8,4}$ \textbf{millions} de types de gamètes différents. Un couple peut donc, en théorie, engendrer $2^{23} \times 2^{23} = 2^{46} \approx 7 \times 10^{13}$ combinaisons génétiques distinctes — et ce \emph{sans compter} le crossing-over (brassage intrachromosomique, section suivante) ni les mutations ! C'est pourquoi, jumeaux monozygotes exceptés, il n'existera jamais deux êtres humains génétiquement identiques : la méiose et la fécondation garantissent l'\textbf{unicité génétique} de chaque individu tout en maintenant l'unité de l'espèce.
\end{savaistu}

\begin{aretenir}[L'essentiel de la section]
\begin{itemize}
\item Le \textbf{dihybridisme} étudie deux couples d'allèles ; le \textbf{trihybridisme} en étudie trois.
\item Gènes indépendants $\Rightarrow$ F2 en \textbf{9 -- 3 -- 3 -- 1} (dihybridisme) ou \textbf{27 -- 9 -- 9 -- 9 -- 3 -- 3 -- 3 -- 1} (trihybridisme) ; croisement-test en classes \textbf{équiprobables} (4 ou 8).
\item La \textbf{3\textsuperscript{e} loi de Mendel} (ségrégation indépendante) s'explique par la \textbf{séparation aléatoire des chromosomes homologues en anaphase I} : c'est le \textbf{brassage interchromosomique}.
\item Nombre de types de gamètes : $2^{n}$ ($n$ = nombre de paires de chromosomes). Ce brassage multiplie les combinaisons alléliques et contribue puissamment à la \textbf{diversité génétique} au sein de l'espèce.
\end{itemize}
\end{aretenir}

Reste une question en suspens : que se passe-t-il lorsque les deux gènes sont portés par le \emph{même} chromosome ? Les proportions s'écartent alors du 9 -- 3 -- 3 -- 1, et un second mécanisme de brassage entre en jeu : le \emph{crossing-over}. C'est l'objet de la section suivante.

\coursec{Linkage et crossing-over : le brassage intrachromosomique}

Dans la section précédente, nous avons vu que lorsque deux gènes sont portés par des chromosomes différents, la migration aléatoire des chromosomes homologues en anaphase I de méiose réalise un \cle{brassage interchromosomique} : un double hétérozygote $A/a ; B/b$ produit quatre types de gamètes équiprobables ($1/4$ chacun), et le croisement-test donne quatre classes phénotypiques à \SI{25}{\percent}. Mais cette troisième loi de Mendel est-elle universelle ? C'est en croisant des drosophiles que Thomas Hunt Morgan, prix Nobel 1933, découvrit que non : certains gènes refusent de se séparer. Cette section explique pourquoi, et montre comment la méiose invente un second mécanisme de diversification : le \cle{crossing-over}.

\courssub{Mise en évidence expérimentale : quand les résultats contredisent la 3\textsuperscript{ème} loi de Mendel}

\begin{experience}[Le croisement-test de Morgan chez la drosophile]
\textbf{Protocole.} Morgan croise des drosophiles de lignée pure à \textbf{corps gris et ailes longues} (génotype $\frac{G\,L}{G\,L}$) avec des drosophiles à \textbf{corps noir et ailes vestigiales} (génotype $\frac{n\,v}{n\,v}$). La F$_1$ est uniforme : tous les individus sont à corps gris et ailes longues, ce qui confirme la dominance des allèles $G$ (gris) sur $n$ (noir) et $L$ (longues) sur $v$ (vestigiales). Les femelles F$_1$ sont donc double hétérozygotes de disposition \emph{cis} : $\frac{G\,L}{n\,v}$.

Il réalise ensuite un \textbf{croisement-test} : une femelle F$_1$ $\frac{G\,L}{n\,v}$ est croisée avec un mâle double récessif $\frac{n\,v}{n\,v}$ (testeur, qui ne produit que des gamètes $n\,v$).

\textbf{Résultats attendus} si les gènes étaient indépendants (3\textsuperscript{ème} loi de Mendel) : 4 classes équiprobables à \SI{25}{\percent} chacune.

\textbf{Résultats observés} (sur 1171 descendants) :
\begin{center}
\begin{tabularx}{\linewidth}{|l|X|X|}
\hline
\rowcolor{popblueL} \textbf{Phénotype} & \textbf{Effectif} & \textbf{Pourcentage} \\
\hline
[gris, longues] (parental) & 483 & \SI{41,2}{\percent} \\
\hline
[noir, vestigiales] (parental) & 482 & \SI{41,2}{\percent} \\
\hline
[gris, vestigiales] (recombiné) & 103 & \SI{8,8}{\percent} \\
\hline
[noir, longues] (recombiné) & 103 & \SI{8,8}{\percent} \\
\hline
\end{tabularx}
\end{center}

\textbf{Interprétation.} Les proportions ne sont \emph{pas} $1/4-1/4-1/4-1/4$ : les phénotypes des grands-parents (parentaux) sont largement majoritaires (\SI{82,4}{\percent}), et deux classes nouvelles (recombinées) apparaissent en faible proportion (\SI{17,6}{\percent}). La ségrégation indépendante des caractères est mise en défaut : les deux gènes ne sont pas indépendants.
\end{experience}

\begin{definition}[Phénotypes parentaux et phénotypes recombinés]
Dans un croisement-test, on appelle :
\begin{itemize}
\item \cle{phénotypes parentaux} : les phénotypes identiques à ceux des parents de la lignée parentale (ici [gris, longues] et [noir, vestigiales]) ;
\item \cle{phénotypes recombinés} : les phénotypes nouveaux, résultant d'une recombinaison des allèles parentaux (ici [gris, vestigiales] et [noir, longues]).
\end{itemize}
\end{definition}

\courssub{Le linkage : des gènes portés par le même chromosome}

L'hypothèse de Morgan est simple et géniale : si les deux gènes ne se séparent pas librement, c'est qu'ils sont \textbf{portés par le même chromosome}. Le chromosome étant une structure physique indivisible lors de sa migration vers les pôles, les allèles qu'il porte migrent ensemble.

\begin{definition}[Linkage total et linkage partiel]
\begin{itemize}
\item Le \cle{linkage} (ou liaison génétique) est la dépendance de deux ou plusieurs gènes portés par le même chromosome, qui tendent à être transmis ensemble aux gamètes.
\item Le linkage est \cle{total} (ou absolu) lorsque les deux gènes restent toujours associés : l'hétérozygote ne produit que 2 types de gamètes parentaux en proportions égales ($1/2-1/2$), et le croisement-test ne donne que \textbf{2 classes phénotypiques équiprobables}.
\item Le linkage est \cle{partiel} lorsque les deux gènes sont parfois séparés par un crossing-over : l'hétérozygote produit 4 types de gamètes inégaux, et le croisement-test donne \textbf{4 classes inégales}, les parentaux étant majoritaires.
\end{itemize}
\end{definition}

\textbf{Cas du linkage total.} Chez la drosophile, le crossing-over n'existe jamais chez le mâle. Si l'on croise un \emph{mâle} F$_1$ $\frac{G\,L}{n\,v}$ avec une femelle testeuse $\frac{n\,v}{n\,v}$, on n'obtient que deux classes : \SI{50}{\percent} [gris, longues] et \SI{50}{\percent} [noir, vestigiales]. Les gamètes du mâle sont uniquement $G\,L$ et $n\,v$ : la disposition des allèles sur les chromosomes parentaux est intégralement conservée.

\begin{attention}[Disposition cis/trans des allèles : le piège classique]
Pour des gènes liés, il ne suffit pas d'écrire le génotype $G/n ; L/v$ : il faut préciser la \textbf{disposition des allèles sur les chromosomes homologues}. On écrit par exemple $\frac{G\,L}{n\,v}$ (disposition \emph{cis} : les deux dominants sur le même chromosome) ou $\frac{G\,v}{n\,L}$ (disposition \emph{trans} : allèles dominants et récessifs en opposition). Les gamètes parentaux ne sont pas les mêmes dans les deux cas ! Écrire les gamètes sans tenir compte de cette disposition est l'erreur la plus fréquente au Baccalauréat. Déduis toujours la disposition à partir des phénotypes des parents de la lignée pure.
\end{attention}

\courssub{Le crossing-over : mécanisme cytologique du linkage partiel}

Si les gènes sont liés, comment expliquer l'apparition de gamètes recombinés ? La réponse se trouve en \textbf{prophase I de méiose}. Lors de l'appariement des chromosomes homologues (synapsis), chaque chromosome est déjà formé de deux chromatides sœurs. Les quatre chromatides du bivalent (tétrade) s'entrecroisent en des points appelés \cle{chiasmas}, visibles au microscope. Au niveau d'un chiasma, deux chromatides \emph{non sœurs} (appartenant à des homologues différents) peuvent se rompre et échanger réciproquement des segments : c'est le \cle{crossing-over} (ou enjambement).

\begin{definition}[Crossing-over et brassage intrachromosomique]
Le \cle{crossing-over} est l'échange réciproque de segments chromatidiens entre chromatides non sœurs de chromosomes homologues, au niveau des chiasmas, au cours de la prophase I de la méiose.

Il réalise un \cle{brassage intrachromosomique} : il produit de \textbf{nouvelles combinaisons d'allèles sur un même chromosome}, c'est-à-dire des chromatides recombinées portant des associations d'allèles différentes de celles des chromosomes parentaux.
\end{definition}

\begin{popfigure}
\begin{center}
\begin{tikzpicture}[scale=0.9, every node/.style={font=\small}]
% Bivalent avant crossing-over
\node[font=\bfseries] at (2.6,4.6) {1. Bivalent en prophase I (chiasma)};
% chromatides homologue 1 (bleu) - chromosome paternel
\draw[line width=5pt, popblue] (1.2,1.2) -- (1.2,4.2);
\draw[line width=5pt, popblue] (1.7,1.2) -- (1.7,4.2);
% chromatides homologue 2 (orange) - chromosome maternel
\draw[line width=5pt, poporange] (3.5,1.2) -- (3.5,4.2);
\draw[line width=5pt, poporange] (4.0,1.2) -- (4.0,4.2);
% allèles sur chromatides sœurs (identiques)
\node[left, popblue] at (1.2,3.7) {$G$};
\node[left, popblue] at (1.2,2.6) {$L$};
\node[left, popblue] at (1.7,3.7) {$G$};
\node[left, popblue] at (1.7,2.6) {$L$};
\node[right, poporange] at (4.0,3.7) {$n$};
\node[right, poporange] at (4.0,2.6) {$v$};
\node[right, poporange] at (3.5,3.7) {$n$};
\node[right, poporange] at (3.5,2.6) {$v$};
% chiasma entre chromatides non sœurs (ex: 2ème bleue et 1ère orange)
\draw[popdark, thick, decorate, decoration={zigzag, segment length=4, amplitude=1.5}] (1.7,3.15) -- (3.5,3.15);
\node[popdark, above] at (2.6,3.25) {chiasma};
% flèche
\draw[-{Stealth}, very thick, popdark] (4.8,2.7) -- (5.8,2.7) node[midway, above]{échange réciproque};
% Après crossing-over
\node[font=\bfseries] at (9.2,4.6) {2. Chromatides après crossing-over (4 gamètes potentiels)};
% Chromatide 1 (bleue, non croisée) -> parental
\draw[line width=5pt, popblue] (6.6,1.2) -- (6.6,4.2);
\node[below, align=center] at (6.6,1.0) {$G\,L$\\parental};
% Chromatide 2 (bleue puis orange) -> recombiné
\draw[line width=5pt, popblue] (7.1,1.2) -- (7.1,3.15);
\draw[line width=5pt, poporange] (7.1,3.15) -- (7.1,4.2);
\node[below, align=center, poppurple] at (7.1,1.0) {$G\,v$\\recombiné};
% Chromatide 3 (orange puis bleue) -> recombiné
\draw[line width=5pt, poporange] (9.9,1.2) -- (9.9,3.15);
\draw[line width=5pt, popblue] (9.9,3.15) -- (9.9,4.2);
\node[below, align=center, poppurple] at (9.9,1.0) {$n\,L$\\recombiné};
% Chromatide 4 (orange, non croisée) -> parental
\draw[line width=5pt, poporange] (10.4,1.2) -- (10.4,4.2);
\node[below, align=center] at (10.4,1.0) {$n\,v$\\parental};
\end{tikzpicture}
\end{center}
\legende{Schéma du crossing-over en prophase I de méiose. 1 : le bivalent comporte deux chromosomes homologues (bleu : d'origine paternelle, portant $G$ et $L$ ; orange : d'origine maternelle, portant $n$ et $v$), chacun formé de deux chromatides sœurs identiques ; un chiasma se forme entre deux chromatides non sœurs. 2 : après échange réciproque de segments au niveau du chiasma, on obtient deux chromatides parentales ($G\,L$ et $n\,v$) et deux chromatides recombinées ($G\,v$ et $n\,L$). Chaque chromatide aboutira à un gamète après les deux divisions méiotiques.}
\end{popfigure}

\begin{propriete}[Conséquences génétiques du linkage et du crossing-over]
\begin{itemize}
\item Les gènes liés \textbf{ne suivent pas la ségrégation indépendante} (3\textsuperscript{ème} loi de Mendel) : celle-ci ne s'applique qu'aux gènes portés par des chromosomes différents (ou très éloignés sur le même chromosome).
\item Le crossing-over \textbf{sépare partiellement} les gènes liés : plus les deux gènes sont éloignés sur le chromosome, plus la probabilité qu'un chiasma se forme entre eux est grande, et plus la proportion de gamètes recombinés est élevée (toujours strictement inférieure à \SI{50}{\percent}).
\item Ces faits ont été établis expérimentalement par Morgan et son équipe (Sturtevant, Bridges) grâce aux croisements chez la drosophile entre 1910 et 1915.
\end{itemize}
\end{propriete}

\begin{popfigure}
\begin{center}
\begin{tikzpicture}
\begin{axis}[
    ybar, bar width=0.6cm,
    width=0.95\linewidth, height=6.5cm,
    ymin=0, ymax=50,
    ylabel={Fréquence (\%)},
    xtick={1,2,3,4},
    xticklabels={{$[G,L]$ parental},{$[n,v]$ parental},{$[G,v]$ recombinant},{$[n,L]$ recombinant}},
    x tick label style={font=\small, align=center},
    nodes near coords, nodes near coords style={font=\small},
    every axis plot/.append style={fill=popblue, draw=popdark},
    axis lines=left,
    enlarge x limits=0.2,
]
\addplot coordinates {(1,41.2) (2,41.2) (3,8.8) (4,8.8)};
\end{axis}
\end{tikzpicture}
\end{center}
\legende{Résultats du croisement-test de Morgan (femelle F$_1$ $\frac{G\,L}{n\,v}$ $\times$ mâle $\frac{n\,v}{n\,v}$) : les deux classes parentales [gris, longues] et [noir, vestigiales] dominent (\SI{41,2}{\percent} chacune), les deux classes recombinées [gris, vestigiales] et [noir, longues] sont minoritaires (\SI{8,8}{\percent} chacune). Cette dissymétrie est la signature d'un linkage partiel.}
\end{popfigure}

\courssub{Reconnaître un linkage dans un exercice}

\begin{methode}[Diagnostiquer un linkage à partir d'un croisement-test]
On croise un double hétérozygote F$_1$ avec un homozygote récessif (testeur). Les phénotypes des descendants reflètent directement les gamètes produits par la F$_1$.
\begin{enumerate}
\item \textbf{4 classes équiprobables} (\SI{25}{\percent} chacune) $\Rightarrow$ gènes \textbf{indépendants} (brassage interchromosomique seul, gènes sur chromosomes différents).
\item \textbf{2 classes équiprobables} (\SI{50}{\percent} chacune), de phénotypes parentaux $\Rightarrow$ \textbf{linkage total} (pas de crossing-over : ex. mâle drosophile, ou gènes très proches).
\item \textbf{4 classes inégales}, avec deux classes majoritaires (parentales) et deux classes minoritaires (recombinées) $\Rightarrow$ \textbf{linkage partiel} (crossing-over en prophase I).
\item Identifier les classes parentales : ce sont les plus fréquentes ; elles révèlent la disposition \emph{cis} ou \emph{trans} des allèles sur les chromosomes de la F$_1$.
\end{enumerate}
\end{methode}

\begin{exoresolu}[Croisement-test chez la drosophile]
Énoncé. On croise une femelle drosophile de génotype $\frac{G\,L}{n\,v}$ avec un mâle double récessif $\frac{n\,v}{n\,v}$. On obtient 1171 descendants : 483 [gris, longues], 482 [noir, vestigiales], 103 [gris, vestigiales] et 103 [noir, longues]. Interpréter ces résultats.
\tcblower
\textbf{Solution.}
\begin{enumerate}
\item Le testeur ne produit que des gamètes $n\,v$ : les phénotypes des descendants traduisent directement les gamètes de la femelle F$_1$.
\item On observe 4 classes \textbf{inégales} : les classes parentales totalisent $483+482=965$ individus, soit $\frac{965}{1171}\approx\SI{82,4}{\percent}$, et les classes recombinées $103+103=206$ individus, soit $\frac{206}{1171}\approx\SI{17,6}{\percent}$.
\item Les gènes sont donc \textbf{liés} (portés par le même chromosome) avec \textbf{linkage partiel} : dans environ \SI{17,6}{\percent} des méioses, un crossing-over entre les deux loci a produit des chromatides recombinées $G\,v$ et $n\,L$.
\item La disposition des allèles chez la F$_1$ est \emph{cis} : $\frac{G\,L}{n\,v}$, confirmée par les classes majoritaires.
\end{enumerate}
\end{exoresolu}

\courssub{Comparaison des deux brassages et rôle dans la diversité}

\begin{center}
\begin{tabularx}{\linewidth}{|l|X|X|}
\hline
\rowcolor{popblueL} \textbf{Critère} & \textbf{Brassage interchromosomique} & \textbf{Brassage intrachromosomique} \\
\hline
Mécanisme cytologique & Migration aléatoire des homologues en anaphase I & Crossing-over entre chromatides non sœurs en prophase I (chiasmas) \\
\hline
Gènes concernés & Gènes indépendants (chromosomes différents) & Gènes liés (même chromosome) \\
\hline
Gamètes d'un double hétérozygote & 4 types équiprobables ($1/4$ chacun) & 4 types inégaux : parentaux majoritaires, recombinés minoritaires \\
\hline
Croisement-test (phénotypes) & 4 classes à \SI{25}{\percent} & 4 classes inégales (ou 2 classes si linkage total) \\
\hline
Fréquence max des recombinés & \SI{50}{\percent} (indépendance) & $< \SI{50}{\percent}$ (tend vers 50\% si gènes très éloignés) \\
\hline
\end{tabularx}
\end{center}

\begin{aretenir}[Méiose, fécondation et diversité génétique]
Le crossing-over multiplie encore le nombre de combinaisons gamétiques : même pour des gènes liés, la méiose fabrique des chromosomes portant des associations d'allèles inédites. Combiné au brassage interchromosomique ($2^n$ combinaisons pour $n$ paires de chromosomes, soit $2^{23}\approx 8,4 \times 10^6$ chez l'Homme) et à la rencontre aléatoire des gamètes à la fécondation, il garantit que chaque individu (hors vrais jumeaux monozygotes) est génétiquement unique : c'est le fondement de la diversité génétique au sein d'une espèce.
\end{aretenir}

\begin{savaistu}[La drosophile mâle, un cas particulier précieux]
Chez la drosophile mâle, il n'y a \textbf{jamais de crossing-over} : les méioses spermatiques se déroulent sans chiasma. Morgan a exploité cette particularité : un croisement-test utilisant un mâle hétérozygote donne toujours 2 classes équiprobables (linkage total apparent), ce qui permettait de vérifier simplement que deux gènes sont sur le même chromosome. À l'inverse, les femelles, qui réalisent des crossing-over, ont permis à son élève Sturtevant de construire en 1913 les premières cartes génétiques — l'objet de la section suivante.
\end{savaistu}

\coursec{Pourcentage de recombinaison et carte génétique}

Dans la section précédente, nous avons vu que deux gènes portés par le même chromosome ne se transmettent pas toujours ensemble : le \cle{crossing-over}, survenu en prophase I de la méiose, réalise un \cle{brassage intrachromosomique} qui produit des gamètes dits « recombinés ». Une question surgit alors naturellement : peut-on \emph{mesurer} ce brassage ? Et si oui, cette mesure peut-elle renseigner sur la \emph{position} des gènes le long du chromosome ? C'est toute l'ambition de cette section : transformer les proportions obtenues dans les croisements en une véritable \emph{carte} des chromosomes, exploit que réalisèrent Thomas Hunt Morgan et son élève Alfred Sturtevant au début du XX\textsuperscript{e} siècle.

\courssub{Le taux de recombinaison : définition et calcul}

\textbf{Une intuition pour commencer.} Imagine deux villages reliés par une longue route. Plus la route est longue, plus il est probable qu'un voyageur s'arrête en chemin et change de véhicule. De même, plus deux gènes sont éloignés l'un de l'autre sur un chromosome, plus il est probable qu'un crossing-over survienne \emph{entre eux} au cours de la méiose. Le taux de recombinaison mesure précisément cette probabilité, estimée à partir des descendants d'un croisement.

\begin{definition}[Taux (ou pourcentage) de recombinaison]
Le \cle{taux de recombinaison} $p$ entre deux gènes est le pourcentage d'individus \emph{recombinés} (c'est-à-dire présentant une combinaison de caractères différente de celle des parents) parmi l'ensemble des descendants d'un croisement :
\[ p = \frac{\text{nombre d'individus recombinés}}{\text{effectif total des descendants}} \times 100 \]
Il s'exprime en pourcentage (\%).
\end{definition}

\begin{attention}[Croisement-test obligatoire]
Le taux de recombinaison se calcule à partir d'un \cle{croisement-test} (croisement d'un individu double hétérozygote avec un homozygote récessif). En effet, chez le parent testeur récessif, chaque gamète ne porte que des allèles récessifs : le phénotype du descendant révèle donc \emph{directement} la nature du gamète fourni par l'hybride. Les classes phénotypiques \emph{les plus nombreuses} correspondent aux gamètes \emph{parentaux} ; les classes \emph{les moins nombreuses} correspondent aux gamètes \emph{recombinés} (issus d'un crossing-over).
\end{attention}

\begin{exemplebox}[Un premier calcul guidé]
On croise une drosophile double hétérozygote $[AB//ab]$ avec un mâle double récessif $[ab//ab]$. On obtient 400 descendants :
\begin{itemize}
\item 160 drosophiles $[AB]$ (phénotype parental) ;
\item 160 drosophiles $[ab]$ (phénotype parental) ;
\item 45 drosophiles $[Ab]$ (phénotype recombiné) ;
\item 35 drosophiles $[aB]$ (phénotype recombiné).
\end{itemize}
Les classes parentales sont les plus nombreuses (160 + 160 = 320) ; les classes recombinées totalisent $45 + 35 = 80$ individus. D'où :
\[ p = \frac{80}{400} \times 100 = 20\,\% \]
Le taux de recombinaison entre les gènes $A$ et $B$ vaut 20\,\%.
\end{exemplebox}

\courssub{Taux de recombinaison et distance génétique : le centimorgan}

\textbf{Observation fondatrice.} En croisant systématiquement des drosophiles portant différentes paires de gènes liés, Morgan constata que le taux de recombinaison est \emph{constant pour une paire de gènes donnée}, mais \emph{variable d'une paire à l'autre} : 1,3\,\% entre deux gènes, 32,6\,\% entre deux autres, etc. L'hypothèse de Sturtevant fut audacieuse et simple : si le crossing-over est un événement dont la probabilité augmente avec la longueur du segment chromosomique, alors le taux de recombinaison est une \emph{mesure indirecte de la distance} séparant les gènes.

\begin{propriete}[Relation distance--recombinaison (admise)]
Le taux de recombinaison entre deux gènes est \emph{proportionnel à leur distance} sur le chromosome : plus deux gènes sont éloignés, plus la probabilité qu'un crossing-over survienne entre eux est grande, et plus le taux de recombinaison est élevé. Cette propriété, établie expérimentalement par les travaux de Morgan et Sturtevant (1911--1913), est admise au programme ; sa justification (additivité des distances) sera utilisée dans la construction des cartes.
\end{propriete}

\begin{definition}[Le centimorgan]
L'unité de distance génétique est le \cle{centimorgan} (symbole : \cle{cM}), en hommage à Morgan. Par définition :
\[ 1\,\% \text{ de recombinaison} = 1\ \text{cM} \]
Ainsi, deux gènes présentant un taux de recombinaison de 20\,\% sont distants de 20 cM sur la carte génétique.
\end{definition}

\begin{attention}[Le plafond des 50 \%]
Le taux de recombinaison \emph{ne peut jamais dépasser 50\,\%}. En effet, lorsque deux gènes sont très éloignés, les crossing-over deviennent si fréquents (simples, doubles, triples...) que les gamètes recombinés finissent par représenter la moitié des gamètes : les gènes se comportent alors \emph{comme s'ils étaient indépendants} (situation de brassage interchromosomique). Un taux calculé supérieur à 50\,\% dans un exercice signale une erreur d'identification des classes parentales et recombinées.
\end{attention}

\begin{popfigure}
\begin{center}
\begin{tikzpicture}
\begin{axis}[
  width=0.85\linewidth, height=6cm,
  xlabel={Distance r\'eelle entre les g\`enes (unit\'es arbitraires)},
  ylabel={Taux de recombinaison (\%)},
  xmin=0, xmax=100, ymin=0, ymax=60,
  xtick={0,20,40,60,80,100},
  ytick={0,10,20,30,40,50},
  grid=both, grid style={popline!40},
  axis line style={popdark},
]
\addplot[popcurve,domain=0:100,samples=200]{50*(1-exp(-0.02*x))};
\addplot[dashed,poporange,thick,domain=0:100]{50};
\node[poporange,anchor=south west] at (axis cs:52,50.5) {plafond : 50\,\%};
\addplot[dashed,popgreen,thick,domain=0:35]{x};
\node[popgreen,anchor=west,rotate=32] at (axis cs:12,14) {zone de proportionnalit\'e};
\end{axis}
\end{tikzpicture}
\end{center}
\legende{Relation entre la distance s\'eparant deux g\`enes sur le chromosome et le taux de recombinaison observ\'e. Pour des distances faibles, la relation est quasi proportionnelle (droite verte) ; pour des distances grandes, les crossing-over multiples font plafonner le taux \`a 50\,\% (courbe bleue). C'est pourquoi les distances en cM ne sont additives que pour des g\`enes proches.}
\end{popfigure}

\courssub{Construire une carte génétique : cas du dihybridisme}

\begin{definition}[Carte génétique (ou carte factorielle)]
Une \cle{carte génétique} (ou \cle{carte factorielle}) est la représentation schématique de la \emph{position relative} des gènes le long d'un chromosome, les distances étant exprimées en centimorgans à partir des taux de recombinaison mesurés deux à deux.
\end{definition}

Pour deux gènes liés, la construction est immédiate : le taux de recombinaison donne directement la distance. Si $p = 20\,\%$ entre $A$ et $B$, on place les deux loci à 20 cM l'un de l'autre sur un segment représentant le chromosome.

\courssub{Construire une carte génétique : cas du trihybridisme avec gènes liés}

Avec trois gènes liés, le raisonnement se complique mais devient beaucoup plus riche : il faut \emph{ordonner} les gènes avant de calculer les distances.

\begin{methode}[Construire une carte génétique en 4 étapes]
\begin{enumerate}
\item \textbf{Identifier les classes.} Dans les résultats du croisement-test, les deux classes \emph{les plus nombreuses} sont les classes \emph{parentales} ; les deux classes \emph{les plus rares} correspondent aux \emph{doubles crossing-over} (DCO).
\item \textbf{Ordonner les gènes.} Comparer une classe parentale à une classe DCO : le gène dont l'allèle a \emph{changé de compagnon} est le gène \emph{situé au milieu}. (Le double crossing-over ne « déplace » que le gène central.)
\item \textbf{Calculer les distances.} Pour chaque intervalle : 
\[ d = \frac{\text{recombinés dans l'intervalle} + \text{DCO}}{\text{effectif total}} \times 100 \]
Les DCO comptent dans \emph{chacun} des deux intervalles, car chaque double crossing-over implique un échange dans l'intervalle 1 \emph{et} un échange dans l'intervalle 2.
\item \textbf{Tracer la carte.} Placer les trois loci sur un segment gradué en cM et vérifier : la distance entre les gènes extrêmes est approximativement égale à la somme des distances intermédiaires (aux doubles crossing-over près).
\end{enumerate}
\end{methode}

\begin{attention}[Le piège classique du double crossing-over]
Erreur fréquente au Baccalauréat : oublier de compter les individus issus de \emph{doubles crossing-over} dans le calcul des distances des intervalles. Un double crossing-over produit un échange dans \emph{chacun} des deux intervalles : ces individus doivent donc être comptés \emph{deux fois} au total (une fois pour chaque intervalle). Si on les oublie, la distance entre les gènes extrêmes calculée directement (recombinés apparents) sera \emph{inférieure} à la somme des distances intermédiaires, et la carte sera incohérente.
\end{attention}

\begin{exoresolu}[Trois gènes liés chez la drosophile]
Chez la drosophile, on considère trois gènes liés : $vg$ (ailes vestigiales) / $vg^{+}$ (ailes normales) ; $b$ (corps noir) / $b^{+}$ (corps gris) ; $cn$ (yeux cinabre) / $cn^{+}$ (yeux rouges). Une femelle triple hétérozygote est croisée avec un mâle triple récessif (croisement-test). On obtient 1000 descendants :

\begin{center}
\begin{tabularx}{\linewidth}{|l|X|c|}\hline
\rowcolor{popblueL} Phénotype du descendant & Gamète fourni par la femelle & Effectif \\ \hline
ailes normales, corps gris, yeux rouges & $vg^{+}\ b^{+}\ cn^{+}$ & 400 \\ \hline
ailes vestigiales, corps noir, yeux cinabre & $vg\ b\ cn$ & 390 \\ \hline
ailes normales, corps noir, yeux rouges & $vg^{+}\ b\ cn^{+}$ & 60 \\ \hline
ailes vestigiales, corps gris, yeux cinabre & $vg\ b^{+}\ cn$ & 70 \\ \hline
ailes normales, corps gris, yeux cinabre & $vg^{+}\ b^{+}\ cn$ & 35 \\ \hline
ailes vestigiales, corps noir, yeux rouges & $vg\ b\ cn^{+}$ & 25 \\ \hline
ailes normales, corps noir, yeux cinabre & $vg^{+}\ b\ cn$ & 10 \\ \hline
ailes vestigiales, corps gris, yeux rouges & $vg\ b^{+}\ cn^{+}$ & 10 \\ \hline
\end{tabularx}
\end{center}

Déterminer l'ordre des gènes, calculer les distances et établir la carte factorielle.

\tcblower

\textbf{Étape 1 : identification des classes.} Les classes les plus nombreuses (400 et 390) sont parentales : $vg^{+}\ b^{+}\ cn^{+}$ et $vg\ b\ cn$. Les classes les plus rares (10 et 10) sont issues de doubles crossing-over : $vg^{+}\ b\ cn$ et $vg\ b^{+}\ cn^{+}$.

\textbf{Étape 2 : ordre des gènes.} Comparons la classe parentale $vg^{+}\ b^{+}\ cn^{+}$ à la classe DCO $vg^{+}\ b\ cn$ : seuls les allèles du gène $b$ ont changé de compagnons ($b^{+}$ est devenu $b$, associé à $vg^{+}$ et $cn$). Le gène $b$ est donc le gène \emph{central}. L'ordre est : $vg$ --- $b$ --- $cn$.

\textbf{Étape 3 : calcul des distances.}

Intervalle $vg$--$b$ : recombinés = gamètes $vg^{+}\ b\ cn^{+}$ (60) et $vg\ b^{+}\ cn$ (70), plus les DCO (10 + 10) :
\[ d_{vg-b} = \frac{60 + 70 + 10 + 10}{1000} \times 100 = 15\,\% \Rightarrow 15\ \text{cM} \]

Intervalle $b$--$cn$ : recombinés = gamètes $vg^{+}\ b^{+}\ cn$ (35) et $vg\ b\ cn^{+}$ (25), plus les DCO (10 + 10) :
\[ d_{b-cn} = \frac{35 + 25 + 10 + 10}{1000} \times 100 = 8\,\% \Rightarrow 8\ \text{cM} \]

\textbf{Étape 4 : vérification.} Distance directe entre les extrêmes $vg$ et $cn$ (recombinés apparents : 60 + 70 + 35 + 25 = 190, les DCO redonnant la combinaison parentale pour ces deux gènes) :
\[ d_{vg-cn} = \frac{190}{1000} \times 100 = 19\,\% \]
Or $15 + 8 = 23$ cM. L'écart ($23 - 19 = 4$) correspond exactement aux DCO comptés deux fois ($2 \times 20/1000 \times 100 = 4$). La carte est cohérente : $vg$ --- 15 cM --- $b$ --- 8 cM --- $cn$.
\end{exoresolu}

\begin{popfigure}
\begin{center}
\begin{tikzpicture}[scale=1]
% chromosome
\draw[line width=3pt,popblue!60] (0,0) -- (11.5,0);
% loci
\foreach \x/\name/\allele in {1/vg/vg^{+}\ \text{ou}\ vg, 6/b/b^{+}\ \text{ou}\ b, 9.4/cn/cn^{+}\ \text{ou}\ cn}{
  \fill[poporange] (\x,0) circle (0.14);
  \node[above=6pt,popdark] at (\x,0) {\textbf{\name}};
}
% distances
\draw[<->,thick,poppurple] (1,-0.9) -- (6,-0.9);
\node[below,poppurple] at (3.5,-0.9) {15 cM};
\draw[<->,thick,poppurple] (6,-0.9) -- (9.4,-0.9);
\node[below,poppurple] at (7.7,-0.9) {8 cM};
\draw[<->,thick,popgreen] (1,-1.9) -- (9.4,-1.9);
\node[below,popgreen] at (5.2,-1.9) {$vg$--$cn$ : 19 cM (direct) $\approx$ 23 cM (somme, aux DCO pr\`es)};
\end{tikzpicture}
\end{center}
\legende{Carte factorielle du chromosome portant les gènes $vg$, $b$ et $cn$ chez la drosophile. Les loci (points orange) sont placés à des distances proportionnelles aux taux de recombinaison : 15 cM entre $vg$ et $b$, 8 cM entre $b$ et $cn$. La distance directe entre les gènes extrêmes (19 cM) est légèrement inférieure à la somme des intervalles (23 cM) car les doubles crossing-over, invisibles entre les extrêmes, ne sont pas détectés par le calcul direct.}
\end{popfigure}

\begin{savaistu}[La première carte génétique de l'histoire]
En 1913, Alfred \cle{Sturtevant}, alors étudiant de seulement 19 ans dans le laboratoire de Morgan à l'université Columbia, passa une nuit à analyser les données de croisements portant sur six gènes du chromosome X de la drosophile — délaissant, dit-on, ses devoirs de mathématiques. En une seule nuit de calculs, il établit la \emph{première carte génétique de l'histoire}, démontrant que les gènes sont disposés \emph{linéairement} le long des chromosomes. Cette idée révolutionnaire — mesurer des distances entre des entités invisibles à l'œil nu grâce à des statistiques de croisements — vaut à Morgan le prix Nobel en 1933 et fonde la génétique moderne, jusqu'aux cartes du génome humain utilisées aujourd'hui dans le diagnostic des maladies héréditaires.
\end{savaistu}

\begin{aretenir}[L'essentiel de la section]
\begin{itemize}
\item Le taux de recombinaison $p = \dfrac{\text{recombinés}}{\text{effectif total}} \times 100$ se calcule sur un croisement-test ; les classes parentales sont les plus nombreuses.
\item $1\,\%$ de recombinaison $= 1$ cM ; le taux ne dépasse jamais 50\,\%.
\item Pour trois gènes liés : les classes les plus rares sont les doubles crossing-over ; le gène qui « change » entre parentaux et DCO est le gène central.
\item Chaque intervalle se calcule en incluant les DCO ; la distance entre les extrêmes vaut la somme des intervalles aux DCO près.
\end{itemize}
\end{aretenir}

\begin{exercice}[À toi de jouer]
Un croisement-test chez le maïs donne 500 descendants : 210 $[AB]$, 200 $[ab]$, 50 $[Ab]$, 40 $[aB]$. 1) Identifier les classes parentales et recombinées. 2) Calculer le taux de recombinaison et la distance entre $A$ et $B$. 3) Un troisième gène $C$ donne 12\,\% de recombinaison avec $A$ et 30\,\% avec $B$. Proposer l'ordre des trois gènes et tracer la carte.
\end{exercice}

\begin{corrige}[Exercice : carte à trois gènes]
1) Classes parentales (les plus nombreuses) : $[AB]$ et $[ab]$ ; classes recombinées : $[Ab]$ et $[aB]$.
2) $p = \dfrac{50 + 40}{500} \times 100 = 18\,\%$, donc $A$ et $B$ sont distants de 18 cM.
3) $A$--$C$ = 12 cM et $B$--$C$ = 30 cM. Comme $12 + 18 = 30$, le gène $A$ est situé entre $C$ et $B$. Ordre : $C$ --- 12 cM --- $A$ --- 18 cM --- $B$, avec $C$--$B$ = 30 cM, ce qui vérifie l'additivité des distances.
\end{corrige}

\coursec{Exceptions à la monogénie : polygénie et pléiotropie}

Dans les sections précédentes, nous avons vu comment le brassage interchromosomique (gènes indépendants) et le brassage intrachromosomique (linkage et crossing-over) expliquent les proportions observées en F2 et permettent d'établir des cartes génétiques. Toutes ces analyses reposaient sur un postulat implicite : \textbf{chaque caractère étudié est gouverné par un seul gène, et chaque gène ne gouverne qu'un seul caractère}. Or l'observation attentive du vivant — la couleur de la peau humaine, la taille des individus, le rendement d'un champ de maïs — montre des caractères qui \emph{varient de façon continue}, sans classes nettes. Inversement, certaines maladies génétiques, comme la drépanocytose fréquente au Cameroun, affectent simultanément plusieurs organes. Ces faits obligent à dépasser le cadre mendélien strict : c'est l'objet de cette section.

\courssub{La monogénie : le cadre étudié jusqu'ici}

\begin{definition}[Monogénie]
Il y a \cle{monogénie} lorsqu'un caractère héréditaire est gouverné par \textbf{un seul gène} (existant en plusieurs allèles dans la population). Le génotype au niveau de ce locus détermine alors directement le phénotype pour ce caractère.
\end{definition}

Tous les croisements étudiés précédemment relèvent de la monogénie : la couleur des fleurs du pois, la forme des graines, la couleur du corps de la drosophile. En monohybridisme, un croisement entre deux races pures donne une F1 homogène et une F2 à \textbf{deux classes phénotypiques} dans les proportions $3/4 - 1/4$ (dominance complète) ou \textbf{trois classes} $1/4 - 2/4 - 1/4$ (dominance incomplète). La question qui guide cette section est la suivante : \emph{tous les caractères obéissent-ils à ce schéma ?}

\courssub{La polygénie : un caractère gouverné par plusieurs gènes}

\textbf{Situation-problème.} Un agriculteur de la plaine des Mbo croise une variété de blé à grains \textbf{rouge foncé} avec une variété à grains \textbf{blancs}. En F1, tous les grains sont \textbf{rouge clair} (intermédiaires). Jusqu'ici, rien d'anormal : on pourrait penser à une dominance incomplète. Mais en F2, au lieu des trois classes attendues ($1/4 - 2/4 - 1/4$), il observe \textbf{cinq classes} de couleur, du rouge foncé au blanc, avec les proportions suivantes :

\begin{center}
\begin{tabularx}{\linewidth}{|l|X|X|X|X|X|}
\hline
\rowcolor{popblueL} \textbf{Phénotype F2} & Rouge foncé & Rouge & Intermédiaire & Clair & Blanc \\
\hline
\textbf{Proportion} & $1/16$ & $4/16$ & $6/16$ & $4/16$ & $1/16$ \\
\hline
\end{tabularx}
\end{center}

\textbf{Hypothèse et interprétation.} Ces proportions sont incompatibles avec la monogénie. Elles s'expliquent si la couleur du grain est gouvernée par \textbf{deux gènes indépendants} $A/a$ et $B/b$ dont les allèles sont \emph{additifs} : chaque allèle « fort » ($A$ ou $B$) apporte une dose de pigment, chaque allèle « faible » ($a$ ou $b$) n'en apporte pas. La couleur dépend alors du \textbf{nombre total d'allèles forts} (de 0 à 4) :

\begin{center}
\begin{tabularx}{\linewidth}{|l|X|X|}
\hline
\rowcolor{popblueL} \textbf{Nombre d'allèles forts} & \textbf{Génotypes F2} & \textbf{Proportion} \\
\hline
4 & $AABB$ & $1/16$ \\
\hline
3 & $AABb$, $AaBB$ & $4/16$ \\
\hline
2 & $AAbb$, $aaBB$, $AaBb$ & $6/16$ \\
\hline
1 & $Aabb$, $aaBb$ & $4/16$ \\
\hline
0 & $aabb$ & $1/16$ \\
\hline
\end{tabularx}
\end{center}

On retrouve exactement les proportions observées : l'hypothèse est validée. La distribution $1-4-6-4-1$ n'est autre que la \textbf{loi binomiale} (coefficients du binôme de Newton), car il s'agit de compter les façons de répartir 4 allèles forts parmi 4 locus alléliques.

\begin{definition}[Polygénie]
Il y a \cle{polygénie} lorsqu'un \textbf{même caractère} est gouverné par \textbf{plusieurs gènes indépendants} (souvent situés sur des chromosomes différents ou très éloignés sur un même chromosome) dont les effets \textbf{s'additionnent} : on parle de \cle{gènes additifs} ou à effet cumulatif. Chaque allèle actif apporte une « dose » qui s'ajoute aux autres.
\end{definition}

\begin{propriete}[Nombre de classes phénotypiques en polygénie]
Pour un caractère gouverné par $n$ gènes additifs indépendants à deux allèles chacun, la F2 issue du croisement de deux races pures extrêmes comporte $2n+1$ \textbf{classes phénotypiques}, dont les effectifs suivent une \textbf{distribution binomiale} (coefficients de $(a+b)^{2n}$). Ainsi : $n=1 \Rightarrow 3$ classes ($1-2-1$) ; $n=2 \Rightarrow 5$ classes ($1-4-6-4-1$) ; $n=3 \Rightarrow 7$ classes ($1-6-15-20-15-6-1$).
\end{propriete}

\begin{popfigure}
\begin{center}
\begin{tikzpicture}
\begin{axis}[
    width=0.85\linewidth, height=6.5cm,
    ybar, bar width=22pt,
    xlabel={Nombre d'allèles forts (doses de pigment)},
    ylabel={Fréquence en F2},
    xtick={0,1,2,3,4},
    xmin=-0.4, xmax=4.4,
    ymin=0, ymax=0.45,
    ytick={0,0.0625,0.25,0.375},
    yticklabels={$0$,$1/16$,$4/16$,$6/16$},
    axis lines=left,
    nodes near coords,
    every node near coord/.append style={font=\small},
]
\addplot[fill=poporange!70, draw=popdark] coordinates {(0,0.0625) (1,0.25) (2,0.375) (3,0.25) (4,0.0625)};
\addplot[popcurve, popblue, domain=-0.4:4.4, samples=200] {0.4*exp(-((x-2)^2)/1.6)};
\end{axis}
\end{tikzpicture}
\end{center}
\legende{Histogramme des proportions phénotypiques en F2 pour un caractère polygénique à deux gènes additifs (couleur du grain de blé). Les effectifs $1/16 - 4/16 - 6/16 - 4/16 - 1/16$ dessinent une distribution en cloche (courbe bleue) : quand le nombre de gènes augmente et que l'environnement ajoute ses fluctuations, le caractère paraît varier de façon \emph{continue} (couleur de la peau, taille, rendement).}
\end{popfigure}

\textbf{Conséquence majeure : la variation continue.} Avec 2 gènes, on distingue 5 classes ; avec 3 gènes, 7 classes ; avec une dizaine de gènes, les classes se chevauchent et le caractère semble varier \emph{continûment} dans la population : c'est le cas de la \textbf{couleur de la peau humaine}, de la \textbf{taille}, du \textbf{rendement grain} chez le maïs ou le manioc. La courbe en cloche (distribution gaussienne) observée dans les populations naturelles est la signature statistique de la polygénie.

\begin{methode}[Reconnaître une polygénie]
\begin{enumerate}
\item Calculer les proportions phénotypiques de la F2 (ou de la descendance étudiée).
\item Si l'on obtient \textbf{plus de 3 classes} avec des proportions symétriques de type $1-4-6-4-1$ (ou $1-6-15-20-15-6-1$), écarter la monogénie.
\item Vérifier que les classes extrêmes ont la fréquence $(1/4)^n$, ce qui donne le nombre $n$ de gènes : $1/16 = (1/4)^2 \Rightarrow 2$ gènes.
\item Conclure : caractère polygénique à gènes additifs ; le phénotype dépend du \emph{nombre d'allèles forts}, non de leur identité.
\end{enumerate}
\end{methode}

\courssub{La pléiotropie : un gène qui influence plusieurs caractères}

\textbf{Observation clinique.} Dans les hôpitaux de Yaoundé et de Douala, les médecins suivent des patients atteints de \textbf{drépanocytose} (anémie falciforme). Chez un même malade, on observe simultanément : des globules rouges en forme de faucille, une anémie chronique, des crises douloureuses (occlusions vasculaires), une rate hypertrophiée... Or l'analyse génétique montre que tous ces signes découlent d'\textbf{une seule mutation} du gène codant la chaîne $\beta$ de l'hémoglobine (substitution d'un acide glutamique par une valine en position 6). Un seul gène, de nombreux caractères affectés.

\begin{definition}[Pléiotropie]
Il y a \cle{pléiotropie} lorsqu'\textbf{un seul gène} influence \textbf{plusieurs caractères} du phénotype. La protéine codée par ce gène intervient dans plusieurs tissus ou plusieurs voies métaboliques, si bien qu'une mutation a des conséquences multiples.
\end{definition}

\textbf{Autres exemples.} Chez la drosophile, certains gènes affectent à la fois la taille des ailes et la fertilité. Chez l'Homme, le syndrome de Marfan (gène de la fibrilline-1) affecte le squelette, les yeux et le système cardiovasculaire. Remarquons que la pléiotropie s'explique par la \emph{cascade} des effets : l'hémoglobine anormale (HbS) déforme le globule rouge (caractère cellulaire), ce qui provoque l'anémie hémolytique (caractère physiologique) et les occlusions vasculaires (caractère clinique).

\begin{attention}[Ne pas confondre polygénie et pléiotropie]
C'est l'erreur la plus fréquente au Baccalauréat. Retiens le \textbf{sens des flèches} :
\begin{itemize}
\item \textbf{Polygénie} : \emph{plusieurs gènes} $\rightarrow$ \emph{un seul caractère} (ex. couleur de la peau, taille, rendement).
\item \textbf{Pléiotropie} : \emph{un seul gène} $\rightarrow$ \emph{plusieurs caractères} (ex. drépanocytose : forme des hématies, anémie, résistance au paludisme).
\end{itemize}
Dans un exercice, demande-toi toujours : « combien de gènes ? combien de caractères ? » avant de conclure.
\end{attention}

\begin{popfigure}
\begin{center}
\begin{tikzpicture}[font=\small,
  gene/.style={rectangle, rounded corners, draw=popblue, fill=popblueL, minimum width=1.5cm, minimum height=0.7cm},
  car/.style={rectangle, rounded corners, draw=poporange, fill=poporangeL, minimum width=1.9cm, minimum height=0.7cm},
  titre/.style={font=\bfseries},
  fleche/.style={-{Stealth[length=2.5mm]}, thick, popdark}]
% Monogénie
\node[titre] at (0,3.4) {Monogénie};
\node[gene] (g1) at (0,2.4) {Gène 1};
\node[car] (c1) at (0,1.0) {Caractère A};
\draw[fleche] (g1) -- (c1);
% Polygénie
\node[titre] at (5,3.4) {Polygénie};
\node[gene] (g2) at (3.6,2.4) {Gène 1};
\node[gene] (g3) at (5,2.4) {Gène 2};
\node[gene] (g4) at (6.4,2.4) {Gène 3};
\node[car] (c2) at (5,1.0) {Caractère A};
\draw[fleche] (g2) -- (c2);
\draw[fleche] (g3) -- (c2);
\draw[fleche] (g4) -- (c2);
% Pléiotropie
\node[titre] at (10.4,3.4) {Pléiotropie};
\node[gene] (g5) at (10.4,2.4) {Gène 1};
\node[car] (c3) at (8.8,1.0) {Caractère A};
\node[car] (c4) at (10.4,0.2) {Caractère B};
\node[car] (c5) at (12,1.0) {Caractère C};
\draw[fleche] (g5) -- (c3);
\draw[fleche] (g5) -- (c4);
\draw[fleche] (g5) -- (c5);
\end{tikzpicture}
\end{center}
\legende{Schéma comparatif des relations gènes--caractères. \textbf{Monogénie} : un gène gouverne un caractère. \textbf{Polygénie} : plusieurs gènes additifs gouvernent un même caractère (ex. couleur du grain de blé). \textbf{Pléiotropie} : un gène unique influence plusieurs caractères (ex. gène de l'hémoglobine $\beta$ : forme des globules rouges, anémie, résistance au paludisme).}
\end{popfigure}

\begin{exoresolu}[Couleur des grains de blé : polygénie ou monogénie ?]
On croise une lignée pure de blé à grains rouge foncé avec une lignée pure à grains blancs. La F1 est uniformément intermédiaire. L'autofécondation des F1 donne en F2 : 12 plants à grains rouge foncé, 47 à grains rouges, 71 à grains intermédiaires, 50 à grains clairs et 11 à grains blancs (total : 191). Interpréter ces résultats.
\tcblower
\textbf{1. Analyse des effectifs.} Rapportés à 192 (multiple de 16 proche de 191), les effectifs théoriques d'une répartition $1-4-6-4-1$ seraient : $12 - 48 - 72 - 48 - 12$. Les valeurs observées ($12 - 47 - 71 - 50 - 11$) en sont très proches : la F2 comporte \textbf{5 classes} en proportions $1/16 - 4/16 - 6/16 - 4/16 - 1/16$.

\textbf{2. Rejet de la monogénie.} Un seul gène donnerait au plus 3 classes ($1/4 - 2/4 - 1/4$). L'hypothèse monogénique est rejetée.

\textbf{3. Hypothèse polygénique.} La classe extrême a pour fréquence $1/16 = (1/4)^2$ : le caractère est gouverné par \textbf{deux gènes additifs indépendants} $A/a$ et $B/b$. Les parents sont $AABB$ (rouge foncé) $\times$ $aabb$ (blanc) ; la F1 $AaBb$ (2 allèles forts) est intermédiaire, ce qui correspond à l'observation.

\textbf{4. Conclusion.} Il s'agit d'une \cle{polygénie} à deux gènes à effet cumulatif : le phénotype dépend du nombre d'allèles forts (0 à 4), d'où 5 classes phénotypiques en distribution binomiale.
\end{exoresolu}

\begin{savaistu}[Drépanocytose et paludisme : une pléiotropie qui protège]
Au Cameroun, environ 20 à 25 \% de la population porte l'allèle drépanocytaire $S$ (génotype $A/S$, dit « trait drépanocytaire »). Ce gène muté est \textbf{pléiotrope} : il modifie la forme des globules rouges, mais il confère aussi aux hétérozygotes $A/S$ une \textbf{résistance partielle au paludisme grave} — le parasite \emph{Plasmodium falciparum} se développe mal dans les hématies falciformes. C'est pourquoi l'allèle $S$, délétère à l'état homozygote ($S/S$ : drépanocytose majeure), a été maintenu par la sélection naturelle dans les zones d'endémie palustre comme le bassin du Congo et le Cameroun. Un bel exemple de la complexité des relations génotype--phénotype--environnement !
\end{savaistu}

\courssub{Bilan : les limites des lois de Mendel}

La monogénie, la polygénie et la pléiotropie montrent que les relations entre génotype et phénotype sont plus riches que le schéma « un gène -- un caractère » des premiers travaux de Mendel. Les lois de Mendel restent \textbf{exactes pour la transmission des allèles} (ségrégation, indépendance), mais le \emph{phénotype} résulte souvent de l'action conjointe de plusieurs gènes (polygénie), et un gène peut avoir des effets multiples (pléiotropie). À cela s'ajoute l'influence du milieu : la taille d'un individu dépend de ses gènes \emph{et} de sa nutrition.

\begin{aretenir}[L'essentiel]
\begin{itemize}
\item \textbf{Monogénie} : un caractère $\leftarrow$ un gène (F2 : 2 ou 3 classes).
\item \textbf{Polygénie} : un caractère $\leftarrow$ plusieurs gènes additifs ; F2 à $2n+1$ classes en distribution binomiale ($1-4-6-4-1$ pour $n=2$) ; explique la \textbf{variation continue} (peau, taille, rendement).
\item \textbf{Pléiotropie} : un gène $\rightarrow$ plusieurs caractères (drépanocytose : hématies falciformes + anémie + résistance au paludisme).
\item Méthode de diagnostic : compter les classes F2 et repérer les proportions binomiales symétriques.
\end{itemize}
\end{aretenir}

\coursec{Synthèse : méiose, fécondation et diversité génétique}

Nous venons de voir, avec la polygénie et la pléiotropie, que les relations entre gènes et caractères peuvent être complexes. Mais une question plus profonde demeure, celle qui a guidé toute cette leçon : pourquoi, au sein d'une même espèce, deux individus issus des mêmes parents sont-ils presque toujours différents ? Pourquoi des caractères nouveaux — et parfois des anomalies — réapparaissent-ils de génération en génération ? La réponse tient en trois mécanismes que nous allons maintenant rassembler : deux brassages qui ont lieu pendant la méiose, et le hasard de la fécondation. Cette synthèse est le point d'aboutissement du module : elle relie directement les lois statistiques des croisements (sections 1 à 5) aux événements cellulaires de la méiose.

\courssub{Les trois sources de brassage génétique}

\begin{definition}[Diversité génétique et unicité génétique]
La \cle{diversité génétique} d'une espèce est l'ensemble des combinaisons alléliques différentes présentes dans les génotypes des individus de cette espèce. L'\cle{unicité génétique} de l'individu désigne le fait que, sauf cas exceptionnel des jumeaux monozygotes (vrais jumeaux), chaque individu possède un génotype différent de celui de tout autre individu de l'espèce.
\end{definition}

\textbf{Première source : le brassage interchromosomique (anaphase I).} Au cours de l'anaphase I de la méiose, les paires de chromosomes homologues se séparent. Or l'orientation de chaque paire sur la plaque équatoriale est \emph{indépendante} de celle des autres paires : le chromosome d'origine paternelle d'une paire peut migrer vers un pôle ou vers l'autre, indépendamment du sort des chromosomes des autres paires. Il en résulte une répartition aléatoire des chromosomes parentaux dans les gamètes. C'est la base cellulaire de la \cle{troisième loi de Mendel} (ségrégation indépendante des caractères) observée dans le dihybridisme et le trihybridisme avec gènes indépendants : un hybride $A/a$ ; $B/b$ produit quatre types de gamètes $AB$, $Ab$, $aB$, $ab$ en proportions égales ($1/4$ chacun).

\textbf{Deuxième source : le brassage intrachromosomique (prophase I).} Lors de la prophase I, les chromosomes homologues s'apparient (synapsis) et des échanges de fragments peuvent se produire entre chromatides non s\oe urs : ce sont les \cle{crossing-over} (enjambements). Ces échanges créent des chromosomes \emph{recombinés}, portant des associations d'allèles différentes de celles des chromosomes parentaux. C'est la base cellulaire du \cle{linkage partiel} : un individu $AB/ab$ produit, outre les gamètes parentaux $AB$ et $ab$ (majoritaires), des gamètes recombinés $Ab$ et $aB$ (minoritaires), dont la proportion est le taux de recombinaison.

\textbf{Troisième source : la rencontre aléatoire des gamètes à la fécondation.} La fécondation est l'union de deux gamètes. Chaque parent produisant plusieurs types de gamètes, et la rencontre d'un spermatozoïde et d'un ovule étant un événement dû au hasard, le nombre de zygotes possibles est le produit des nombres de types de gamètes des deux parents. La fécondation ne crée pas de nouveaux allèles, mais elle réalise une \emph{association nouvelle} des gènes des deux parents.

\begin{popfigure}
\begin{center}
\begin{tikzpicture}[font=\small, >=Stealth, node distance=0.5cm]
% Cellule mère
\node[draw=popblue, fill=popblueL, rounded corners, minimum width=3.6cm, minimum height=1cm, align=center] (mere) {Cellule mère diploïde\\ $2n$ chromosomes};
% Méiose
\node[draw=popdark, fill=popcream, rounded corners, below=0.9cm of mere, minimum width=4.6cm, minimum height=2.6cm, align=center] (meiose) {};
\node[anchor=south, align=center, font=\bfseries] at (meiose.north) {MÉIOSE};
\node[below=0.15cm of meiose.north, align=left, text width=4.3cm] {\textbullet\ \textbf{Prophase I} : crossing-over\\ \hspace{0.3cm}$\Rightarrow$ brassage \emph{intra}chromosomique\\[2pt] \textbullet\ \textbf{Anaphase I} : migration au hasard\\ \hspace{0.3cm}des homologues $\Rightarrow$ brassage \emph{inter}chromosomique};
% Gamètes
\node[draw=popgreen, fill=popgreenL, rounded corners, below left=1.1cm and -0.3cm of meiose, minimum width=2.6cm, minimum height=0.9cm, align=center] (gam1) {Gamètes du parent 1\\ ($2^{n}$ types possibles)};
\node[draw=popgreen, fill=popgreenL, rounded corners, below right=1.1cm and -0.3cm of meiose, minimum width=2.6cm, minimum height=0.9cm, align=center] (gam2) {Gamètes du parent 2\\ ($2^{n}$ types possibles)};
% Fécondation
\node[draw=poporange, fill=poporangeL, rounded corners, below=2.6cm of meiose, minimum width=4.2cm, minimum height=0.9cm, align=center] (fecond) {FÉCONDATION : rencontre au hasard\\ de deux gamètes};
% Zygote
\node[draw=poppurple, fill=poppurpleL, rounded corners, below=0.8cm of fecond, minimum width=4.6cm, minimum height=0.9cm, align=center] (zygote) {Zygote diploïde ($2n$)\\ \textbf{combinaison unique} d'allèles};
% Flèches
\draw[->, thick, popblue] (mere) -- (meiose);
\draw[->, thick, popgreen] (meiose.south west) ++(0.4,0) -- (gam1.north);
\draw[->, thick, popgreen] (meiose.south east) ++(-0.4,0) -- (gam2.north);
\draw[->, thick, poporange] (gam1.south) |- ($(fecond.west)+(-0.2,0.15)$) -- (fecond.west);
\draw[->, thick, poporange] (gam2.south) |- ($(fecond.east)+(0.2,0.15)$) -- (fecond.east);
\draw[->, thick, poppurple] (fecond) -- (zygote);
\end{tikzpicture}
\end{center}
\legende{Schéma-bilan : les trois sources de diversité génétique. Deux brassages ont lieu pendant la méiose (crossing-over en prophase I ; migration aléatoire des chromosomes homologues en anaphase I) ; le troisième est la rencontre aléatoire des gamètes à la fécondation. Le zygote reçoit ainsi une combinaison d'allèles unique.}
\end{popfigure}

\courssub{Le calcul du nombre de combinaisons : une diversité astronomique}

Raisonnons d'abord simplement. Considérons un organisme dont les $n$ paires de chromosomes portent chacune deux allèles différents (individu hétérozygote pour $n$ gènes situés sur $n$ paires différentes). Pour chaque paire, un gamète reçoit l'un ou l'autre chromosome : 2 possibilités. Les paires étant indépendantes, le nombre de types de gamètes est :

\[ N_{\text{gamètes}} = 2^{n} \]

La fécondation réunissant un gamète de chaque parent, le nombre de zygotes génétiquement distincts qu'un couple peut produire est :

\[ N_{\text{zygotes}} = 2^{n} \times 2^{n} = \left(2^{n}\right)^{2} = 4^{n} \]

À cela s'ajoutent les crossing-over, qui multiplient encore le nombre de types de gamètes en créant des chromosomes recombinés pour les gènes liés.

\begin{exemplebox}[Le cas de l'espèce humaine]
Chez l'Homme, $n = 23$ paires de chromosomes. Par le seul brassage interchromosomique, chaque individu peut produire :
\[ 2^{23} = \num{8388608} \approx 8,4 \times 10^{6} \text{ types de gamètes différents.} \]
Un couple peut donc engendrer, sans même compter les crossing-over :
\[ \left(2^{23}\right)^{2} = 2^{46} \approx 7 \times 10^{13} \text{ combinaisons possibles,} \]
soit environ \textbf{70 000 milliards} de zygotes génétiquement distincts. En tenant compte des crossing-over (en moyenne 1 à 3 par paire de chromosomes à chaque méiose), le nombre réel est immensément plus grand. Aucun couple humain ne pourra jamais épuiser cette diversité.
\end{exemplebox}

\begin{savaistu}[La probabilité d'avoir un jumeau génétique]
La probabilité que deux enfants d'un même couple (naissances distinctes) reçoivent exactement la même combinaison d'allèles est de l'ordre de $1 / 7 \times 10^{13}$, avant même les crossing-over : elle est quasi nulle. Seuls les \cle{jumeaux monozygotes}, issus de la scission d'un même zygote, partagent le même génotype (à de rares mutations près). C'est pourquoi, au sein d'une même famille camerounaise comme partout ailleurs, des frères et s\oe urs peuvent présenter des groupes sanguins, des teints ou des sensibilités aux maladies différents.
\end{savaistu}

\courssub{Double rôle de la méiose et rôle de la fécondation}

\begin{propriete}[Rôles complémentaires de la méiose et de la fécondation]
La \cle{méiose} joue un double rôle :
\begin{itemize}
\item \textbf{conservation} : elle ramène le stock chromosomique de $2n$ (diploïdie) à $n$ (haploïdie) dans les gamètes, ce qui est indispensable pour que le nombre de chromosomes reste constant de génération en génération ;
\item \textbf{diversification} : par les brassages interchromosomique et intrachromosomique, elle produit des gamètes génétiquement variés.
\end{itemize}
La \cle{fécondation} joue aussi un double rôle :
\begin{itemize}
\item elle \textbf{restaure la diploïdie} ($n + n = 2n$) dans le zygote ;
\item elle réalise une \textbf{association nouvelle et unique} des gènes parentaux, par la rencontre au hasard de deux gamètes.
\end{itemize}
Conclusion de la leçon (savoir admis) : \emph{la diversité génétique des individus au sein d'une espèce résulte de la conjugaison des brassages méiotiques et du hasard de la fécondation.}
\end{propriete}

\begin{attention}[La diversité ne vient pas seulement des mutations]
Erreur fréquente au Bac : écrire que « les mutations sont la source de la diversité génétique des individus d'une fratrie ». Les mutations créent bien des \emph{allèles nouveaux}, mais elles sont rares (de l'ordre de $10^{-6}$ à $10^{-5}$ par gène et par génération). La source \textbf{majeure et immédiate} de la diversité entre individus, à génotypes parentaux constants, ce sont les \cle{brassages} (interchromosomique, intrachromosomique) et le hasard de la fécondation : ils ne créent aucun allèle nouveau, mais des \emph{combinaisons nouvelles} d'allèles existants. Rédige ainsi : « les caractères nouveaux observés chez les descendants résultent de la recombinaison d'allèles préexistants chez les parents ».
\end{attention}

\courssub{Tableau de synthèse des trois brassages}

\begin{popfigure}
\begin{center}
\begin{tikzpicture}[font=\scriptsize]
% Dimensions
\def\colA{0}
\def\colB{3.15}
\def\colC{7.35}
\def\colD{11.15}
\def\colE{15.2}
\def\rowTop{0.7}
\def\rowMid{0}
\def\rowOne{-1.6}
\def\rowTwo{-3.4}
\def\rowThree{-5.2}
% En-têtes
\fill[popblue] (\colA,\rowMid) rectangle (\colB,\rowTop); \node[white, align=center] at ($(\colA,\rowMid)!0.5!(\colB,\rowTop)$) {\textbf{Source de diversité}};
\fill[popblue] (\colB,\rowMid) rectangle (\colC,\rowTop); \node[white, align=center] at ($(\colB,\rowMid)!0.5!(\colC,\rowTop)$) {\textbf{Mécanisme (moment)}};
\fill[popblue] (\colC,\rowMid) rectangle (\colD,\rowTop); \node[white, align=center] at ($(\colC,\rowMid)!0.5!(\colD,\rowTop)$) {\textbf{Résultat}};
\fill[popblue] (\colD,\rowMid) rectangle (\colE,\rowTop); \node[white, align=center] at ($(\colD,\rowMid)!0.5!(\colE,\rowTop)$) {\textbf{Conséquence}};
% Ligne 1
\fill[popgreenL] (\colA,\rowOne) rectangle (\colB,\rowMid); \node[align=center, text width=2.9cm] at ($(\colA,\rowOne)!0.5!(\colB,\rowMid)$) {Brassage \textbf{inter}chromosomique};
\fill[popcream] (\colB,\rowOne) rectangle (\colC,\rowMid); \node[align=center, text width=3.9cm] at ($(\colB,\rowOne)!0.5!(\colC,\rowMid)$) {Migration au hasard des chromosomes homologues vers les pôles (\textbf{anaphase I})};
\fill[white] (\colC,\rowOne) rectangle (\colD,\rowMid); \node[align=center, text width=3.5cm] at ($(\colC,\rowOne)!0.5!(\colD,\rowMid)$) {$2^{n}$ types de gamètes ; gamètes équiprobables};
\fill[poporangeL] (\colD,\rowOne) rectangle (\colE,\rowMid); \node[align=center, text width=3.8cm] at ($(\colD,\rowOne)!0.5!(\colE,\rowMid)$) {Ségrégation indépendante des gènes situés sur des chromosomes différents ($3^{\text{ème}}$ loi de Mendel)};
% Ligne 2
\fill[popgreenL] (\colA,\rowTwo) rectangle (\colB,\rowOne); \node[align=center, text width=2.9cm] at ($(\colA,\rowTwo)!0.5!(\colB,\rowOne)$) {Brassage \textbf{intra}chromosomique};
\fill[popcream] (\colB,\rowTwo) rectangle (\colC,\rowOne); \node[align=center, text width=3.9cm] at ($(\colB,\rowTwo)!0.5!(\colC,\rowOne)$) {Crossing-over entre chromatides non s\oe urs de chromosomes homologues (\textbf{prophase I})};
\fill[white] (\colC,\rowTwo) rectangle (\colD,\rowOne); \node[align=center, text width=3.5cm] at ($(\colC,\rowTwo)!0.5!(\colD,\rowOne)$) {Gamètes recombinés minoritaires ; gamètes parentaux majoritaires};
\fill[poporangeL] (\colD,\rowTwo) rectangle (\colE,\rowOne); \node[align=center, text width=3.8cm] at ($(\colD,\rowTwo)!0.5!(\colE,\rowOne)$) {Linkage partiel ; taux de recombinaison $p$ \% ; permet la carte génétique};
% Ligne 3
\fill[popgreenL] (\colA,\rowThree) rectangle (\colB,\rowTwo); \node[align=center, text width=2.9cm] at ($(\colA,\rowThree)!0.5!(\colB,\rowTwo)$) {Fécondation};
\fill[popcream] (\colB,\rowThree) rectangle (\colC,\rowTwo); \node[align=center, text width=3.9cm] at ($(\colB,\rowThree)!0.5!(\colC,\rowTwo)$) {Rencontre \textbf{au hasard} d'un gamète mâle et d'un gamète femelle};
\fill[white] (\colC,\rowThree) rectangle (\colD,\rowTwo); \node[align=center, text width=3.5cm] at ($(\colC,\rowThree)!0.5!(\colD,\rowTwo)$) {$(2^{n})^{2}$ zygotes possibles par couple};
\fill[poporangeL] (\colD,\rowThree) rectangle (\colE,\rowTwo); \node[align=center, text width=3.8cm] at ($(\colD,\rowThree)!0.5!(\colE,\rowTwo)$) {Restauration de la diploïdie ; association unique des gènes parentaux};
% Contours et séparateurs
\draw[popdark, thick] (\colA,\rowTop) rectangle (\colE,\rowThree);
\draw[popdark] (\colA,\rowMid) -- (\colE,\rowMid);
\draw[popdark] (\colA,\rowOne) -- (\colE,\rowOne);
\draw[popdark] (\colA,\rowTwo) -- (\colE,\rowTwo);
\draw[popdark] (\colB,\rowTop) -- (\colB,\rowThree);
\draw[popdark] (\colC,\rowTop) -- (\colC,\rowThree);
\draw[popdark] (\colD,\rowTop) -- (\colD,\rowThree);
\end{tikzpicture}
\end{center}
\legende{Tableau de synthèse comparatif des trois sources de diversité génétique : mécanisme, résultat immédiat et conséquence génétique.}
\end{popfigure}

\courssub{Lien avec la famille de situations : récurrence des caractères nouveaux et des anomalies}

La famille de situations de cette leçon évoque la \emph{récurrence des caractères nouveaux et/ou des anomalies au sein des espèces}. Nous pouvons maintenant l'expliquer complètement :

\begin{itemize}
\item Les \textbf{caractères nouveaux} (nouvelles combinaisons phénotypiques chez les descendants, par exemple un poulet à crête frisée et plumage noir alors que les parents étaient crête simple/noir et frisée/blanc) résultent des \emph{recombinaisons} : brassage interchromosomique, crossing-over et hasard de la fécondation. Ces caractères « nouveaux » ne font que révéler des allèles déjà présents, masqués, chez les parents.
\item Certaines \textbf{anomalies} récurrentes s'expliquent de la même façon : un allèle récessif délétère, caché à l'état hétérozygote chez les parents (porteurs sains), peut se retrouver à l'état homozygote chez un descendant (probabilité $1/4$ si les deux parents sont hétérozygotes) — c'est le cas de la drépanocytose ou de l'albinisme, fréquemment observés dans les consultations des hôpitaux camerounais.
\end{itemize}

\begin{savaistu}[Hors programme strict : quand la méiose faillit — la trisomie 21]
Les mécanismes de la méiose, si perfectionnés soient-ils, connaissent des limites. Il arrive qu'une paire de chromosomes homologues ne se sépare pas en anaphase I (ou que les chromatides ne se séparent pas en anaphase II) : c'est une \cle{non-disjonction}. Il en résulte un gamète portant deux exemplaires d'un chromosome au lieu d'un. Après fécondation par un gamète normal, le zygote possède trois exemplaires de ce chromosome : c'est une \textbf{trisomie}. La trisomie 21 (environ 1 naissance sur 700) en est l'exemple le plus connu ; sa fréquence augmente avec l'âge de la mère. Ce complément, hors programme strict de Terminale D, montre que les anomalies chromosomiques sont la « limite » des mécanismes de la méiose — et constitue une excellente ouverture dans une copie de Bac.
\end{savaistu}

\courssub{Démarche type Bac : interpréter un croisement et le relier à la méiose}

\begin{methode}[Interpréter les résultats d'un croisement (démarche en 5 étapes)]
\begin{enumerate}
\item \textbf{Analyser les proportions} obtenues en F1 et F2 (ou dans un test-cross) : calculer les rapports ($3/4$--$1/4$ ; $9/16$--$3/16$--$3/16$--$1/16$ ; $1/4$--$1/4$--$1/4$--$1/4$ ; proportions inégales avec deux classes majoritaires\ldots).
\item \textbf{Dénombrer les gènes} en jeu : un seul caractère (monohybridisme) ou plusieurs (di- ou trihybridisme) ; vérifier la monogénie ou suspecter polygénie/pléiotropie si les proportions sont atypiques.
\item \textbf{Identifier le type de transmission} : gènes indépendants (proportions équilibrées) ou gènes liés (linkage total : 2 types de gamètes ; linkage partiel : 4 types de gamètes inégaux, parentaux majoritaires).
\item \textbf{Justifier par les mécanismes méiotiques} : ségrégation des allèles en anaphase I ; brassage interchromosomique pour les gènes indépendants ; crossing-over en prophase I pour le linkage partiel ; échiquier de croisement pour la fécondation aléatoire.
\item \textbf{Conclure et calculer} si demandé : taux de recombinaison $p = \dfrac{\text{effectif des recombinés}}{\text{effectif total}} \times 100$, puis carte génétique (distance en centimorgans entre les loci).
\end{enumerate}
\end{methode}

\begin{exoresolu}[Des proportions au mécanisme : exemple récapitulatif]
On croise deux drosophiles de race pure : femelle à corps gris et ailes longues $\times$ mâle à corps ébène et ailes vestigiales. La F1 est entièrement à corps gris et ailes longues. Le croisement-test d'une femelle F1 avec un mâle double récessif donne : 415 [gris, longues], 405 [ébène, vestigiales], 92 [gris, vestigiales], 88 [ébène, longues]. Interpréter ces résultats en les reliant aux mécanismes de la méiose.
\tcblower
\textbf{Étape 1 — Analyse des proportions.} Effectif total : $415 + 405 + 92 + 88 = 1000$. Deux classes parentales majoritaires ([gris, longues] et [ébène, vestigiales], soit $82\ \%$) et deux classes recombinées minoritaires ($18\ \%$).

\textbf{Étape 2 — Nombre de gènes.} Deux caractères (couleur du corps, forme des ailes) : dihybridisme. La F1 homogène montre que \emph{gris} ($G$) domine \emph{ébène} ($g$) et que \emph{longues} ($L$) domine \emph{vestigiales} ($l$).

\textbf{Étape 3 — Type de transmission.} Les quatre phénotypes ne sont pas équiprobables : les gènes sont \textbf{liés} (même chromosome), avec linkage \textbf{partiel}. La femelle F1 a le génotype $GL/gl$.

\textbf{Étape 4 — Justification méiotique.} En prophase I de la méiose, des crossing-over entre les deux loci produisent les gamètes recombinés $Gl$ et $gL$ ; les méioses sans crossing-over entre ces loci donnent les gamètes parentaux $GL$ et $gl$. Le mâle double récessif ne fournissant que des gamètes $gl$, les phénotypes de la descendance révèlent directement les gamètes de la femelle F1.

\textbf{Étape 5 — Calcul et carte.} Taux de recombinaison :
\[ p = \frac{92 + 88}{1000} \times 100 = 18\ \%. \]
Les deux gènes sont distants de \textbf{18 centimorgans} sur la carte génétique du chromosome. Conclusion : la diversité des descendants traduit le brassage intrachromosomique (crossing-over) combiné au hasard de la fécondation.
\end{exoresolu}

\begin{aretenir}[L'essentiel de la synthèse]
\begin{itemize}
\item Trois sources de diversité : brassage \textbf{inter}chromosomique (anaphase I, $2^{n}$ gamètes), brassage \textbf{intra}chromosomique (crossing-over en prophase I), rencontre \textbf{aléatoire} des gamètes à la fécondation.
\item Chez l'Homme : $(2^{23})^{2} \approx 7 \times 10^{13}$ combinaisons par couple, avant les crossing-over.
\item La méiose \textbf{conserve} le nombre de chromosomes (haploïdie des gamètes) et \textbf{diversifie} les combinaisons alléliques ; la fécondation \textbf{restaure} la diploïdie et crée une association unique des gènes parentaux.
\item Conséquence : hors jumeaux monozygotes, \textbf{chaque individu est génétiquement unique}.
\item Les caractères nouveaux et certaines anomalies récurrentes (maladies récessives) s'expliquent par la recombinaison d'allèles préexistants — et non principalement par les mutations.
\end{itemize}
\end{aretenir}

\begin{exercice}[Pour t'entraîner]
Chez une plante à fleurs, un individu double hétérozygote $AB/ab$ est soumis à un croisement-test. On obtient 640 descendants : 272 $[AB]$, 264 $[ab]$, 56 $[Ab]$, 48 $[aB]$.
\begin{enumerate}
\item Les deux gènes sont-ils indépendants ou liés ? Justifier.
\item Calculer le pourcentage de recombinaison et donner la distance génétique entre les deux loci.
\item Expliquer, en deux phrases, quels événements de la méiose ont produit les quatre types de gamètes.
\end{enumerate}
\end{exercice}

\begin{corrige}[Exercice]
\begin{enumerate}
\item Si les gènes étaient indépendants, le test-cross donnerait quatre classes équiprobables ($1/4$ chacune, soit 160 descendants par classe). Or deux classes parentales sont nettement majoritaires ($272 + 264 = 536$, soit $83,75\ \%$) et deux classes recombinées minoritaires : les gènes sont \textbf{liés}, avec linkage partiel.
\item Taux de recombinaison : $p = \dfrac{56 + 48}{640} \times 100 = \dfrac{104}{640} \times 100 = 16,25\ \%$. Les deux loci sont distants de \textbf{16,25 centimorgans} sur la carte génétique.
\item Les méioses sans crossing-over entre les deux loci ont produit les gamètes parentaux $AB$ et $ab$ (ségrégation des homologues en anaphase I) ; les méioses avec un crossing-over en prophase I entre les deux loci ont produit les gamètes recombinés $Ab$ et $aB$.
\end{enumerate}
\end{corrige}

\newpage

\coursec{Activité d'intégration}

\coursec{Leçon 04 : Sensibilisation sur les rôles de la méiose et de la fécondation dans le maintien de la diversité génétique des individus au sein d'une espèce}

\courssub{Activité d'intégration : Le mystère des drosophiles de Mvolyé}

\courssub{Objectifs de l'activité}

À l'issue de cette activité d'intégration, tu seras capable de :

\begin{itemize}
\item analyser les proportions issues de croisements-tests et déterminer si des gènes sont \cle{indépendants} ou \cle{liés} ;
\item identifier, dans une descendance, les \cle{classes parentales} et les \cle{classes recombinées} ;
\item calculer un \cle{pourcentage de recombinaison} et l'exprimer en centimorgans (cM) ;
\item construire une \cle{carte génétique} factorielle à partir des distances de recombinaison ;
\item expliquer, par les mécanismes de la \cle{méiose} (brassage interchromosomique et intrachromosomique) et de la \cle{fécondation}, l'apparition de phénotypes nouveaux ;
\item relier ces mécanismes à la diversité génétique observée dans les élevages et les cultures du Cameroun.
\end{itemize}

\courssub{Situation}

Au laboratoire de génétique d'une université de Yaoundé, une équipe d'enseignants-chercheurs élève des drosophiles (\emph{Drosophila melanogaster}) dans l'amphithéâtre de biologie du quartier Mvolyé. Ces petites mouches, qui se reproduisent en douze jours à \SI{25}{\celsius} et pondent des centaines d'{\oe}ufs, sont un matériel idéal pour étudier la transmission des caractères héréditaires.

L'équipe s'intéresse à trois caractères :

\begin{itemize}
\item la \textbf{couleur du corps} : corps \cle{gris} (allèle $G$, dominant) ou corps \cle{ébène} (allèle $g$, récessif) ;
\item la \textbf{forme des ailes} : ailes \cle{normales} (allèle $N$, dominant) ou ailes \cle{vestigiales} (allèle $n$, récessif) ;
\item la \textbf{couleur des yeux} : yeux \cle{rouges} (allèle $R$, dominant) ou yeux \cle{bruns} (allèle $r$, récessif).
\end{itemize}

Les chercheurs disposent de mouches de race pure et réalisent trois \cle{croisements-tests} successifs : à chaque fois, une femelle double hétérozygote (issue du croisement entre une race pure dominante et une race pure récessive pour les deux caractères considérés) est croisée avec un mâle double récessif. Les descendances, dénombrées sur plusieurs pontes, sont les suivantes.

\textbf{Croisement-test n°1} : femelle hétérozygote pour les gènes « couleur du corps » et « forme des ailes » $\times$ mâle à corps ébène et ailes vestigiales.

\begin{center}
\rowcolors{1}{popblueL}{white}
\begin{tabularx}{\linewidth}{|l|X|X|X|X|}
\hline
\rowcolor{popblue!40} Phénotype & [gris, normales] & [ébène, vestigiales] & [gris, vestigiales] & [ébène, normales] \\
\hline
Effectif & 415 & 415 & 85 & 85 \\
\hline
\end{tabularx}
\end{center}

\textbf{Croisement-test n°2} : femelle hétérozygote pour les gènes « couleur du corps » et « couleur des yeux » $\times$ mâle à corps ébène et yeux bruns.

\begin{center}
\rowcolors{1}{popgreenL}{white}
\begin{tabularx}{\linewidth}{|l|X|X|X|X|}
\hline
\rowcolor{popgreen!40} Phénotype & [gris, rouges] & [ébène, bruns] & [gris, bruns] & [ébène, rouges] \\
\hline
Effectif & 248 & 252 & 255 & 245 \\
\hline
\end{tabularx}
\end{center}

\textbf{Croisement-test n°3} : femelle hétérozygote pour les gènes « forme des ailes » et « couleur des yeux » $\times$ mâle à ailes vestigiales et yeux bruns.

\begin{center}
\rowcolors{1}{poporangeL}{white}
\begin{tabularx}{\linewidth}{|l|X|X|X|X|}
\hline
\rowcolor{poporange!40} Phénotype & [normales, rouges] & [vestigiales, bruns] & [normales, bruns] & [vestigiales, rouges] \\
\hline
Effectif & 260 & 240 & 255 & 245 \\
\hline
\end{tabularx}
\end{center}

Le chef de laboratoire confie ces résultats à ta classe : « Retrouvez l'organisation de ces trois gènes sur les chromosomes de la drosophile, et expliquez-moi d'où viennent ces mouches aux phénotypes nouveaux que je n'avais pas dans mes souches de départ ! »

\courssub{Consignes}

Par groupes de quatre, et en rédigeant un rapport qui sera affiché puis discuté en débat corrigé :

\begin{enumerate}
\item Pour chacun des trois croisements-tests, calcule les proportions (en \%) de chaque classe phénotypique et indique si les deux gènes étudiés sont \cle{indépendants} ou \cle{liés}. Justifie ta réponse en comparant les proportions obtenues aux proportions théoriques attendues.
\item Dans le (ou les) croisement(s) révélant un linkage, identifie les \cle{classes parentales} et les \cle{classes recombinées}. Écris le génotype de la femelle hétérozygote en précisant la disposition des allèles (en \emph{cis} ou en \emph{trans}).
\item Calcule le \cle{pourcentage de recombinaison} entre les gènes liés et exprime la distance génétique correspondante en centimorgans.
\item Construis la \cle{carte génétique} des trois gènes étudiés (position relative des gènes sur le ou les chromosomes).
\item Schématise, par un schéma légendé, le mécanisme méiotique (crossing-over en prophase I) qui a produit les classes recombinées du croisement-test n°1.
\item Rédige une conclusion d'une dizaine de lignes expliquant comment la méiose (brassages inter- et intrachromosomique) et la fécondation engendrent la diversité génétique, et illustre ce propos par un exemple camerounais (élevage de volailles, culture du cacao ou du maïs).
\end{enumerate}

\courssub{Ressources à mobiliser}

\begin{aretenir}[Boîte à outils de l'activité]
\begin{itemize}
\item \textbf{Croisement-test} : croisement d'un individu de phénotype dominant (génotype inconnu) avec un individu homozygote récessif. Les phénotypes de la descendance reflètent directement les \cle{gamètes} produits par l'individu testé.
\item \textbf{Gènes indépendants} : un double hétérozygote produit 4 types de gamètes équiprobables ($25\,\%$ chacun) ; le croisement-test donne 4 classes à $\approx 25\,\%$ (brassage interchromosomique).
\item \textbf{Gènes liés} : les classes \emph{parentales} sont majoritaires, les classes \emph{recombinées} minoritaires (crossing-over en prophase I de méiose).
\item \textbf{Pourcentage de recombinaison} :
\[ p = \frac{\text{nombre d'individus recombinés}}{\text{effectif total}} \times 100 \]
\item \textbf{Distance génétique} : $1\,\%$ de recombinaison $= 1$ centimorgan (cM). Deux gènes sont dits indépendants lorsque $p \approx 50\,\%$.
\end{itemize}
\end{aretenir}

\courssub{Démarche et corrigé commenté}

\begin{exoresolu}[Le mystère des drosophiles de Mvolyé : résolution complète]
Énoncé : à partir des trois croisements-tests ci-dessus, déterminer l'organisation des gènes $G/g$, $N/n$ et $R/r$, calculer les distances génétiques, construire la carte factorielle et expliquer l'origine des phénotypes nouveaux.

\tcblower

\textbf{Étape 1 : analyse des proportions de chaque croisement.}

\emph{Croisement-test n°1} (corps $\times$ ailes). Effectif total : $415 + 415 + 85 + 85 = \num{1000}$ mouches. Les proportions sont :
\[ [\text{gris, normales}] = 41,5\,\% \quad ; \quad [\text{ébène, vestigiales}] = 41,5\,\% \]
\[ [\text{gris, vestigiales}] = 8,5\,\% \quad ; \quad [\text{ébène, normales}] = 8,5\,\% \]
On obtient 4 classes \textbf{inégales}, avec deux classes majoritaires et deux classes minoritaires : ce n'est ni le $25/25/25/25$ de l'indépendance, ni le $50/50$ du linkage absolu. Les gènes $G/g$ et $N/n$ sont donc \cle{liés} (même chromosome), avec un \cle{linkage partiel} : il y a eu crossing-over.

\emph{Croisement-test n°2} (corps $\times$ yeux). Total : $\num{1000}$ mouches. Proportions : $24,8\,\%$ ; $25,2\,\%$ ; $25,5\,\%$ ; $24,5\,\%$, soit $\approx 25\,\%$ par classe. Les gènes $G/g$ et $R/r$ sont \cle{indépendants}.

\emph{Croisement-test n°3} (ailes $\times$ yeux). Total : $\num{1000}$ mouches. Proportions : $26\,\%$ ; $24\,\%$ ; $25,5\,\%$ ; $24,5\,\%$, soit $\approx 25\,\%$ par classe. Les gènes $N/n$ et $R/r$ sont \cle{indépendants}.

\textbf{Étape 2 : classes parentales et recombinées (croisement n°1).}

Les classes les plus fréquentes correspondent aux gamètes non modifiés de la femelle : $GN$ et $gn$. Les classes parentales sont donc [gris, normales] et [ébène, vestigiales] ; les classes recombinées sont [gris, vestigiales] (gamète $Gn$) et [ébène, normales] (gamète $gN$). La femelle hétérozygote a ses allèles en position \emph{cis} :
\[ \dfrac{G\ \ N}{g\ \ n} \]

\textbf{Étape 3 : pourcentage de recombinaison.}

\[ p = \frac{85 + 85}{1000} \times 100 = 17\,\% \]
La distance entre les gènes $G/g$ et $N/n$ est donc de \cle{17 cM} (17 unités de recombinaison).

\textbf{Étape 4 : carte génétique.}

$G$ et $N$ sont liés à 17 cM l'un de l'autre ; $R$ est indépendant des deux : il est porté par une \textbf{autre paire de chromosomes} (ou, à titre de discussion, très éloigné sur le même chromosome, ce qui donnerait $p \approx 50\,\%$ — au niveau Terminale, on retient : chromosome différent).

\begin{popfigure}
\begin{center}
\begin{tikzpicture}[scale=1]
% Chromosome porteur de G et N
\draw[line width=6pt, popblue!60, rounded corners=3pt] (0,0) -- (6,0);
\fill[popdark] (1,0) circle (2.5pt);
\fill[popdark] (4.4,0) circle (2.5pt);
\node[above, popdark] at (1,0.15) {\small $G/g$};
\node[above, popdark] at (4.4,0.15) {\small $N/n$};
\draw[<->, poporange, thick] (1,-0.45) -- (4.4,-0.45) node[midway, below, poporange]{\small 17 cM};
\node[left, popblue] at (-0.2,0) {\small Chromosome I};
% Autre chromosome
\draw[line width=6pt, popgreen!60, rounded corners=3pt] (0,-1.8) -- (6,-1.8);
\fill[popdark] (3,-1.8) circle (2.5pt);
\node[above, popdark] at (3,-1.65) {\small $R/r$};
\node[left, popgreen!80!black] at (-0.2,-1.8) {\small Chromosome II};
\end{tikzpicture}
\end{center}
\legende{Carte génétique des trois gènes étudiés : $G/g$ et $N/n$ sont liés sur le chromosome I, distants de 17 cM ; $R/r$ est porté par une autre paire de chromosomes (indépendance vis-à-vis des deux autres gènes).}
\end{popfigure}

\textbf{Étape 5 : mécanisme méiotique.}

En prophase I de la méiose, les chromosomes homologues s'apparient (formation des bivalents). Entre les locus de $G/g$ et de $N/n$, un \cle{crossing-over} (enjambement) échange des segments de chromatides non-s{\oe}urs dans environ $34\,\%$ des méiocytes (chaque crossing-over ne recombinaisonne que 2 chromatides sur 4, d'où $17\,\%$ de gamètes recombinés). Il en résulte 4 types de gamètes : $GN$ et $gn$ (parentaux, majoritaires), $Gn$ et $gN$ (recombinés, minoritaires). C'est le \cle{brassage intrachromosomique}. Parallèlement, la répartition au hasard des chromosomes I et II en anaphase I réalise le \cle{brassage interchromosomique} pour le gène $R/r$.

\textbf{Étape 6 : conclusion.}

Les mouches « nouvelles » du laboratoire de Mvolyé (corps gris à ailes vestigiales, par exemple) ne résultent d'aucune mutation : elles proviennent du \textbf{brassage génétique}. La méiose, par le crossing-over (brassage intrachromosomique) et par la ségrégation indépendante des paires de chromosomes (brassage interchromosomique), produit des gamètes génétiquement variés ; la fécondation, en réunissant au hasard deux de ces gamètes, multiplie encore les combinaisons. Ainsi, chaque individu (hors jumeaux vrais) est génétiquement unique, tout en restant diploïde à chaque génération : méiose et fécondation maintiennent à la fois la \textbf{stabilité} du caryotype et la \textbf{diversité} des génotypes au sein de l'espèce. C'est ce même principe qu'exploitent les éleveurs de volailles de l'Ouest camerounais ou les sélectionneurs de cacaoyers hybrides à la station IRAD de Nkoemvone : en croisant des variétés aux caractères complémentaires (vigueur, résistance aux maladies, rendement), ils créent par recombinaison des descendances nouvelles, mieux adaptées, sans modifier le nombre de chromosomes de l'espèce.
\end{exoresolu}

\begin{attention}[Pièges fréquents rencontrés pendant l'activité]
\begin{itemize}
\item Ne conclus à l'indépendance que si les 4 classes sont \emph{sensiblement égales} ($\approx 25\,\%$) : de petites fluctuations (24,5 \% contre 25,5 \%) sont normales, ce sont des écarts d'échantillonnage.
\item Les classes parentales sont toujours les \textbf{plus nombreuses} ; ne les confonds pas avec les classes de phénotype dominant.
\item Le pourcentage de recombinaison se calcule \emph{uniquement} avec les individus recombinés, pas avec une seule classe recombinée : ici $(85 + 85)/1000$ et non $85/1000$.
\item $1\,\%$ de recombinaison $= 1$ cM : n'écris pas « 17 \% de distance », mais « distance de 17 cM ».
\end{itemize}
\end{attention}

\courssub{Bilan de l'activité}

Cette activité t'a permis de mobiliser, dans une situation authentique de laboratoire, l'ensemble des savoir-faire du brassage génétique : l'\textbf{analyse de proportions} issues de croisements-tests pour distinguer linkage et indépendance, l'\textbf{identification des classes parentales et recombinées}, le \textbf{calcul d'un taux de recombinaison} et la \textbf{construction d'une carte génétique}. Au-delà des calculs, elle a mis en évidence l'idée centrale du module : la diversité génétique des individus au sein d'une espèce ne naît pas du hasard aveugle, mais de mécanismes précis et reproductibles — le brassage interchromosomique, le brassage intrachromosomique par crossing-over, puis la rencontre aléatoire des gamètes à la fécondation. Ces mécanismes, que tu sais désormais quantifier, expliquent aussi bien l'apparition de drosophiles aux phénotypes inédits dans un laboratoire de Yaoundé que la variabilité des maïs, des cacaoyers ou des volailles qui font la richesse agricole du Cameroun.

\newpage

\coursec{Méthodes et exercices résolus}

\courssub{Méthode 1 : Interpréter un croisement de monohybridisme}

\begin{methode}[Analyser un croisement portant sur un seul caractère]
\begin{enumerate}
\item \textbf{Identifier le caractère et ses allèles.} Repérer les deux phénotypes parentaux (ex. : graine lisse / graine ridée) et vérifier qu'il s'agit bien de \cle{lignées pures} (homozygotes).
\item \textbf{Déterminer la dominance.} Si la F$_1$ est homogène et présente un seul des deux phénotypes parentaux, l'allèle exprimé en F$_1$ est \cle{dominant} ; l'autre est \cle{récessif}. Choisir une notation claire : lettre majuscule pour le dominant ($L$), minuscule pour le récessif ($l$).
\item \textbf{Écrire les génotypes.} Parents : $L/L$ et $l/l$. Hybride F$_1$ : $L/l$. Toujours écrire les génotypes avant l'échiquier.
\item \textbf{Construire l'échiquier de croisement} de F$_1 \times$ F$_1$ : gamètes $L$ et $l$ de chaque parent, proportions $\frac{1}{2}$ et $\frac{1}{2}$.
\item \textbf{Conclure sur F$_2$.} Génotypes : $\frac{1}{4}\,L/L$ ; $\frac{1}{2}\,L/l$ ; $\frac{1}{4}\,l/l$. Phénotypes : $\frac{3}{4}$ dominant, $\frac{1}{4}$ récessif (proportions $3:1$), conformément à la \cle{loi de disjonction des allèles} (2\textsuperscript{e} loi de Mendel, dite aussi « loi de pureté des gamètes ») : les deux allèles d'un couple se séparent lors de la formation des gamètes.
\item \textbf{Valider par un test-cross} si demandé : croisement de l'hybride avec le parent récessif $l/l$ donne $\frac{1}{2}$ [$L$] et $\frac{1}{2}$ [$l$].
\end{enumerate}
\textbf{Réflexes :} toute proportion théorique doit être comparée aux effectifs observés ; un écart faible s'explique par les fluctuations d'échantillonnage.
\end{methode}

\begin{exoresolu}[Monohybridisme chez le haricot]
On croise deux lignées pures de haricots (\emph{Phaseolus vulgaris}), l'une à graines lisses, l'autre à graines ridées. La F$_1$ est entièrement constituée de plantes à graines lisses. Le croisement F$_1 \times$ F$_1$ donne en F$_2$ : 5474 graines lisses et 1850 graines ridées.
\begin{enumerate}
\item Déterminer la relation de dominance entre les allèles.
\item Établir les génotypes des parents et de la F$_1$.
\item Donner les proportions théoriques attendues en F$_2$ et les comparer aux résultats observés.
\end{enumerate}
\tcblower
\textbf{1. Dominance.} Les parents sont de lignées pures et la F$_1$ est homogène, de phénotype [lisse]. L'allèle « lisse » est donc \cle{dominant} ; on le note $L$. L'allèle « ridé » est récessif, noté $l$.

\textbf{2. Génotypes.} Les parents, homozygotes, s'écrivent $L/L$ (lisse) et $l/l$ (ridé). Chacun ne produit qu'un type de gamète ($L$ ou $l$) par disjonction des allèles à la méiose. Tous les individus F$_1$ reçoivent $L$ d'un parent et $l$ de l'autre : génotype $L/l$, phénotype [lisse] puisque $L$ domine $l$.

\textbf{3. Proportions en F$_2$.} L'hybride F$_1$ $L/l$ produit deux types de gamètes en proportions égales : $\frac{1}{2}\,L$ et $\frac{1}{2}\,l$. L'échiquier de croisement F$_1 \times$ F$_1$ donne :
\begin{center}
\begin{tabular}{|c|c|c|}
\hline
\rowcolor{popblueL} Gamètes & $\frac{1}{2}\,L$ & $\frac{1}{2}\,l$ \\
\hline
$\frac{1}{2}\,L$ & $\frac{1}{4}\,L/L$ [lisse] & $\frac{1}{4}\,L/l$ [lisse] \\
\hline
$\frac{1}{2}\,l$ & $\frac{1}{4}\,L/l$ [lisse] & $\frac{1}{4}\,l/l$ [ridé] \\
\hline
\end{tabular}
\end{center}
Proportions théoriques : $\frac{3}{4}$ [lisse] et $\frac{1}{4}$ [ridé].

Comparaison avec les observations : effectif total $= 5474 + 1850 = 7324$ graines.
\begin{align*}
\text{lisses attendues} &= \tfrac{3}{4} \times 7324 = 5493 \quad (\text{observées : } 5474)\\
\text{ridées attendues} &= \tfrac{1}{4} \times 7324 = 1831 \quad (\text{observées : } 1850)
\end{align*}
Les écarts relatifs ($\approx 0{,}3\,\%$ et $\approx 1\,\%$) sont faibles : ils relèvent des fluctuations d'échantillonnage.

\textbf{Conclusion :} les résultats vérifient la 2\textsuperscript{e} loi de Mendel : les deux allèles d'un même gène se séparent (disjonctent) lors de la formation des gamètes, et le hasard de la fécondation restitue en F$_2$ les proportions $3/4 - 1/4$.
\end{exoresolu}

\courssub{Méthode 2 : Interpréter un dihybridisme avec gènes indépendants}

\begin{methode}[Reconnaître et exploiter le brassage interchromosomique]
\begin{enumerate}
\item \textbf{Vérifier qu'on étudie deux caractères} (deux gènes), chacun avec deux allèles, et des parents de race pure différant par les deux caractères.
\item \textbf{Observer la F$_1$} : elle doit être homogène (double hétérozygote $A/a\,;\,B/b$), ce qui fixe les dominances.
\item \textbf{Analyser la F$_2$ ou le test-cross.} Quatre phénotypes en F$_2$ dans les proportions $9:3:3:1$ (soit $\frac{9}{16}, \frac{3}{16}, \frac{3}{16}, \frac{1}{16}$) ou quatre phénotypes équiprobables ($\frac{1}{4}$ chacun) au test-cross $\Rightarrow$ les deux gènes sont \cle{indépendants} (portés par deux paires de chromosomes différentes).
\item \textbf{Justifier par la méiose} : l'hybride produit 4 types de gamètes équiprobables ($AB$, $Ab$, $aB$, $ab$) grâce à la \cle{ségrégation indépendante des paires de chromosomes homologues} en anaphase I : c'est le \cle{brassage interchromosomique}.
\item \textbf{Repérer les phénotypes recombinés} (associations de caractères absentes chez les parents) : ils prouvent le brassage.
\end{enumerate}
\textbf{Formule à retenir :} pour $n$ gènes indépendants, l'hybride F$_1$ produit $2^n$ types de gamètes équiprobables ; la F$_2$ compte $2^n$ classes phénotypiques (dominance totale).
\end{methode}

\begin{exoresolu}[Dihybridisme chez la souris]
On croise une souris de race pure à poil noir et lisse avec une souris de race pure à poil blanc et frisé. La F$_1$ est entièrement noire et lisse. Le croisement d'une F$_1$ avec une souris blanche et frisée (test-cross) donne : 52 [noires, lisses] ; 48 [noires, frisées] ; 51 [blanches, lisses] ; 49 [blanches, frisées].
\begin{enumerate}
\item Préciser les dominances et écrire le génotype de la F$_1$.
\item Interpréter les résultats du test-cross : les deux gènes sont-ils liés ou indépendants ?
\end{enumerate}
\tcblower
\textbf{1. Dominances et génotype F$_1$.} La F$_1$, issue de parents de race pure, est homogène [noire, lisse] : l'allèle « noir » ($N$) domine « blanc » ($n$) ; l'allèle « lisse » ($L$) domine « frisé » ($l$). Génotypes parentaux : $N/N\,;\,L/L$ et $n/n\,;\,l/l$. La F$_1$ est $N/n\,;\,L/l$.

\textbf{2. Interprétation du test-cross.} Le parent testeur $n/n\,;\,l/l$ ne fournit que des gamètes $n\,l$ : les phénotypes des descendants révèlent directement les gamètes produits par la F$_1$. Effectif total : $52+48+51+49 = 200$.
\begin{center}
\begin{tabular}{|l|c|c|}
\hline
\rowcolor{popblueL} Phénotype (gamète F$_1$) & Effectif & Fréquence \\
\hline
[noir, lisse] ($N\,L$) & 52 & $0{,}26$ \\
\hline
[noir, frisé] ($N\,l$) & 48 & $0{,}24$ \\
\hline
[blanc, lisse] ($n\,L$) & 51 & $0{,}255$ \\
\hline
[blanc, frisé] ($n\,l$) & 49 & $0{,}245$ \\
\hline
\end{tabular}
\end{center}
Les quatre classes sont \textbf{équiprobables} ($\approx \frac{1}{4}$ chacune). Or, si les gènes étaient liés sur le même chromosome, on obtiendrait deux classes parentales nettement majoritaires. L'équiprobabilité des quatre gamètes montre que les deux gènes sont portés par \textbf{deux paires de chromosomes différentes} : ils sont \cle{indépendants}.

\textbf{Conclusion :} en anaphase I de méiose, les deux paires de chromosomes homologues s'orientent indépendamment l'une de l'autre (brassage interchromosomique), d'où 4 gamètes équiprobables et un test-cross à $4 \times \frac{1}{4}$. Les phénotypes [noir, frisé] et [blanc, lisse] sont des phénotypes \cle{recombinés}, témoins de la diversité génétique créée par la méiose.
\end{exoresolu}

\courssub{Méthode 3 : Interpréter un trihybridisme}

\begin{methode}[Généraliser à trois caractères]
\begin{enumerate}
\item \textbf{Identifier les trois gènes} et fixer les dominances grâce à l'homogénéité de la F$_1$ (triple hétérozygote $A/a\,;\,B/b\,;\,C/c$).
\item \textbf{Compter les types de gamètes} : si les trois gènes sont indépendants, la F$_1$ produit $2^3 = 8$ gamètes équiprobables ($\frac{1}{8}$ chacun).
\item \textbf{Prévoir les proportions} : test-cross $\rightarrow$ 8 classes à $\frac{1}{8}$ ; F$_2$ $\rightarrow$ 8 classes phénotypiques en proportions $27:9:9:9:3:3:3:1$ (sur 64).
\item \textbf{Tester l'indépendance deux à deux} si les proportions ne sont pas celles attendues : regrouper les données par paires de gènes pour déceler un éventuel linkage entre deux d'entre eux.
\item \textbf{Conclure} en précisant quels gènes sont indépendants (brassage interchromosomique) et lesquels sont liés (voir méthode 4).
\end{enumerate}
\textbf{Réflexe :} face à un gros tableau de résultats, toujours commencer par additionner les classes deux à deux pour retrouver les ségrégations $3:1$ (en F$_2$) ou $1:1$ (en test-cross) de chaque gène pris isolément — cela vérifie la validité des dominances.
\end{methode}

\begin{exoresolu}[Trihybridisme chez la drosophile]
Chez la drosophile, on croise une femelle de race pure à corps gris, ailes longues et yeux rouges avec un mâle à corps ébène, ailes vestigiales et yeux bruns. La F$_1$ est entièrement [grise, ailes longues, yeux rouges]. Le test-cross des femelles F$_1$ donne (sur 1600 descendants) environ 200 individus dans chacune des 8 classes phénotypiques possibles.
\begin{enumerate}
\item Écrire le génotype de la F$_1$.
\item Montrer que les trois gènes sont indépendants et justifier les proportions observées.
\end{enumerate}
\tcblower
\textbf{1. Génotype F$_1$.} L'homogénéité de la F$_1$ impose : gris ($G$) dominant sur ébène ($g$) ; ailes longues ($L$) dominantes sur vestigiales ($l$) ; yeux rouges ($R$) dominants sur bruns ($r$). La F$_1$ est triple hétérozygote : $G/g\,;\,L/l\,;\,R/r$.

\textbf{2. Indépendance des gènes.} Le parent testeur $g/g\,;\,l/l\,;\,r/r$ ne produit que des gamètes $g\,l\,r$ ; chaque classe phénotypique du test-cross reflète donc un type de gamète de la F$_1$. On observe 8 classes d'effectifs égaux : $\frac{1600}{8} = 200$ individus, soit une fréquence de $\frac{1}{8} = 0{,}125$ par classe.

La F$_1$ produit donc 8 gamètes équiprobables :
\[ G\,L\,R \;;\; G\,L\,r \;;\; G\,l\,R \;;\; G\,l\,r \;;\; g\,L\,R \;;\; g\,L\,r \;;\; g\,l\,R \;;\; g\,l\,r \]
Or $8 = 2^3$ : c'est le nombre maximal de gamètes pour trois gènes, obtenu uniquement si chaque paire de chromosomes s'oriente indépendamment des deux autres en anaphase I. Vérification gène par gène : pour le caractère couleur du corps, chaque phénotype apparaît dans 4 des 8 classes, soit $200 \times 4 = 800$ [gris] et $800$ [ébène], c'est-à-dire $\frac{1}{2}-\frac{1}{2}$, conforme à la ségrégation d'un couple d'allèles en test-cross ; idem pour les deux autres caractères.

\textbf{Conclusion :} les trois gènes sont portés par trois paires de chromosomes différentes ; le brassage interchromosomique à la méiose génère à lui seul $2^3 = 8$ types de gamètes, donc $8$ combinaisons phénotypiques équiprobables au test-cross.
\end{exoresolu}

\courssub{Méthode 4 : Mettre en évidence un linkage et un crossing-over}

\begin{methode}[Distinguer linkage total et linkage partiel]
\begin{enumerate}
\item \textbf{Réaliser (ou analyser) un test-cross} de l'hybride F$_1$ avec le double récessif : c'est le seul croisement qui révèle directement les gamètes de la F$_1$.
\item \textbf{Comparer les effectifs des 4 classes :}
\begin{itemize}
\item 2 classes seulement (parentales), en proportions égales $\Rightarrow$ \cle{linkage total} (gènes sur le même chromosome, jamais séparés ; l'hybride ne produit que 2 types de gamètes) ;
\item 4 classes \textbf{inégales}, avec 2 classes parentales \textbf{majoritaires} et 2 classes recombinées \textbf{minoritaires} $\Rightarrow$ \cle{linkage partiel} avec \cle{crossing-over} (brassage intrachromosomique en prophase I) ;
\item 4 classes égales $\Rightarrow$ gènes indépendants (pas de linkage).
\end{itemize}
\item \textbf{Identifier les gamètes parentaux} (les plus fréquents) : ils donnent la disposition des allèles sur les chromosomes de la F$_1$ (en \emph{cis} $\frac{AB}{ab}$ ou en \emph{trans} $\frac{Ab}{aB}$).
\item \textbf{Expliquer les recombinés} : échange de segments chromatidiens entre chromosomes homologues (chiasma) au cours de la prophase I.
\end{enumerate}
\textbf{Piège classique :} ne jamais conclure à un linkage à partir d'une F$_2$ seule ; le test-cross est l'outil décisif.
\end{methode}

\begin{exoresolu}[Linkage partiel chez la drosophile]
Une drosophile femelle F$_1$ à corps gris et ailes longues (génotype $G/g\,;\,L/l$, parents de race pure [gris, longues] $\times$ [ébène, vestigiales]) est croisée avec un mâle à corps ébène et ailes vestigiales. Résultats : 415 [gris, longues] ; 405 [ébène, vestigiales] ; 92 [gris, vestigiales] ; 88 [ébène, longues].
\begin{enumerate}
\item Montrer que les deux gènes sont liés.
\item Expliquer l'apparition des classes minoritaires.
\end{enumerate}
\tcblower
\textbf{1. Mise en évidence du linkage.} Effectif total : $415+405+92+88 = 1000$. Fréquences :
\begin{center}
\begin{tabular}{|l|c|c|c|}
\hline
\rowcolor{popblueL} Phénotype & Effectif & Fréquence & Type \\
\hline
[gris, longues] & 415 & $0{,}415$ & parental \\
\hline
[ébène, vestigiales] & 405 & $0{,}405$ & parental \\
\hline
[gris, vestigiales] & 92 & $0{,}092$ & recombiné \\
\hline
[ébène, longues] & 88 & $0{,}088$ & recombiné \\
\hline
\end{tabular}
\end{center}
Les quatre classes ne sont \textbf{pas équiprobables} : si les gènes étaient indépendants, on attendrait $\frac{1}{4}$ ($250$ individus) par classe. Les deux classes parentales représentent $\frac{415+405}{1000} = 0{,}82$, soit $82\,\%$ des descendants : les allèles restent associés comme chez les parents. Les gènes $G/g$ et $L/l$ sont donc \textbf{liés} sur la même paire de chromosomes ; la F$_1$ a ses allèles en position \emph{cis} : $\dfrac{G\,L}{g\,l}$.

\textbf{2. Origine des classes minoritaires.} Les phénotypes recombinés [gris, vestigiales] et [ébène, longues] ($18\,\%$ au total) proviennent de gamètes $G\,l$ et $g\,L$. Ceux-ci naissent d'un \cle{crossing-over} : en prophase I de méiose, les chromosomes homologues s'apparient (synapsis) et deux chromatides non sœurs échangent des segments au niveau d'un chiasma situé entre les deux loci. Il en résulte des chromatides recombinées portant $G\,l$ et $g\,L$.

\textbf{Conclusion :} le linkage n'est pas absolu ; le crossing-over réalise un \cle{brassage intrachromosomique} qui complète le brassage interchromosomique et accroît encore la diversité des gamètes, donc celle des individus.
\end{exoresolu}

\begin{savaistu}[Morgan, la drosophile et la naissance des cartes génétiques]
C'est Thomas Hunt Morgan et son équipe (Université Columbia, vers 1910--1915) qui découvrirent le linkage et le crossing-over en travaillant sur la mouche \emph{Drosophila melanogaster}, un insecte facile à élever et au cycle court (environ 10 jours). En 1911, Alfred Sturtevant, alors étudiant, eut l'idée d'utiliser les taux de recombinaison comme mesures de distances pour tracer la première \cle{carte génétique} d'un chromosome. Morgan reçut le prix Nobel de physiologie ou médecine en 1933. La drosophile reste aujourd'hui un organisme modèle majeur de la génétique, y compris dans les laboratoires de recherche africains qui étudient la résistance des moustiques vecteurs du paludisme aux insecticides.
\end{savaistu}

\courssub{Méthode 5 : Calculer un pourcentage de recombinaison}

\begin{methode}[Déterminer le taux de recombinaison entre deux gènes liés]
\begin{enumerate}
\item \textbf{Utiliser obligatoirement un test-cross} de l'hybride (les phénotypes des descendants = gamètes de l'hybride).
\item \textbf{Repérer les classes recombinées} : ce sont les deux classes \textbf{minoritaires} (ou celles dont le phénotype ne correspond à aucun parent de la F$_1$).
\item \textbf{Appliquer la formule :}
\[ \text{Taux de recombinaison} = \frac{\text{nombre d'individus recombinés}}{\text{effectif total}} \times 100 \quad (\text{en } \%) \]
\item \textbf{Interpréter} : ce pourcentage estime la \cle{distance génétique} entre les deux loci, exprimée en \cle{centimorgans} (cM) : $1\,\%$ de recombinaison $= 1$ cM.
\item \textbf{Vérifier la cohérence} : le taux est toujours compris entre $0\,\%$ (linkage total) et $50\,\%$ (indépendance apparente).
\end{enumerate}
\textbf{Réflexe :} additionner les DEUX classes recombinées avant de diviser ; oublier l'une d'elles divise le résultat par deux.
\end{methode}

\begin{exoresolu}[Calcul d'un taux de recombinaison chez le maïs]
Chez le maïs, le gène $C/c$ contrôle la couleur du grain (coloré/incolore) et le gène $S/s$ sa forme (plein/rétracté). Une plante F$_1$ issue du croisement [coloré, plein] $\times$ [incolore, rétracté] est soumise à un test-cross. On obtient : 4032 [colorés, pleins] ; 149 [colorés, rétractés] ; 152 [incolores, pleins] ; 4035 [incolores, rétractés]. Déterminer le pourcentage de recombinaison entre les deux gènes et la distance qui les sépare.
\tcblower
\textbf{Étape 1 — Identification des classes.} Les phénotypes parentaux de la F$_1$ sont [coloré, plein] et [incolore, rétracté] : ce sont les classes majoritaires ($4032$ et $4035$), ce qui confirme le linkage. Les classes recombinées sont [coloré, rétracté] ($149$) et [incolore, plein] ($152$).

\textbf{Étape 2 — Calcul.}
\begin{align*}
\text{Recombinés} &= 149 + 152 = 301\\
\text{Effectif total} &= 4032 + 149 + 152 + 4035 = 8368\\
\text{Taux de recombinaison} &= \frac{301}{8368} \times 100 \approx 3{,}6\,\%
\end{align*}

\textbf{Étape 3 — Interprétation.} Un taux de $3{,}6\,\%$ signifie qu'environ $3{,}6$ méioses sur $100$ chez la F$_1$ ont été le siège d'un crossing-over entre les deux loci. La distance génétique entre les gènes $C/c$ et $S/s$ est donc estimée à \textbf{3,6 cM} (ou 3,6 unités de recombinaison).

\textbf{Conclusion :} les deux gènes sont très proches sur le chromosome (linkage fort) ; seuls $3{,}6\,\%$ des gamètes de la F$_1$ sont recombinés.
\end{exoresolu}

\courssub{Méthode 6 : Établir une carte génétique}

\begin{methode}[Construire la carte factorielle d'un chromosome]
\begin{enumerate}
\item \textbf{Calculer le taux de recombinaison pour chaque paire de gènes} à partir des test-cross (méthode 5). Pour un trihybridisme avec gènes liés, identifier d'abord les classes parentales (les plus fréquentes) et les doubles recombinants (les plus rares).
\item \textbf{Déterminer l'ordre des gènes} : le gène situé \emph{au milieu} est celui qui change de position dans les doubles recombinants par rapport aux parentaux. Autre critère : les deux gènes les plus éloignés sont ceux dont le taux de recombinaison est le plus élevé.
\item \textbf{Placer les gènes sur une droite graduée en cM} : la distance entre les gènes extrêmes est (approximativement) la somme des distances intermédiaires.
\item \textbf{Schématiser la carte} : segment de droite, loci marqués par des traits, distances annotées en cM.
\item \textbf{Vérifier} : taux extrême $\approx$ somme des taux intermédiaires (l'écart vient des doubles crossing-over, qui « masquent » des recombinations).
\end{enumerate}
\end{methode}

\begin{exoresolu}[Carte génétique de trois gènes chez la drosophile]
Trois gènes liés de la drosophile donnent, par test-cross de la triple hétérozygote, les taux de recombinaison suivants : entre $A$ et $B$ : $12\,\%$ ; entre $B$ et $C$ : $7\,\%$ ; entre $A$ et $C$ : $18\,\%$. Établir la carte génétique du chromosome.
\tcblower
\textbf{Étape 1 — Gènes extrêmes.} Le taux le plus élevé ($18\,\%$) sépare $A$ et $C$ : ce sont les deux gènes \textbf{extrêmes}. Le gène $B$ occupe donc la position \textbf{centrale}.

\textbf{Étape 2 — Distances.} $A$--$B = 12$ cM et $B$--$C = 7$ cM. On vérifie : $12 + 7 = 19$ cM, légèrement supérieur aux $18$ cM observés entre $A$ et $C$. L'écart de $1$ point s'explique par les \cle{doubles crossing-over} entre $A$ et $C$ : deux échanges successifs rétablissent la configuration parentale pour les marqueurs extrêmes, qui échappent ainsi au comptage.

\textbf{Étape 3 — Carte.}
\begin{popfigure}
\begin{center}
\begin{tikzpicture}[scale=1]
\draw[very thick, popblue] (0,0) -- (9.5,0);
\foreach \x/\name in {0/A, 6/B, 9.5/C}{
  \draw[very thick, popdark] (\x,-0.15) -- (\x,0.15);
  \node[above, popdark] at (\x,0.15) {\textbf{\name}};
}
\draw[<->, poporange, thick] (0,-0.6) -- (6,-0.6) node[midway, below]{12 cM};
\draw[<->, poporange, thick] (6,-0.6) -- (9.5,-0.6) node[midway, below]{7 cM};
\draw[<->, popgreen, thick] (0,-1.3) -- (9.5,-1.3) node[midway, below]{$A$--$C$ : 18 cM observés (19 cM attendus)};
\end{tikzpicture}
\end{center}
\legende{Carte génétique du chromosome : ordre des loci $A$ -- $B$ -- $C$ et distances en centimorgans (cM), déduites des taux de recombinaison.}
\end{popfigure}

\textbf{Conclusion :} l'ordre des gènes sur le chromosome est $A$ -- $B$ -- $C$, avec $d(A,B) = 12$ cM et $d(B,C) = 7$ cM ; la carte factorielle traduit les positions relatives des loci, établies à partir des fréquences de crossing-over.
\end{exoresolu}

\courssub{Méthode 7 : Reconnaître polygénie et pléiotropie}

\begin{methode}[Identifier les exceptions à la monogénie]
\begin{enumerate}
\item \textbf{Compter les classes phénotypiques en F$_2$} et comparer aux proportions mendéliennes classiques ($3:1$, $9:3:3:1$).
\item \textbf{Polygénie} : un \emph{seul} caractère est gouverné par \emph{plusieurs gènes} à effets additifs. Indices : variation \textbf{continue} du caractère (taille, teint de la peau, rendement), distribution en cloche en F$_2$, proportions du type $1:4:6:4:1$ pour deux gènes additifs.
\item \textbf{Pléiotropie} : un \emph{seul} gène influence \emph{plusieurs caractères}. Indice : des anomalies ou caractères apparemment sans rapport sont \textbf{toujours associés} et se transmettent ensemble (ex. : drépanocytose : anémie + déformation des globules rouges en faucille + résistance partielle au paludisme chez l'hétérozygote).
\item \textbf{Conclure} en rappelant que la relation « un gène — un caractère » (monogénie) est un cas particulier : le génotype agit sur le phénotype par des voies souvent complexes.
\end{enumerate}
\end{methode}

\begin{exoresolu}[Polygénie de la couleur du grain de blé]
On croise deux lignées pures de blé, l'une à grains rouge foncé, l'autre à grains blancs. La F$_1$ est uniformément rouge intermédiaire. En F$_2$, on observe une gamme continue de teintes, du rouge foncé au blanc, réparties approximativement ainsi : 1/16 rouge foncé ; 4/16 rouge ; 6/16 rouge clair ; 4/16 rose pâle ; 1/16 blanc. Interpréter ces résultats.
\tcblower
\textbf{1. Écart avec la monogénie.} Un seul gène à deux allèles donnerait au plus 3 classes en F$_2$ ($1:2:1$ en dominance incomplète). Or on compte \textbf{5 classes} en proportions $1:4:6:4:1$ : il faut faire intervenir \textbf{deux gènes} $A/a$ et $B/b$ indépendants, à effets \textbf{additifs} sur le même caractère (couleur du grain).

\textbf{2. Modèle.} Chaque allèle dominant ($A$ ou $B$) apporte une dose de pigment ; les allèles récessifs n'en apportent pas. La F$_1$ $A/a\,;\,B/b$ possède 2 doses : teinte intermédiaire. En F$_2$, le nombre de doses varie de 0 à 4 :
\begin{center}
\begin{tabular}{|l|c|c|}
\hline
\rowcolor{popblueL} Doses de pigment & Génotypes & Fréquence \\
\hline
4 & $A/A\,;\,B/B$ & $1/16$ \\
\hline
3 & $A/A\,;\,B/b$ ; $A/a\,;\,B/B$ & $4/16$ \\
\hline
2 & $A/A\,;\,b/b$ ; $a/a\,;\,B/B$ ; $A/a\,;\,B/b$ & $6/16$ \\
\hline
1 & $A/a\,;\,b/b$ ; $a/a\,;\,B/b$ & $4/16$ \\
\hline
0 & $a/a\,;\,b/b$ & $1/16$ \\
\hline
\end{tabular}
\end{center}
On retrouve exactement les proportions observées $1:4:6:4:1$ (coefficients du binôme $(a+b)^4$).

\textbf{Conclusion :} la couleur du grain de blé est un caractère \cle{polygénique} : plusieurs gènes à effets cumulatifs déterminent un même caractère, ce qui engendre une variation continue — situation fréquente pour les caractères quantitatifs (taille, rendement agricole, teint de la peau humaine).
\end{exoresolu}

\begin{aretenir}[L'essentiel sur le brassage génétique]
\begin{itemize}
\item \textbf{Brassage interchromosomique} : ségrégation indépendante des paires de chromosomes homologues en anaphase I $\Rightarrow$ $2^n$ gamètes équiprobables pour $n$ gènes indépendants (test-cross : classes équiprobables ; F$_2$ : $9:3:3:1$ pour deux gènes).
\item \textbf{Brassage intrachromosomique} : crossing-over entre chromatides non sœurs en prophase I $\Rightarrow$ gamètes recombinés minoritaires pour des gènes liés (test-cross : 2 classes parentales majoritaires).
\item \textbf{Taux de recombinaison} $= \dfrac{\text{recombinés}}{\text{total}} \times 100$ ; il mesure la distance entre loci en centimorgans ($1\,\% = 1$ cM) et permet d'établir la \cle{carte génétique}.
\item \textbf{Exceptions à la monogénie} : polygénie (plusieurs gènes $\rightarrow$ un caractère, variation continue) et pléiotropie (un gène $\rightarrow$ plusieurs caractères).
\item \textbf{Enjeu} : brassages méiotiques + hasard de la fécondation = unicité génétique de chaque individu et diversité au sein de l'espèce, matière première de l'évolution.
\end{itemize}
\end{aretenir}

\begin{attention}[Les 5 erreurs qui coûtent des points au Bac]
\begin{enumerate}
\item \textbf{Confondre « gènes liés » et « gènes indépendants » à partir d'une F$_2$ seule.} La preuve décisive du linkage repose sur le \cle{test-cross} : 4 classes équiprobables = indépendance ; 2 classes parentales majoritaires = linkage. Cite toujours le croisement qui justifie ta conclusion.
\item \textbf{Oublier une classe recombinée dans le calcul du taux de recombinaison.} Le numérateur est la \emph{somme des deux classes minoritaires} : $\text{taux} = \dfrac{\text{recombinés}}{\text{total}} \times 100$. Diviser par l'effectif d'une seule classe ou oublier le $\times 100$ fausse tout le résultat.
\item \textbf{Écrire les génotypes sans préciser la disposition des allèles en cas de linkage.} Pour des gènes liés, écris $\dfrac{A\,B}{a\,b}$ (cis) ou $\dfrac{A\,b}{a\,B}$ (trans), et non $A/a\,;\,B/b$ qui sous-entend l'indépendance.
\item \textbf{Affirmer qu'un crossing-over a lieu « à chaque méiose » ou « entre chromatides sœurs ».} Le crossing-over est un événement \emph{aléatoire} de prophase I, entre \cle{chromatides non sœurs} de chromosomes \emph{homologues} ; sa fréquence dépend de la distance entre les loci.
\item \textbf{Mal rédiger la conclusion génétique.} Une conclusion attendue comporte trois éléments : la \emph{relation de dominance} (avec justification par la F$_1$), la \emph{localisation des gènes} (liés ou indépendants, avec les proportions qui le prouvent) et le \emph{brassage mis en jeu} (interchromosomique par ségrégation des chromosomes en anaphase I, intrachromosomique par crossing-over en prophase I). Une conclusion sans justification chiffrée perd la moitié des points.
\end{enumerate}
\end{attention}

\newpage

\coursec{Exercices}

\courssub{Je vérifie mes connaissances}

\begin{exercice}[QCM : brassages génétiques]
Pour chaque question, une seule réponse est exacte. Recopie la lettre correspondante.
\begin{enumerate}
\item Le brassage interchromosomique a lieu :
\begin{enumerate}
\item entre chromatides d'un même chromosome ;
\item lors de la séparation aléatoire des paires de chromosomes homologues en anaphase I de méiose ;
\item lors de la fécondation uniquement ;
\item entre gènes situés sur le même chromosome.
\end{enumerate}
\item Un croisement-test entre un individu double hétérozygote et un double récessif donne 4 phénotypes en proportions inégales (42 \%, 42 \%, 8 \%, 8 \%). On en déduit que :
\begin{enumerate}
\item les gènes sont indépendants ;
\item les gènes sont totalement liés ;
\item les gènes sont partiellement liés avec recombinaison de 16 \% ;
\item il s'agit d'une polygénie.
\end{enumerate}
\item Le crossing-over se produit :
\begin{enumerate}
\item en prophase I de méiose, entre chromatides non sœurs ;
\item en anaphase II de méiose ;
\item pendant la mitose ;
\item au moment de la fécondation.
\end{enumerate}
\item La pléiotropie désigne :
\begin{enumerate}
\item le contrôle d'un caractère par plusieurs gènes ;
\item l'action d'un même gène sur plusieurs caractères ;
\item la liaison de deux gènes sur un chromosome ;
\item l'apparition d'allèles multiples.
\end{enumerate}
\end{enumerate}
\end{exercice}

\begin{exercice}[Vrai ou faux justifié]
Indique si chaque affirmation est vraie ou fausse, puis justifie en une ou deux phrases.
\begin{enumerate}
\item Deux gènes indépendants donnent toujours, en F2 d'un dihybridisme, quatre phénotypes en proportions 9/16, 3/16, 3/16, 1/16.
\item Le linkage absolu supprime toute diversité gamétique chez un double hétérozygote : il ne produit que deux types de gamètes.
\item Plus deux gènes sont éloignés sur un chromosome, plus leur pourcentage de recombinaison est faible.
\item La fécondation contribue à la diversité génétique par la rencontre aléatoire de deux gamètes génétiquement différents.
\item Dans une polygénie à deux gènes additifs, la F2 d'un croisement entre lignées pures extrêmes comporte cinq classes phénotypiques.
\end{enumerate}
\end{exercice}

\begin{exercice}[Définitions et vocabulaire]
Définis précisément, en deux lignes maximum, chacun des termes suivants et donne pour chacun un exemple :
\begin{enumerate}
\item monogénie et polygénie ;
\item linkage partiel ;
\item pourcentage (taux) de recombinaison ;
\item carte génétique d'un chromosome ;
\item gamète parental et gamète recombiné.
\end{enumerate}
\end{exercice}

\begin{exercice}[Calculs directs]
\begin{enumerate}
\item Chez la drosophile, le croisement-test d'une femelle double hétérozygote donne 500 descendants : 210 de phénotype parental 1, 205 de phénotype parental 2, 43 et 42 de phénotypes recombinés. Calcule le pourcentage de recombinaison entre les deux gènes.
\item Trois gènes A, B et C sont liés. Le taux de recombinaison entre A et B vaut 12 \% ; celui entre B et C vaut 8 \%. Sachant que B est le gène central, calcule la distance A--C en centimorgans (cM) et représente la carte factorielle.
\item Un double hétérozygote $AB//ab$ produit 1\,000 gamètes avec un taux de recombinaison de 20 \%. Donne le nombre de gamètes de chaque type ($AB$, $ab$, $Ab$, $aB$).
\end{enumerate}
\end{exercice}

\courssub{Je m'entraîne}

\begin{exercice}[Croisement-test chez la drosophile]
Chez la drosophile, le corps gris ($G$) domine le corps ébène ($e$) et les ailes normales ($N$) dominent les ailes vestigiales ($n$). Une femelle double hétérozygote issue du croisement de lignées pures [gris, normales] $\times$ [ébène, vestigiales] est croisée avec un mâle double récessif. On obtient 800 descendants :
\begin{center}
\begin{tabular}{|l|c|c|c|c|}
\hline
\rowcolor{popblueL} Phénotype & [gris, normales] & [ébène, vestigiales] & [gris, vestigiales] & [ébène, normales] \\
\hline
Effectif & 332 & 328 & 74 & 66 \\
\hline
\end{tabular}
\end{center}
\begin{enumerate}
\item Écris le génotype de la femelle double hétérozygote en précisant la disposition (cis ou trans) des allèles.
\item Montre, en comparant les proportions observées à celles d'un croisement-test avec gènes indépendants, que les deux gènes sont liés.
\item Identifie les gamètes parentaux et les gamètes recombinés produits par la femelle.
\item Calcule le pourcentage de recombinaison entre les deux gènes et donne la distance génétique en cM.
\end{enumerate}
\end{exercice}

\begin{exercice}[Dihybridisme chez la souris]
On croise deux lignées pures de souris : souris à pelage noir et lisse $\times$ souris à pelage blanc et frisé. La F1 est entièrement [noir, lisse]. Le croisement des souris F1 entre elles donne en F2, sur 320 souriceaux : 180 [noir, lisse], 62 [noir, frisé], 59 [blanc, lisse], 19 [blanc, frisé].
\begin{enumerate}
\item Détermine les relations de dominance entre les allèles de chaque gène.
\item Vérifie que les proportions F2 correspondent au rapport 9:3:3:1 et conclus sur la localisation des deux gènes.
\item Établis l'échiquier de croisement de la F2 (gamètes et génotypes).
\item On réalise le croisement-test d'une souris F1. Prévois les phénotypes attendus et leurs proportions, puis les effectifs théoriques sur 200 descendants.
\end{enumerate}
\end{exercice}

\begin{exercice}[Linkage absolu chez le haricot]
Chez le haricot, on croise une lignée pure à fleurs violettes et graines lisses avec une lignée pure à fleurs blanches et graines ridées. La F1 est homogène. Le croisement-test de la F1 ne donne que deux types de descendants : 248 [violettes, lisses] et 252 [blanches, ridées].
\begin{enumerate}
\item Détermine les allèles dominants et écris les génotypes des parents et de la F1.
\item Explique pourquoi on obtient seulement deux phénotypes au croisement-test. Quel nom donne-t-on à cette situation ?
\item Quels gamètes la F1 produit-elle et dans quelles proportions ?
\item Quel phénomène méiotique, absent ici, aurait permis l'apparition de phénotypes recombinés ? Précise le stade de la méiose concerné.
\end{enumerate}
\end{exercice}

\begin{exercice}[Polygénie : couleur des grains de blé]
On croise une variété de blé à grains rouge foncé avec une variété à grains blancs. La F1 est uniformément rouge clair. En F2, on obtient sur 100 plants : 6 à grains rouge foncé, 25 à grains rouge, 38 à grains rouge clair, 25 à grains rose pâle, 6 à grains blancs (proportions voisines de 1:4:6:4:1).
\begin{enumerate}
\item Explique pourquoi ces résultats ne peuvent pas s'interpréter par un monohybridisme classique.
\item Propose un modèle à deux gènes additifs $A/a$ et $B/b$ (chaque allèle dominant ajoute une dose de pigment) et écris les génotypes des parents, de la F1 et des cinq classes de F2.
\item Vérifie que l'échiquier de croisement F1 $\times$ F1 rend compte des proportions 1:4:6:4:1.
\item Comment appelle-t-on ce type d'hérédité ? Cite un autre exemple chez l'Homme.
\end{enumerate}
\end{exercice}

\courssub{Je me prépare au Bac}

\begin{exercice}[Type Bac — Trihybridisme avec gènes liés chez la drosophile]
Chez \emph{Drosophila melanogaster}, on étudie trois gènes : corps gris ($g^+$) / corps jaune ($g$), œil rouge ($r^+$) / œil écarlate ($r$), ailes normales ($a^+$) / ailes coupées ($a$). Une femelle triple hétérozygote issue de lignées pures est croisée avec un mâle triple récessif (croisement-test). On obtient 1\,000 descendants :
\begin{center}
\begin{tabular}{|l|c|}
\hline
\rowcolor{popblueL} Phénotype du descendant & Effectif \\
\hline
[gris, rouge, normales] & 358 \\
\hline
[jaune, écarlate, coupées] & 362 \\
\hline
[gris, écarlate, coupées] & 92 \\
\hline
[jaune, rouge, normales] & 88 \\
\hline
[gris, rouge, coupées] & 41 \\
\hline
[jaune, écarlate, normales] & 39 \\
\hline
[gris, écarlate, normales] & 11 \\
\hline
[jaune, rouge, coupées] & 9 \\
\hline
\end{tabular}
\end{center}
\begin{enumerate}
\item Justifie que les trois gènes sont liés sur le même chromosome.
\item Identifie les classes parentales et les classes issues de doubles crossing-over. Justifie ta réponse.
\item Déduis de la comparaison de ces deux classes l'ordre des trois gènes sur le chromosome.
\item Calcule le pourcentage de recombinaison dans chacun des deux intervalles, en tenant compte des doubles crossing-over.
\item Établis la carte génétique complète du chromosome (positions et distances en cM).
\item Écris le génotype de la femelle triple hétérozygote en respectant l'ordre des gènes trouvé.
\end{enumerate}
\end{exercice}

\begin{exercice}[Type Bac — La drépanocytose au Cameroun : monohybridisme et pléiotropie]
La drépanocytose (anémie falciforme) est fréquente dans les régions camerounaises d'endémie palustre. Elle est due à un allèle muté $S$ du gène de l'hémoglobine : les individus $SS$ sont atteints d'une anémie grave, les $AA$ sont sains, et les hétérozygotes $AS$ (porteurs du « trait drépanocytaire ») sont cliniquement quasi normaux mais présentent, en conditions d'hypoxie, une falciformation partielle des globules rouges, une légère anémie et une résistance accrue au paludisme grave.

\textbf{Document 1.} Deux parents $AS$ consultent au centre de santé avant un mariage. Ils ont déjà un enfant.

\textbf{Document 2.} Dans une zone de forte transmission palustre, les fréquences génotypiques observées chez des nouveau-nés sont : $AA$ : \num{0,64} ; $AS$ : \num{0,32} ; $SS$ : \num{0,04}, alors qu'en zone sans paludisme la fréquence de $AS$ est beaucoup plus faible.
\begin{enumerate}
\item À partir du document 1, établis l'échiquier de croisement et détermine les probabilités pour chaque génotype et phénotype de la descendance.
\item Représente l'arbre généalogique d'une famille où deux parents $AS$ ont quatre enfants dont un $SS$ (convention : cercle = femme, carré = homme, symbole rempli = atteint).
\item Montre, en t'appuyant sur les multiples effets de l'allèle $S$ (forme des globules rouges, anémie, résistance au paludisme), que ce gène illustre la pléiotropie.
\item Explique, à l'aide du document 2, pourquoi l'allèle $S$ se maintient à fréquence élevée dans les zones d'endémie palustre (avantage sélectif des hétérozygotes).
\item Propose un message de sensibilisation génétique à destination des couples $AS$ au Cameroun.
\end{enumerate}
\end{exercice}

\begin{exercice}[Type Bac — Synthèse : brassage génétique et diversité chez la volaille]
Un éleveur de la localité de Bafoussam croise des poules pondeuses. Il étudie deux caractères : la couleur du plumage (noir $N$ dominant / blanc $b$) et la forme de la crête (crête en rose $R$ dominante / crête simple $r$).

\textbf{Premier croisement.} Une lignée pure [noir, crête en rose] $\times$ [blanc, crête simple] donne une F1 homogène. Les hybrides F1 croisés entre eux donnent en F2 : 140 [noir, rose], 48 [noir, simple], 45 [blanc, rose], 15 [blanc, simple].

\textbf{Deuxième croisement.} Un troisième gène est étudié : pattes emplumées ($E$ dominant) / pattes nues ($e$). Le croisement-test d'une poule double hétérozygote pour les gènes $N/b$ et $E/e$ donne : 190 [noir, emplumées], 185 [blanc, nues], 32 [noir, nues], 33 [blanc, emplumées].
\begin{enumerate}
\item Interprète les résultats du premier croisement : dominance, proportions F2, conclusion sur la localisation des gènes $N/b$ et $R/r$. Quel type de brassage explique la diversité des phénotypes F2 ?
\item Interprète les résultats du deuxième croisement : les gènes $N/b$ et $E/e$ sont-ils indépendants ou liés ? Justifie.
\item Calcule le pourcentage de recombinaison entre $N/b$ et $E/e$ et donne la distance en cM.
\item Écris le génotype de la poule double hétérozygote (disposition des allèles) et la liste de ses gamètes avec leurs proportions.
\item En mobilisant l'ensemble des résultats, explique comment la méiose et la fécondation assurent le maintien de la diversité génétique au sein de l'espèce, et précise l'intérêt de cette diversité pour l'élevage et l'adaptation des populations.
\end{enumerate}
\end{exercice}

\newpage

\coursec{Corrigés des exercices}

\begin{corrige}[Exercice 1 — QCM : brassages génétiques]
\begin{enumerate}
\item \textbf{Réponse b.} Le brassage interchromosomique résulte de la \cle{séparation aléatoire} des paires de chromosomes homologues en \cle{anaphase I} de méiose : chaque paire s'oriente indépendamment des autres, ce qui mélange les chromosomes d'origine paternelle et maternelle dans les gamètes. Les réponses a et d décrivent le brassage \emph{intra}chromosomique ; la fécondation amplifie la diversité mais n'est pas le brassage interchromosomique.
\item \textbf{Réponse c.} Un croisement-test avec gènes indépendants donne 4 phénotypes équiprobables (25~\% chacun) ; un linkage absolu donne 2 phénotypes (50~\%/50~\%). Ici, deux classes parentales majoritaires (42~\% + 42~\%) et deux classes recombinées minoritaires (8~\% + 8~\%) : les gènes sont \cle{partiellement liés}, avec un taux de recombinaison de $8 + 8 = 16~\%$.
\item \textbf{Réponse a.} Le crossing-over a lieu en \cle{prophase I} de méiose (stade pachytène), entre \cle{chromatides non sœurs} de deux chromosomes homologues appariés en bivalent, au niveau des chiasmas.
\item \textbf{Réponse b.} La \cle{pléiotropie} est l'action d'un \emph{même} gène sur \emph{plusieurs} caractères (ex. : allèle $S$ de l'hémoglobine). La réponse a définit la polygénie.
\end{enumerate}
\end{corrige}

\begin{corrige}[Exercice 2 — Vrai ou faux justifié]
\begin{enumerate}
\item \textbf{Vrai.} Avec deux gènes indépendants à dominance complète, chaque hybride F1 produit 4 gamètes équiprobables ; l'échiquier 4 $\times$ 4 donne 9/16 [A--B--], 3/16 [A--bb], 3/16 [aaB--], 1/16 [aabb].
\item \textbf{Vrai.} En linkage absolu, aucun crossing-over ne sépare les deux gènes : le double hétérozygote $AB//ab$ ne produit que les gamètes parentaux $AB$ et $ab$ en proportions égales (50~\%/50~\%).
\item \textbf{Faux.} C'est l'inverse : plus deux gènes sont \emph{éloignés} sur le chromosome, plus la probabilité qu'un crossing-over survienne entre eux est grande, donc plus le pourcentage de recombinaison est \emph{élevé}. C'est le principe de la cartographie génétique.
\item \textbf{Vrai.} La fécondation réunit au hasard deux gamètes parmi des millions de types possibles : c'est un tirage aléatoire qui multiplie les combinaisons génétiques (chez l'Homme, $2^{23} \times 2^{23}$ combinaisons avant même de compter les crossing-over).
\item \textbf{Vrai.} Avec deux gènes additifs, le nombre de doses d'allèles dominants varie de 0 à 4, ce qui donne 5 classes phénotypiques en proportions 1:4:6:4:1.
\end{enumerate}
\end{corrige}

\begin{corrige}[Exercice 3 — Définitions et vocabulaire]
\begin{enumerate}
\item \textbf{Monogénie} : un caractère est contrôlé par un seul gène (ex. : couleur des fleurs du pois chez Mendel). \textbf{Polygénie} : un caractère est contrôlé par plusieurs gènes à effets additifs (ex. : couleur des grains de blé, couleur de la peau humaine).
\item \textbf{Linkage partiel} : deux gènes sont situés sur le même chromosome mais des crossing-over peuvent les séparer ; le double hétérozygote produit alors des gamètes parentaux majoritaires et des gamètes recombinés minoritaires (ex. : gènes corps/ailes chez la drosophile).
\item \textbf{Pourcentage de recombinaison} : proportion (en \%) d'individus recombinés parmi les descendants d'un croisement-test ; il estime la distance entre deux gènes liés ($1~\% = 1$ cM).
\item \textbf{Carte génétique} : représentation schématique de l'ordre et des distances (en cM) des gènes le long d'un chromosome, établie à partir des taux de recombinaison deux à deux.
\item \textbf{Gamète parental} : gamète portant la même association d'allèles que les chromosomes parentaux reçus par l'hybride (ex. : $AB$ ou $ab$ chez $AB//ab$). \textbf{Gamète recombiné} : gamète portant une association nouvelle issue d'un crossing-over (ex. : $Ab$ ou $aB$).
\end{enumerate}
\end{corrige}

\begin{corrige}[Exercice 4 — Calculs directs]
\begin{enumerate}
\item Les phénotypes recombinés sont les classes minoritaires : $43 + 42 = 85$ individus sur 500.
\[ \text{Taux de recombinaison} = \frac{85}{500} \times 100 = 17~\% \]
Les deux gènes sont distants de \textbf{17 cM}.
\item B étant le gène central, les distances s'additionnent :
\[ d(A-C) = d(A-B) + d(B-C) = 12 + 8 = 20~\text{cM} \]
\begin{popfigure}
\begin{tikzpicture}[scale=1]
\draw[thick, popblue] (0,0) -- (6,0);
\fill[popink] (0,0) circle (2.5pt) node[above=2pt] {\textbf{A}};
\fill[poporange] (3.6,0) circle (2.5pt) node[above=2pt] {\textbf{B}};
\fill[poppurple] (6,0) circle (2.5pt) node[above=2pt] {\textbf{C}};
\draw[<->, popmuted] (0,-0.5) -- (3.6,-0.5) node[midway, below] {12 cM};
\draw[<->, popmuted] (3.6,-0.5) -- (6,-0.5) node[midway, below] {8 cM};
\draw[<->, popmuted] (0,-1.2) -- (6,-1.2) node[midway, below] {20 cM};
\end{tikzpicture}
\legende{Carte génétique des gènes A, B et C. Les distances sont proportionnelles aux pourcentages de recombinaison (1 cM = 1 \%).}
\end{popfigure}
\item Avec un taux de recombinaison de 20~\% sur 1\,000 gamètes : les gamètes \emph{recombinés} représentent 200 gamètes, soit 100 $Ab$ et 100 $aB$ ; les gamètes \emph{parentaux} représentent 800 gamètes, soit 400 $AB$ et 400 $ab$.
\begin{center}
\begin{tabular}{|l|c|c|c|c|}
\hline
\rowcolor{popblueL} Gamète & $AB$ (parental) & $ab$ (parental) & $Ab$ (recombiné) & $aB$ (recombiné) \\
\hline
Effectif & 400 & 400 & 100 & 100 \\
\hline
Proportion & 40~\% & 40~\% & 10~\% & 10~\% \\
\hline
\end{tabular}
\end{center}
\end{enumerate}
\end{corrige}

\begin{corrige}[Exercice 5 — Croisement-test chez la drosophile]
\begin{enumerate}
\item La femelle F1 est issue de lignées pures $GN//GN \times en//en$ : elle a reçu le chromosome $GN$ d'un parent et $en$ de l'autre. Son génotype est donc $GN//en$, allèles dominants en \cle{disposition cis} (sur le même chromosome).
\item Si les gènes étaient indépendants, le croisement-test donnerait 4 phénotypes équiprobables, soit 200 descendants par classe. Or on observe deux classes très majoritaires (332 et 328) et deux classes minoritaires (74 et 66) : les proportions sont très différentes de 1:1:1:1. Les gènes sont donc \cle{liés} (linkage partiel, car il existe des recombinés).
\item Gamètes \textbf{parentaux} (majoritaires) : $GN$ et $en$. Gamètes \textbf{recombinés} (minoritaires, issus de crossing-over) : $Gn$ et $eN$.
\item 
\[ \text{Taux de recombinaison} = \frac{74 + 66}{800} \times 100 = \frac{140}{800} \times 100 = 17{,}5~\% \]
La distance génétique entre les deux gènes est de \textbf{17,5 cM}.
\end{enumerate}
\end{corrige}

\begin{corrige}[Exercice 6 — Dihybridisme chez la souris]
\begin{enumerate}
\item La F1 est homogène [noir, lisse] : les allèles \cle{noir} ($N$) et \cle{lisse} ($L$) sont dominants ; \cle{blanc} ($n$) et \cle{frisé} ($l$) sont récessifs.
\item Rapport théorique 9:3:3:1 sur 320 descendants : 180 ; 60 ; 60 ; 20. Effectifs observés : 180 ; 62 ; 59 ; 19. La concordance est excellente (écarts inférieurs à 5~\%) : les proportions correspondent au 9:3:3:1, donc les deux gènes sont \cle{indépendants} (portés par des chromosomes différents) ; la diversité F2 résulte du \cle{brassage interchromosomique}.
\item Chaque souris F1 $NL//nl$ (gènes indépendants : génotype $N/n ; L/l$) produit 4 gamètes équiprobables : $NL$, $Nl$, $nL$, $nl$.
\begin{center}
\begin{tabular}{|l|c|c|c|c|}
\hline
\rowcolor{popblueL} Gamètes & $NL$ & $Nl$ & $nL$ & $nl$ \\
\hline
$NL$ & $NNLL$ & $NNLl$ & $NnLL$ & $NnLl$ \\
\hline
$Nl$ & $NNLl$ & $NNll$ & $NnLl$ & $Nnll$ \\
\hline
$nL$ & $NnLL$ & $NnLl$ & $nnLL$ & $nnLl$ \\
\hline
$nl$ & $NnLl$ & $Nnll$ & $nnLl$ & $nnll$ \\
\hline
\end{tabular}
\end{center}
Soit 9/16 [noir, lisse], 3/16 [noir, frisé], 3/16 [blanc, lisse], 1/16 [blanc, frisé].
\item Croisement-test $N/n ; L/l \times n/n ; l/l$ : le parent récessif ne fournit que des gamètes $nl$, donc les phénotypes des descendants révèlent directement les gamètes de la F1. Attendu : 4 phénotypes équiprobables (25~\% chacun) : [noir, lisse], [noir, frisé], [blanc, lisse], [blanc, frisé], soit \textbf{50 descendants de chaque type sur 200}.
\end{enumerate}
\end{corrige}

\begin{corrige}[Exercice 7 — Linkage absolu chez le haricot]
\begin{enumerate}
\item La F1 homogène [violettes, lisses] montre que \cle{violet} ($V$) et \cle{lisse} ($L$) dominent \cle{blanc} ($v$) et \cle{ridé} ($l$). Parents : $VL//VL \times vl//vl$ ; F1 : $VL//vl$ (disposition cis).
\item Le croisement-test $VL//vl \times vl//vl$ ne donne que deux phénotypes [violettes, lisses] et [blanches, ridées] en proportions égales ($\approx 50~\%/50~\%$) : la F1 ne produit que deux types de gamètes. Les deux gènes sont portés par le même chromosome et \emph{aucun} crossing-over ne les sépare : c'est le \cle{linkage absolu} (ou liaison complète).
\item Gamètes de la F1 : $VL$ (50~\%) et $vl$ (50~\%), uniquement des gamètes parentaux.
\item C'est le \cle{crossing-over} (enjambement entre chromatides non sœurs de chromosomes homologues) qui est absent ici ; il se produit normalement en \cle{prophase I} de la méiose, au niveau des chiasmas. Son absence explique l'absence de gamètes recombinés $Vl$ et $vL$.
\end{enumerate}
\end{corrige}

\begin{corrige}[Exercice 8 — Polygénie : couleur des grains de blé]
\begin{enumerate}
\item Un monohybridisme classique donne en F2 \emph{deux} phénotypes (3:1) ou trois (1:2:1 en dominance incomplète). Or on observe ici \emph{cinq} classes phénotypiques en proportions 1:4:6:4:1, avec une variation \emph{continue} d'intensité de couleur : un seul gène ne peut pas rendre compte de ces résultats.
\item Modèle à deux gènes additifs : chaque allèle dominant ($A$ ou $B$) ajoute une dose de pigment rouge. Parents : $AABB$ (rouge foncé, 4 doses) $\times$ $aabb$ (blanc, 0 dose). F1 : $AaBb$ (rouge clair, 2 doses). Classes de F2 :
\begin{center}
\begin{tabular}{|l|c|l|}
\hline
\rowcolor{popblueL} Phénotype & Doses & Génotypes \\
\hline
Rouge foncé & 4 & $AABB$ \\
\hline
Rouge & 3 & $AABb$, $AaBB$ \\
\hline
Rouge clair & 2 & $AAbb$, $aaBB$, $AaBb$ \\
\hline
Rose pâle & 1 & $Aabb$, $aaBb$ \\
\hline
Blanc & 0 & $aabb$ \\
\hline
\end{tabular}
\end{center}
\item La F1 $AaBb$ produit 4 gamètes équiprobables ($AB$, $Ab$, $aB$, $ab$). L'échiquier 4 $\times$ 4 (16 cases) donne : $AABB$ : 1/16 ; 3 doses : $2\ AABb + 2\ AaBB = 4/16$ ; 2 doses : $AAbb + aaBB + 4\ AaBb = 6/16$ ; 1 dose : $2\ Aabb + 2\ aaBb = 4/16$ ; $aabb$ : 1/16. Soit bien \textbf{1:4:6:4:1}, conforme aux observations (6 ; 25 ; 38 ; 25 ; 6 sur 100).
\item Il s'agit d'une \cle{polygénie} (hérédité polygénique ou quantitative), exception à la monogénie. Autre exemple chez l'Homme : la \cle{couleur de la peau} (ou la taille).
\end{enumerate}
\end{corrige}

\begin{corrige}[Exercice 9 — Type Bac : trihybridisme avec gènes liés chez la drosophile]
\begin{enumerate}
\item En croisement-test, si les trois gènes étaient indépendants, on attendrait 8 classes équiprobables ($\approx 125$ chacune). Or les effectifs sont très inégaux (de 9 à 362) : les trois gènes sont \cle{liés} sur le même chromosome.
\item \textbf{Classes parentales} : les plus fréquentes — [gris, rouge, normales] (358) et [jaune, écarlate, coupées] (362). La femelle porte donc les chromosomes $g^+\,r^+\,a^+$ et $g\,r\,a$. \textbf{Classes de doubles crossing-over} : les plus rares — [gris, écarlate, normales] (11) et [jaune, rouge, coupées] (9), car un double enjambement est l'événement le moins probable.
\item Comparons un parental au double crossing-over correspondant : $g^+\,r^+\,a^+ \rightarrow g^+\,r\,a^+$ : seul le gène $r$ « change de camp ». Or, dans un double crossing-over, c'est le gène \emph{central} qui est permuté seul. Le gène central est donc $r$ ; l'ordre est \textbf{$g$ — $r$ — $a$}.
\item En réécrivant les classes dans l'ordre $g-r-a$ :
\begin{itemize}
\item Crossing-over simple dans l'intervalle $g$--$r$ : [gris, écarlate, coupées] (92) et [jaune, rouge, normales] (88) $\rightarrow 180$ individus ;
\item Crossing-over simple dans l'intervalle $r$--$a$ : [gris, rouge, coupées] (41) et [jaune, écarlate, normales] (39) $\rightarrow 80$ individus ;
\item Doubles crossing-over : $11 + 9 = 20$ individus (à compter dans les \emph{deux} intervalles).
\end{itemize}
\begin{align*}
d(g-r) &= \frac{180 + 20}{1000} \times 100 = 20~\% = 20~\text{cM} \\
d(r-a) &= \frac{80 + 20}{1000} \times 100 = 10~\% = 10~\text{cM}
\end{align*}
\item Carte génétique :
\begin{popfigure}
\begin{tikzpicture}[scale=1]
\draw[thick, popblue] (0,0) -- (8,0);
\fill[popink] (0,0) circle (2.5pt) node[above=2pt] {\textbf{$g$}};
\fill[poporange] (5.33,0) circle (2.5pt) node[above=2pt] {\textbf{$r$}};
\fill[poppurple] (8,0) circle (2.5pt) node[above=2pt] {\textbf{$a$}};
\draw[<->, popmuted] (0,-0.5) -- (5.33,-0.5) node[midway, below] {20 cM};
\draw[<->, popmuted] (5.33,-0.5) -- (8,-0.5) node[midway, below] {10 cM};
\draw[<->, popmuted] (0,-1.3) -- (8,-1.3) node[midway, below] {30 cM};
\end{tikzpicture}
\legende{Carte génétique des gènes $g$, $r$, $a$ chez la drosophile. Distance $g-r$ = 20 cM, $r-a$ = 10 cM. La distance cartographique $g-a$ est de 30 cM, supérieure au taux de recombinaison observé (26 \%) car les doubles crossing-over ne sont pas détectés comme recombinés pour les gènes extrêmes.}
\end{popfigure}
(Remarque : le taux de recombinaison \emph{observé} entre $g$ et $a$ vaut $\frac{180+80}{1000} = 26~\%$, inférieur à 30 cM, car les doubles crossing-over « annulent » la recombinaison visible entre les extrémités.)
\item Génotype de la femelle triple hétérozygote, dans l'ordre $g-r-a$ :
\[ \frac{g^+\ r^+\ a^+}{g\ r\ a} \quad \text{(disposition cis des trois allèles sauvages)} \]
\end{enumerate}
\textbf{Barème indicatif (sur 10)} : justification du linkage (1,5) ; classes parentales et DCO justifiées (2) ; ordre des gènes déduit de la comparaison (2) ; calculs des deux intervalles avec prise en compte des DCO (2,5) ; carte (1) ; génotype (1).

\textbf{Points de vigilance} : ne pas oublier d'ajouter les doubles crossing-over dans \emph{chaque} intervalle ; justifier le gène central par la comparaison parental/DCO (seul le gène du milieu est permuté) ; distinguer distance cartographique (30 cM) et taux de recombinaison observé (26~\%).
\end{corrige}

\begin{corrige}[Exercice 10 — Type Bac : drépanocytose, monohybridisme et pléiotropie]
\begin{enumerate}
\item Croisement $AS \times AS$ ; chaque parent produit les gamètes $A$ (1/2) et $S$ (1/2) :
\begin{center}
\begin{tabular}{|l|c|c|}
\hline
\rowcolor{popblueL} Gamètes & $A$ (1/2) & $S$ (1/2) \\
\hline
$A$ (1/2) & $AA$ (1/4) sain & $AS$ (1/4) porteur sain \\
\hline
$S$ (1/2) & $AS$ (1/4) porteur sain & $SS$ (1/4) drépanocytaire \\
\hline
\end{tabular}
\end{center}
Probabilités : $AA$ : 1/4 (sain) ; $AS$ : 1/2 (porteur du trait, quasi normal, résistant au paludisme) ; $SS$ : 1/4 (atteint d'anémie grave). Phénotypes : 3/4 non atteints, 1/4 atteint.
\item Arbre généalogique :
\begin{popfigure}
\begin{tikzpicture}[scale=1.1]
\node[draw, circle, minimum size=0.55cm, fill=popgold!60] (mere) at (0,0) {};
\node[draw, rectangle, minimum size=0.55cm, fill=popgold!60] (pere) at (1.2,0) {};
\draw (mere) -- (pere);
\draw (0.6,0) -- (0.6,-0.6) -- (0.6,-0.6);
\draw (0.6,-0.6) -- ++(-1.35,0) -- ++(2.7,0);
\foreach \x/\sym/\fillc in {-0.75/circle/popgold!60, -0.25/rectangle/white, 0.25/circle/white, 0.75/rectangle/popdark}{
  \node[draw, \sym, minimum size=0.45cm, fill=\fillc] at (\x,-1.3) {};
}
\draw (-0.75,-0.6) -- (-0.75,-1.05); \draw (-0.25,-0.6) -- (-0.25,-1.05);
\draw (0.25,-0.6) -- (0.25,-1.05); \draw (0.75,-0.6) -- (0.75,-1.05);
\node[left] at (-1.6,0) {I};
\node[left] at (-1.6,-1.3) {II};
\end{tikzpicture}
\legende{Arbre généalogique d'une famille pour la drépanocytose. Cercle = femme, carré = homme. Symbole doré = hétérozygote $AS$ (porteur sain du trait). Symbole noir = homozygote $SS$ (atteint d'anémie drépanocytaire). Symbole blanc = phénotype sain ($AA$ ou $AS$ non distingué phénotypiquement ici). Génération I : parents ; Génération II : quatre enfants.}
\end{popfigure}
\item L'allèle $S$ agit simultanément sur \emph{plusieurs} caractères : la \textbf{forme} des globules rouges (falciformation), l'\textbf{anémie} (destruction des hématies) et la \textbf{résistance au paludisme}. Un même gène, plusieurs effets phénotypiques : c'est la définition de la \cle{pléiotropie}.
\item En zone d'endémie palustre, les hétérozygotes $AS$ résistent mieux au paludisme grave que les $AA$ : ils survivent et se reproduisent davantage (\cle{avantage sélectif de l'hétérozygote}). L'allèle $S$ est ainsi transmis préférentiellement malgré l'élimination des $SS$ : c'est un \cle{polymorphisme équilibré}, ce qui explique la fréquence élevée de $AS$ (\num{0,32}) dans le document 2.
\item Message possible : « Avant le mariage, faites le test de dépistage de la drépanocytose dans votre centre de santé. Si vous êtes tous les deux porteurs du trait ($AS$), sachez qu'à chaque grossesse votre enfant a 1 chance sur 4 d'être atteint ($SS$), 1 chance sur 2 d'être porteur sain et 1 chance sur 4 d'être indemne. Un conseil génétique vous aidera à faire des choix éclairés. »
\end{enumerate}
\textbf{Barème indicatif (sur 10)} : échiquier et probabilités (2) ; pedigree conforme aux conventions (2) ; pléiotropie démontrée à partir des trois effets (2) ; explication de l'avantage sélectif des hétérozygotes (2,5) ; message de sensibilisation (1,5).

\textbf{Points de vigilance} : citer les probabilités \emph{à chaque naissance} (les tirages sont indépendants) ; ne pas écrire que $AS$ est « malade » : c'est un porteur quasi sain ; pour le document 2, la justification attendue est explicitement l'\emph{avantage sélectif de l'hétérozygote} en zone palustre.
\end{corrige}

\begin{corrige}[Exercice 11 — Type Bac : synthèse, brassage génétique et diversité chez la volaille]
\begin{enumerate}
\item \textbf{Premier croisement.} La F1 homogène [noir, rose] montre que $N$ (noir) domine $b$ (blanc) et $R$ (rose) domine $r$ (simple). Total F2 : $140+48+45+15 = 248$. Rapport théorique 9:3:3:1 : $139{,}5$ ; $46{,}5$ ; $46{,}5$ ; $15{,}5$ — concordance excellente avec les effectifs observés. Les gènes $N/b$ et $R/r$ sont donc \cle{indépendants} ; la diversité des phénotypes F2 s'explique par le \cle{brassage interchromosomique} (ségrégation indépendante des paires d'homologues en anaphase I), puis par la rencontre aléatoire des gamètes à la fécondation.
\item \textbf{Deuxième croisement.} Croisement-test : si les gènes étaient indépendants, on attendrait 4 classes équiprobables ($\approx 110$ chacune sur 440). Or on observe deux classes parentales majoritaires (190 et 185) et deux classes recombinées minoritaires (32 et 33) : les gènes $N/b$ et $E/e$ sont \cle{liés} (linkage partiel).
\item 
\[ \text{Taux de recombinaison} = \frac{32 + 33}{440} \times 100 = \frac{65}{440} \times 100 \approx 14{,}8~\% \]
Distance génétique : \textbf{$\approx 14{,}8$ cM} (environ 15 cM).
\item Les gamètes parentaux étant $NE$ et $be$, la poule est en disposition \cle{cis} : $NE//be$. Ses gamètes :
\begin{center}
\begin{tabular}{|l|c|c|c|c|}
\hline
\rowcolor{popblueL} Gamète & $NE$ (parental) & $be$ (parental) & $Ne$ (recombiné) & $bE$ (recombiné) \\
\hline
Proportion & 42,6~\% & 42,6~\% & 7,4~\% & 7,4~\% \\
\hline
\end{tabular}
\end{center}
\item \textbf{Synthèse.} La méiose crée la diversité gamétique par deux brassages : le brassage \emph{interchromosomique} (répartition aléatoire des chromosomes homologues — illustré par les gènes $N/b$ et $R/r$) et le brassage \emph{intrachromosomique} (crossing-over en prophase I — illustré par les gamètes recombinés $Ne$ et $bE$). La fécondation réunit ensuite au hasard deux gamètes différents, multipliant les combinaisons : chaque individu (sauf jumeaux vrais) est génétiquement unique. Pour l'éleveur de Bafoussam, cette diversité permet de \emph{sélectionner} des combinaisons favorables (ponde, rusticité) ; pour l'espèce, elle constitue le réservoir de variabilité qui permet l'\emph{adaptation} des populations aux changements du milieu (maladies, climat).
\end{enumerate}
\textbf{Barème indicatif (sur 10)} : premier croisement interprété avec vérification du 9:3:3:1 (2,5) ; linkage démontré par comparaison au 1:1:1:1 (2) ; calcul du taux de recombinaison (1,5) ; génotype cis et gamètes avec proportions (2) ; synthèse méiose + fécondation + intérêt (2).

\textbf{Points de vigilance} : toujours \emph{comparer quantitativement} les proportions observées aux proportions théoriques avant de conclure (indépendance ou liaison) ; dans un croisement-test, les phénotypes des descendants révèlent directement les gamètes de l'hybride — le dire explicitement ; ne pas confondre les deux types de brassage dans la synthèse finale.
\end{corrige}

\newpage

\coursec{Fiche bilan}

\begin{popfigure}
\begin{tikzpicture}[mindmap, grow cyclic, every node/.style={concept, text=white, font=\small\bfseries},
  root concept/.append style={concept color=popdark, font=\normalsize\bfseries},
  level 1/.append style={level distance=3.6cm, sibling angle=60},
  level 2/.append style={level distance=2.2cm, sibling angle=45, font=\scriptsize}]
\node[root concept]{Méiose, fécondation et diversité génétique}
  child[concept color=popblue]{ node{Monohybridisme}
    child{ node{$1^{\text{re}}$ loi : uniformité de $F_1$} }
    child{ node{$2^{\text{nde}}$ loi : ségrégation $3/4 : 1/4$} }
  }
  child[concept color=poporange]{ node{Dihybridisme et trihybridisme}
    child{ node{$3^{\text{ème}}$ loi : ségrégation indépendante} }
    child{ node{$F_2$ : $9/16 : 3/16 : 3/16 : 1/16$} }
  }
  child[concept color=popgreen]{ node{Brassage interchromosomique}
    child{ node{Gènes sur chromosomes différents} }
    child{ node{$2^n$ types de gamètes} }
  }
  child[concept color=poppurple]{ node{Brassage intrachromosomique}
    child{ node{Linkage total : 2 gamètes} }
    child{ node{Crossing-over : 4 gamètes} }
  }
  child[concept color=poppink]{ node{Carte génétique}
    child{ node{Taux de recombinaison $p$} }
    child{ node{Distances en cM} }
  }
  child[concept color=popgold]{ node{Exceptions à la monogénie}
    child{ node{Polygénie} }
    child{ node{Pléiotropie} }
  };
\end{tikzpicture}
\legende{Carte mentale de la leçon : les deux brassages génétiques (interchromosomique et intrachromosomique), les lois de Mendel, la cartographie factorielle et les exceptions à la monogénie.}
\end{popfigure}

\begin{aretenir}[L'essentiel de la leçon : Définitions clés]
\begin{itemize}
  \item \cle{Monohybridisme} : croisement de deux lignées pures différant par un seul caractère (une paire d'allèles).
  \item \cle{Dihybridisme} / \cle{trihybridisme} : croisements portant sur deux ou trois caractères.
  \item \cle{Brassage interchromosomique} : redistribution des allèles de gènes situés sur des chromosomes différents, grâce à la séparation au hasard des chromosomes homologues en anaphase I de méiose.
  \item \cle{Brassage intrachromosomique} : échange de fragments de chromatides entre chromosomes homologues (\cle{crossing-over}) en prophase I, pour des gènes situés sur le même chromosome.
  \item \cle{Linkage} : liaison de gènes portés par un même chromosome ; \emph{total} si aucun crossing-over, \emph{partiel} si des crossing-overs produisent des gamètes recombinés.
  \item \cle{Polygénie} : un caractère dépend de plusieurs gènes. \cle{Pléiotropie} : un seul gène agit sur plusieurs caractères.
\end{itemize}
\end{aretenir}

\begin{propriete}[Résultats et formules à connaître par cœur]
\begin{center}
\begin{tabularx}{\linewidth}{|l|X|}
\hline
\rowcolor{popblueL} \textbf{Situation} & \textbf{Résultat attendu} \\
\hline
$F_1$ de monohybridisme (lignées pures) & 100 \% homogènes ($1^{\text{re}}$ loi de Mendel) \\
\hline
$F_2$ de monohybridisme & $3/4$ dominant : $1/4$ récessif \\
\hline
Test-cross monohybride & $1/2 : 1/2$ \\
\hline
$F_2$ de dihybridisme (gènes indépendants) & $9/16 : 3/16 : 3/16 : 1/16$ \\
\hline
Test-cross dihybride, gènes indépendants & 4 phénotypes équiprobables (25 \% chacun) \\
\hline
Test-cross, linkage total & 2 phénotypes parentaux seulement (50 \% -- 50 \%) \\
\hline
Test-cross, linkage partiel & 2 parentaux majoritaires $+$ 2 recombinés minoritaires \\
\hline
Nombre de types de gamètes & $2^n$ ($n$ = nombre de paires d'allèles en hétérozygotie, gènes indépendants) \\
\hline
Taux de recombinaison & $p = \dfrac{\text{nombre d'individus recombinés}}{\text{effectif total}} \times 100$ (en \%, assimilé à des cM) \\
\hline
\end{tabularx}
\end{center}
\end{propriete}

\begin{methode}[Méthodes express]
\begin{itemize}
  \item \emph{Gènes liés ou indépendants ?} Au test-cross : 4 classes équiprobables $\Rightarrow$ indépendance ; 2 classes $\Rightarrow$ linkage total ; 2 classes majoritaires $+$ 2 minoritaires $\Rightarrow$ linkage partiel.
  \item \emph{Carte génétique} : identifier les classes parentales (les plus fréquentes), repérer les doubles crossing-overs (les plus rares) pour trouver le gène central, calculer chaque taux de recombinaison, puis placer les gènes en cM.
  \item \emph{Unicité génétique} : la diversité vient du brassage interchromosomique ($2^{23}$ combinaisons chez l'Homme), du crossing-over et de la rencontre au hasard des gamètes à la fécondation.
\end{itemize}
\end{methode}

\begin{aretenir}[Checklist avant le Bac]
\begin{itemize}
  \item $\square$ Je sais définir hybride, lignée pure, allèle, génotype, phénotype, test-cross.
  \item $\square$ Je connais les trois lois de Mendel et les proportions types ($3/4$--$1/4$ ; $9/16$--$3/16$--$3/16$--$1/16$).
  \item $\square$ Je sais construire un échiquier de croisement et justifier les proportions de $F_2$.
  \item $\square$ Je sais distinguer brassage interchromosomique et intrachromosomique et les relier aux étapes de la méiose.
  \item $\square$ Je sais reconnaître un linkage total, un linkage partiel ou une ségrégation indépendante à partir des résultats d'un test-cross.
  \item $\square$ Je sais calculer un taux de recombinaison et l'exprimer en \% (ou cM).
  \item $\square$ Je sais repérer les classes parentales (majoritaires) et les doubles crossing-overs (minoritaires) dans un trihybridisme.
  \item $\square$ Je sais établir une carte génétique : ordre des gènes et distances en cM.
  \item $\square$ Je sais définir et illustrer la polygénie (couleur du grain de blé, taille) et la pléiotropie (drépanocytose).
  \item $\square$ Je sais expliquer en quoi méiose $+$ fécondation assurent l'unicité génétique de chaque individu.
  \item $\square$ Je rédige les génotypes correctement : $A/B$ pour des gènes indépendants, $AB/ab$ pour des gènes liés.
  \item $\square$ Je vérifie toujours que la somme des pourcentages d'un croisement vaut 100 \%.
\end{itemize}
\end{aretenir}

\findecours

\end{document}
